 \documentclass[11pt,letterpaper]{report}

\usepackage[utf8]{inputenc}
\usepackage[T1]{fontenc}
\usepackage{lmodern}
\usepackage[spanish,es-tabla]{babel}
\usepackage[margin=2cm]{geometry}
\usepackage{graphicx}
\usepackage{float}
\usepackage{array}
\usepackage{booktabs}
\usepackage{xltabular}
\usepackage{enumitem}
\usepackage{caption}
\usepackage{titlesec}
\usepackage[hidelinks]{hyperref}

\titleformat{\chapter}
{\normalfont\LARGE\bfseries}
{\thechapter.}{0.6em}{}
\titlespacing*{\chapter}{0pt}{0pt}{1.5em}

\newcommand{\nem}[1]{\texttt{#1}}
\newcommand{\us}{\_\allowbreak{}}
\newcommand{\bs}{\textbackslash\allowbreak{}}

% Captura o diagrama: \captura[ancho]{archivo}{título}{etiqueta}
% El ancho es opcional (por defecto, el de la página).
% Todas las imágenes se guardan en la carpeta recursos/
\newcommand{\captura}[4][\textwidth]{%
	\begin{figure}[H]
		\centering
		\IfFileExists{recursos/#2}%
		{\includegraphics[width=#1]{recursos/#2}}%
		{\fbox{\parbox[c][4cm][c]{0.9\textwidth}{\centering\small
					Imagen pendiente: \texttt{recursos/\detokenize{#2}}}}}
		\caption{#3}
		\label{fig:#4}
	\end{figure}
}

\setlength{\emergencystretch}{2em}

\newcolumntype{L}[1]{>{\raggedright\arraybackslash}p{#1}}
\newcolumntype{Y}{>{\raggedright\arraybackslash}X}

\newcommand{\encabezadodiccionario}{%
	\toprule
	\textbf{Nemónico} & \textbf{Descripción} & \textbf{Dominio} &
	\textbf{Formato} & \textbf{Rango} & \textbf{Nulo} &
	\textbf{Llave} & \textbf{Referencia} \\
	\midrule
}

\newenvironment{diccionario}[1]{%
	\footnotesize
	\setlength{\tabcolsep}{3pt}
	\xltabular{\textwidth}{L{2.6cm} Y L{2.2cm} L{2.9cm} L{1.15cm} L{0.9cm} L{1.6cm} L{2.2cm}}
	\caption{Diccionario de datos de la tabla \nem{#1}}\\
	\encabezadodiccionario
	\endfirsthead
	\multicolumn{8}{l}{\textit{Continuación de la tabla \nem{#1}}}\\
	\encabezadodiccionario
	\endhead
	\bottomrule
	\endfoot
}{%
	\endxltabular
}

\begin{document}
	\begin{titlepage}
	\centering
	{\bfseries\LARGE Tecnol\'ogico Nacional de M\'exico \par}
	{\bfseries\LARGE Instituto Tecnol\'ogico de Veracruz \par}
	\vspace{1cm}
	{\scshape\LARGE Ingenier\'ia en Sistemas Computacionales \par}
	
	\vfill
	
	{\scshape\huge Sistema de Inversiones\par}
	{\itshape\LARGE Manual T\'ecnico - Unidad 4\par}
	
	\vfill
	
	{\Large Alumnos:\par}
	{\Large Daniel Landero Arias\par}
	{\Large Diego de Jes\'us Mej\'ia Rodr\'iguez\par}
	{\Large Eduardo Novoa Siles\par}
	{\Large Ximena Zaleta Hern\'andez\par}
	{\Large Pedro Armando Zuñiga Duran\par}
	\vspace{0.5cm}
	{\Large Maestro:\par}
	{\Large Ra\'ul H\'ector Trujillo Jim\'enez\par}
	{\Large Materia:\par}
	{\Large Administraci\'on de Bases de Datos\par}
	
	\vfill
	
	{\Large \today \par}
\end{titlepage}
	
	\tableofcontents
	\listoftables
	\listoffigures
	
	\chapter{Introducci\'on}
\label{cap:introduccion}

\section{Descripci\'on del sistema}

El sistema administra las inversiones documentadas mediante pagar\'es de inversi\'on con rendimientos mensuales.

La empresa emite los pagar\'es a trav\'es de sus sucursales. Los clientes participan como inversionistas (acreedores) y, cuando existe, un aval garantiza el pago de las obligaciones de la empresa. El personal de la empresa es el \'unico que opera el sistema; los clientes y los avales no acceden directamente a la base de datos, sino que reciben de la empresa la informaci\'on que les corresponde. El pagar\'e establece los siguientes elementos:

\begin{itemize}[leftmargin=*]
	\item \textbf{Datos del pagar\'e:} n\'umero de pagar\'e, monto de la inversi\'on y su moneda, y el lugar y la fecha de suscripci\'on.
	\item \textbf{Partes:} el emisor y el inversionista, cada uno con su nombre completo y domicilio.
	\item \textbf{Fecha de vencimiento:} fecha en la que el emisor debe pagar la cantidad principal de la inversi\'on.
	\item \textbf{Intereses ordinarios y pagos mensuales:} el emisor paga un inter\'es ordinario mensual a una tasa fija (o una tasa anual divisible entre doce), calculado sobre el saldo del capital insoluto. Los pagos se realizan de manera mensual y consecutiva en un d\'ia espec\'ifico de cada mes, y concluyen en la fecha de vencimiento junto con la devoluci\'on del capital principal.
	\item \textbf{Intereses moratorios:} en caso de mora o retraso en el pago de cualquiera de las mensualidades de intereses o del capital principal, el emisor paga un inter\'es moratorio a una tasa mensual sobre los saldos vencidos y no pagados, desde la fecha de su vencimiento hasta su total liquidaci\'on.
	\item \textbf{Disposiciones generales:} el pagar\'e se rige por la Ley General de T\'itulos y Operaciones de Cr\'edito, y las partes se someten a la jurisdicci\'on de los tribunales de una ciudad determinada.
	\item \textbf{Aval (opcional):} persona que garantiza el pago de las obligaciones del emisor, con su nombre completo, domicilio e identificaci\'on.
\end{itemize}

\section{Objetivo general}

Generar un esquema que permita tener un control adecuado de un sistema de inversiones, as\'i como la optimizaci\'on de las liquidaciones de las inversiones proporcionadas por los clientes. La base de datos deber\'a almacenar de manera adecuada los datos de los clientes, asegurando su integridad.

\section{Objetivos espec\'ificos}

\begin{enumerate}[leftmargin=*]
	\item Consultar los datos de las personas que tengan inversiones otorgadas en un periodo, por nombre.
	\item Consultar los datos de las inversiones que est\'en vencidas en un periodo.
	\item Consultar los estados de cuenta de los inversionistas.
	\item Consultar la inversi\'on y su tabla de inter\'es en un periodo.
	\item Consultar las inversiones hechas en un periodo.
\end{enumerate}

\section{Consultas de los objetivos espec\'ificos}
\label{sec:consultas-objetivos}

Cada objetivo espec\'ifico se resuelve con una consulta que puede ejecutar cualquier cuenta de operaci\'on (gerente, ejecutivo, tesorer\'ia o auditor), ya que solo requiere el privilegio \nem{SELECT}. Cada objetivo se acompa\~na de una captura que muestra el formato de la respuesta, igual en ambos gestores, obtenida con la cuenta \nem{op\us gerente}. Las consultas solo leen datos, por lo que no requieren identificarse con \nem{SET} (secci\'on~\ref{sec:identificacion}).

\paragraph{Par\'ametros.} Los valores de entrada (el periodo, el nombre buscado, el cliente o el pagar\'e) se escriben directamente en el \nem{WHERE} de cada consulta. Para hacer otra consulta solo se cambian esos valores.

\paragraph{Versiones por instancia.} Los objetivos 1 y 4 se resuelven con la misma consulta en las dos instancias. Los objetivos 2, 3 y 5 necesitan datos calculados, como el estado de una inversi\'on, el inter\'es pendiente o la mora estimada, y tienen una versi\'on por instancia:

\begin{itemize}[leftmargin=*]
	\item \textbf{Primera instancia:} la consulta se resuelve directamente sobre las tablas, por lo que incluye los c\'alculos. Es m\'as larga: obtiene primero el saldo o lo pagado de cada periodo de inter\'es en una expresi\'on com\'un (\nem{periodo}) y despu\'es lo agrega por inversi\'on.
	\item \textbf{Segunda instancia:} la consulta se resuelve sobre las vistas derivadas de la segmentaci\'on, \nem{interes\us saldo} e \nem{inversion\us estado} (secci\'on~\ref{sec:segmentacion-derivada}), que ya contienen esos c\'alculos. La consulta se reduce a elegir columnas y filtrar.
\end{itemize}

Con los datos de prueba, ambas versiones devuelven exactamente el mismo resultado.

\paragraph{Fecha de corte.} El estado de una inversi\'on y los montos vencidos dependen del d\'ia en que se consultan. Las consultas calculan siempre a la fecha del sistema (\nem{CURRENT\us DATE}), por lo que el mismo dato produce resultados distintos conforme pasa el tiempo. Las capturas y los ejemplos de este manual se obtuvieron con los datos de prueba, cuyo \'ultimo movimiento es del 28 de octubre de 2026; los valores que dependen de la fecha, como los d\'ias vencida y la mora estimada, cambian cada d\'ia.

\paragraph{Criterios comunes.} Las consultas aplican las reglas de negocio del modelo conceptual (cap\'itulo~\ref{cap:conceptual}):

\begin{itemize}[leftmargin=*]
	\item El estado de una inversi\'on es \textit{liquidada} si tiene una liquidaci\'on registrada, \textit{vencida} si su fecha de vencimiento ya pas\'o y \textit{vigente} en los dem\'as casos.
	\item El inter\'es de un periodo es exigible cuando su fecha final es anterior o igual a hoy. Se consideran todos los pagos registrados.
	\item El inter\'es moratorio se calcula como en los datos de prueba: sobre la tasa moratoria mensual, en proporci\'on a los d\'ias de atraso y con meses de 30 d\'ias.
\end{itemize}

\paragraph{Diferencias entre gestores.} Solo cambian de un gestor a otro la b\'usqueda por nombre del objetivo~1, la presentaci\'on en formato vertical de los datos del pagar\'e del objetivo~4 y la resta de fechas, que en la primera instancia solo aparece en el objetivo~2 (en PostgreSQL restar dos fechas devuelve los d\'ias; en MySQL se utiliza \nem{DATEDIFF}). Las dem\'as consultas se ejecutan sin cambios en ambos gestores.

\subsection{Objetivo 1}

\textit{Consultar los datos de las personas que tengan inversiones otorgadas en un periodo, por nombre.}

La consulta lista a los inversionistas que suscribieron al menos un pagar\'e en el periodo y cuyo nombre completo contiene el texto buscado. Devuelve un rengl\'on por persona, con sus datos, el n\'umero de inversiones del periodo y el capital total invertido. Es la misma en las dos instancias.

\paragraph{B\'usqueda sin acentos.} La intercalaci\'on de la base de datos en MySQL compara los textos sin distinguir may\'usculas, min\'usculas ni acentos (cap\'itulo~\ref{cap:implementacion}). PostgreSQL s\'i los distingue, por lo que su consulta usa \nem{ILIKE} para ignorar las may\'usculas y la funci\'on \nem{translate} para sustituir las vocales acentuadas y la letra \~n del nombre por su letra base.

\paragraph{PostgreSQL.}

\begin{verbatim}
	SELECT c.cliente_num, c.cliente_rfc,
	       c.cliente_nombre, c.cliente_apellido,
	       c.cliente_direccion, c.cliente_telefono,
	       COUNT(*)                 AS inversiones,
	       SUM(i.inversion_capital) AS capital_total
	FROM inversion i
	  JOIN cliente c ON c.cliente_num = i.cliente_num
	WHERE i.pagare_fecha BETWEEN '2025-01-01' AND '2025-12-31'
	  AND translate(c.cliente_nombre || ' ' || c.cliente_apellido,
	                'ÁÉÍÓÚÜÑáéíóúüñ', 'AEIOUUNaeiouun') ILIKE '%jose%'
	GROUP BY c.cliente_num
	ORDER BY c.cliente_apellido, c.cliente_nombre;
\end{verbatim}

\paragraph{MySQL.}

\begin{verbatim}
	SELECT c.cliente_num, c.cliente_rfc,
	       c.cliente_nombre, c.cliente_apellido,
	       c.cliente_direccion, c.cliente_telefono,
	       COUNT(*)                 AS inversiones,
	       SUM(i.inversion_capital) AS capital_total
	FROM inversion i
	  JOIN cliente c ON c.cliente_num = i.cliente_num
	WHERE i.pagare_fecha BETWEEN '2025-01-01' AND '2025-12-31'
	  AND CONCAT(c.cliente_nombre, ' ', c.cliente_apellido) LIKE '%jose%'
	GROUP BY c.cliente_num
	ORDER BY c.cliente_apellido, c.cliente_nombre;
\end{verbatim}

Con los par\'ametros del ejemplo, el a\~no 2025 y la b\'usqueda \nem{jose}, ambas consultas devuelven un rengl\'on: el cliente 1, con una inversi\'on de \$75{,}000.00.

\captura{objetivo1.png}{Formato de la respuesta del objetivo 1}{objetivo1}

\subsection{Objetivo 2}

\textit{Consultar los datos de las inversiones que est\'en vencidas en un periodo.}

La consulta lista las inversiones vencidas, es decir, las que ya pasaron su fecha de vencimiento sin tener una liquidaci\'on registrada, cuyo vencimiento cae dentro del periodo. Para cada una muestra los d\'ias vencida, los periodos sin pagar, el inter\'es vencido y la mora estimada:

\begin{itemize}[leftmargin=*]
	\item \textbf{D\'ias vencida:} d\'ias transcurridos desde el vencimiento hasta hoy.
	\item \textbf{Periodos sin pagar e inter\'es vencido:} periodos terminados que conservan un saldo, y la suma de esos saldos. Un periodo con pago parcial cuenta por su saldo restante.
	\item \textbf{Mora estimada:} inter\'es moratorio que la empresa adeudar\'ia si liquidara hoy. Corresponde al capital por los d\'ias desde el vencimiento, m\'as el saldo de cada periodo por los d\'ias desde su fecha final. Es un c\'alculo y no un valor almacenado: se registra hasta que la empresa lo paga, en \nem{pago} o en \nem{liquidacion}.
\end{itemize}

\paragraph{Primera instancia, PostgreSQL.}

\begin{verbatim}
	WITH periodo AS (
	  -- Saldo de interés de cada periodo ya terminado
	  SELECT t.inversion_num, t.interes_fecha_final,
	         t.interes_monto - COALESCE(SUM(g.pago_interes), 0) AS saldo
	  FROM interes t
	    LEFT JOIN pago g
	      ON g.inversion_num = t.inversion_num
	     AND g.interes_num = t.interes_num
	  WHERE t.interes_fecha_final < CURRENT_DATE
	  GROUP BY t.inversion_num, t.interes_num,
	           t.interes_fecha_final, t.interes_monto
	)
	SELECT i.inversion_num, i.pagare_num,
	       CONCAT(c.cliente_nombre, ' ', c.cliente_apellido) AS cliente,
	       s.sucursal_nombre,
	       i.inversion_capital AS capital,
	       i.inversion_fecha_vencimiento AS vencimiento,
	       CURRENT_DATE - i.inversion_fecha_vencimiento AS dias_vencida,
	       SUM(CASE WHEN p.saldo > 0 THEN 1 ELSE 0 END) AS periodos_sin_pagar,
	       SUM(CASE WHEN p.saldo > 0 THEN p.saldo ELSE 0 END) AS interes_vencido,
	       ROUND(i.inversion_capital * i.inversion_tasa_moratoria / 100
	             * (CURRENT_DATE - i.inversion_fecha_vencimiento) / 30
	             + SUM(CASE WHEN p.saldo > 0
	                        THEN p.saldo * i.inversion_tasa_moratoria / 100
	                             * (CURRENT_DATE - p.interes_fecha_final) / 30
	                        ELSE 0 END), 2) AS mora_estimada
	FROM inversion i
	  JOIN cliente c  ON c.cliente_num = i.cliente_num
	  JOIN sucursal s ON s.sucursal_num = i.sucursal_num
	  JOIN periodo p  ON p.inversion_num = i.inversion_num
	WHERE i.inversion_fecha_vencimiento BETWEEN '2025-01-01' AND '2025-12-31'
	  AND i.inversion_fecha_vencimiento < CURRENT_DATE
	  AND NOT EXISTS (SELECT 1 FROM liquidacion l
	                  WHERE l.inversion_num = i.inversion_num)
	GROUP BY i.inversion_num, c.cliente_num, s.sucursal_num
	ORDER BY i.inversion_fecha_vencimiento;
\end{verbatim}

\paragraph{Primera instancia, MySQL.} Es igual a la de PostgreSQL, salvo que cada resta de fechas usa \nem{DATEDIFF}.

\begin{verbatim}
	WITH periodo AS (
	  -- Saldo de interés de cada periodo ya terminado
	  SELECT t.inversion_num, t.interes_fecha_final,
	         t.interes_monto - COALESCE(SUM(g.pago_interes), 0) AS saldo
	  FROM interes t
	    LEFT JOIN pago g
	      ON g.inversion_num = t.inversion_num
	     AND g.interes_num = t.interes_num
	  WHERE t.interes_fecha_final < CURRENT_DATE
	  GROUP BY t.inversion_num, t.interes_num,
	           t.interes_fecha_final, t.interes_monto
	)
	SELECT i.inversion_num, i.pagare_num,
	       CONCAT(c.cliente_nombre, ' ', c.cliente_apellido) AS cliente,
	       s.sucursal_nombre,
	       i.inversion_capital AS capital,
	       i.inversion_fecha_vencimiento AS vencimiento,
	       DATEDIFF(CURRENT_DATE, i.inversion_fecha_vencimiento) AS dias_vencida,
	       SUM(CASE WHEN p.saldo > 0 THEN 1 ELSE 0 END) AS periodos_sin_pagar,
	       SUM(CASE WHEN p.saldo > 0 THEN p.saldo ELSE 0 END) AS interes_vencido,
	       ROUND(i.inversion_capital * i.inversion_tasa_moratoria / 100
	             * DATEDIFF(CURRENT_DATE, i.inversion_fecha_vencimiento) / 30
	             + SUM(CASE WHEN p.saldo > 0
	                        THEN p.saldo * i.inversion_tasa_moratoria / 100
	                             * DATEDIFF(CURRENT_DATE, p.interes_fecha_final) / 30
	                        ELSE 0 END), 2) AS mora_estimada
	FROM inversion i
	  JOIN cliente c  ON c.cliente_num = i.cliente_num
	  JOIN sucursal s ON s.sucursal_num = i.sucursal_num
	  JOIN periodo p  ON p.inversion_num = i.inversion_num
	WHERE i.inversion_fecha_vencimiento BETWEEN '2025-01-01' AND '2025-12-31'
	  AND i.inversion_fecha_vencimiento < CURRENT_DATE
	  AND NOT EXISTS (SELECT 1 FROM liquidacion l
	                  WHERE l.inversion_num = i.inversion_num)
	GROUP BY i.inversion_num, c.cliente_num, s.sucursal_num
	ORDER BY i.inversion_fecha_vencimiento;
\end{verbatim}

\paragraph{Segunda instancia.} Igual en ambos gestores. Las columnas provienen de la vista \nem{inversion\us estado} (tabla~\ref{tab:vista-estado}).

\begin{verbatim}
	SELECT inversion_num, pagare_num, cliente, sucursal_nombre, capital,
	       vencimiento, dias_vencida, periodos_sin_pagar,
	       interes_vencido, mora_estimada
	FROM inversion_estado
	WHERE estado = 'Vencida'
	  AND vencimiento BETWEEN '2025-01-01' AND '2025-12-31'
	ORDER BY vencimiento;
\end{verbatim}

Con el periodo del ejemplo, el a\~no 2025, ambas versiones devuelven las inversiones 16, 18, 20 y 17, en ese orden de vencimiento, cada una con dos periodos sin pagar. La mora estimada de la inversi\'on 16 es el capital de \$500{,}000 por 2.5\,\% mensual y los d\'ias transcurridos desde el 5 de julio de 2025 entre 30, m\'as el saldo de \$4{,}250 de cada periodo pendiente por la misma tasa y sus d\'ias de atraso; crece cada d\'ia.

\captura{objetivo2.png}{Formato de la respuesta del objetivo 2}{objetivo2}

\subsection{Objetivo 3}

\textit{Consultar los estados de cuenta de los inversionistas.}

El estado de cuenta de un inversionista, identificado por su RFC, muestra un rengl\'on por cada una de sus inversiones, con:

\begin{itemize}[leftmargin=*]
	\item El estado de la inversi\'on.
	\item El inter\'es exigible, el pagado y el pendiente.
	\item El inter\'es moratorio pagado, tanto en los pagos mensuales como en la liquidaci\'on.
	\item El capital devuelto.
\end{itemize}

Si se unieran directamente \nem{interes} y \nem{pago}, el monto de cada periodo se repetir\'ia por cada pago y las sumas resultar\'ian multiplicadas. Por ello, la versi\'on de la primera instancia obtiene antes lo pagado de cada periodo (\nem{periodo}). Ambas versiones son iguales en los dos gestores.

\paragraph{Primera instancia.}

\begin{verbatim}
	WITH periodo AS (
	  -- Lo pagado de cada periodo de interés
	  SELECT t.inversion_num, t.interes_fecha_final, t.interes_monto,
	         COALESCE(SUM(g.pago_interes), 0) AS pagado,
	         COALESCE(SUM(g.pago_mora), 0)    AS mora_pagada
	  FROM interes t
	    LEFT JOIN pago g
	      ON g.inversion_num = t.inversion_num
	     AND g.interes_num = t.interes_num
	  GROUP BY t.inversion_num, t.interes_num,
	           t.interes_fecha_final, t.interes_monto
	)
	SELECT i.inversion_num, i.pagare_num, i.pagare_fecha AS suscripcion,
	       i.inversion_fecha_vencimiento AS vencimiento,
	       i.inversion_capital AS capital,
	       CASE WHEN l.inversion_num IS NOT NULL THEN 'Liquidada'
	            WHEN i.inversion_fecha_vencimiento < CURRENT_DATE THEN 'Vencida'
	            ELSE 'Vigente'
	       END AS estado,
	       SUM(CASE WHEN p.interes_fecha_final <= CURRENT_DATE
	                THEN p.interes_monto ELSE 0 END) AS interes_exigible,
	       SUM(p.pagado) AS interes_pagado,
	       SUM(CASE WHEN p.interes_fecha_final <= CURRENT_DATE
	                THEN p.interes_monto ELSE 0 END) - SUM(p.pagado)
	         AS interes_pendiente,
	       SUM(p.mora_pagada) + COALESCE(l.liquidacion_mora, 0) AS mora_pagada,
	       COALESCE(l.liquidacion_capital, 0.00) AS capital_devuelto
	FROM cliente c
	  JOIN inversion i ON i.cliente_num = c.cliente_num
	  JOIN periodo p   ON p.inversion_num = i.inversion_num
	  LEFT JOIN liquidacion l ON l.inversion_num = i.inversion_num
	WHERE c.cliente_rfc = 'MEAP850304BV0'
	GROUP BY i.inversion_num, l.inversion_num
	ORDER BY i.pagare_fecha;
\end{verbatim}

\paragraph{Segunda instancia.} Todos los datos provienen de la vista \nem{inversion\us estado}.

\begin{verbatim}
	SELECT inversion_num, pagare_num, suscripcion, vencimiento, capital,
	       estado, interes_exigible, interes_pagado, interes_pendiente,
	       mora_pagada, capital_devuelto
	FROM inversion_estado
	WHERE cliente_rfc = 'MEAP850304BV0'
	ORDER BY suscripcion;
\end{verbatim}

El ejemplo corresponde a la clienta Patricia Mendoza Aguilar, que tiene dos inversiones. La inversi\'on 10 est\'a liquidada: se devolvi\'o con 20 d\'ias de atraso, con un inter\'es moratorio de $75{,}000 \times 2.5\mkern 3mu\% \times 20/30 = 1{,}250.00$. La inversi\'on 35 est\'a vigente y al corriente; su mora pagada de \$7.50 proviene del pago parcial del periodo de junio de 2026.

\captura{objetivo3.png}{Formato de la respuesta del objetivo 3}{objetivo3}

\subsection{Objetivo 4}

\textit{Consultar la inversi\'on y su tabla de inter\'es en un periodo.}

El objetivo se resuelve con dos consultas, que reciben el n\'umero de pagar\'e, y son las mismas en las dos instancias. La primera muestra los datos de la inversi\'on, que corresponden a los elementos del pagar\'e. Como devuelve un solo rengl\'on con muchas columnas, se presenta en formato vertical: en PostgreSQL con el comando \nem{\textbackslash x} de \nem{psql}, y en MySQL terminando la sentencia con \nem{\textbackslash G} en lugar del punto y coma. Cuando la inversi\'on no tiene aval, los datos del aval aparecen como \nem{NULL}.

\paragraph{PostgreSQL.}

\begin{verbatim}
	\x on
	SELECT i.inversion_num, i.pagare_num, i.pagare_fecha,
	       s.sucursal_ciudad AS lugar_suscripcion,
	       c.cliente_nombre, c.cliente_apellido,
	       a.aval_nombre, a.aval_apellido,
	       i.inversion_capital, i.inversion_moneda, i.inversion_tasa,
	       i.inversion_tasa_moratoria, i.inversion_dia_pago,
	       i.inversion_fecha_vencimiento, i.pagare_jurisdiccion
	FROM inversion i
	  JOIN sucursal s ON s.sucursal_num = i.sucursal_num
	  JOIN cliente c  ON c.cliente_num = i.cliente_num
	  LEFT JOIN aval a ON a.aval_num = i.aval_num
	WHERE i.pagare_num = '0018/2024';
	\x off
\end{verbatim}

\paragraph{MySQL.}

\begin{verbatim}
	SELECT i.inversion_num, i.pagare_num, i.pagare_fecha,
	       s.sucursal_ciudad AS lugar_suscripcion,
	       c.cliente_nombre, c.cliente_apellido,
	       a.aval_nombre, a.aval_apellido,
	       i.inversion_capital, i.inversion_moneda, i.inversion_tasa,
	       i.inversion_tasa_moratoria, i.inversion_dia_pago,
	       i.inversion_fecha_vencimiento, i.pagare_jurisdiccion
	FROM inversion i
	  JOIN sucursal s ON s.sucursal_num = i.sucursal_num
	  JOIN cliente c  ON c.cliente_num = i.cliente_num
	  LEFT JOIN aval a ON a.aval_num = i.aval_num
	WHERE i.pagare_num = '0018/2024'\G
\end{verbatim}

La segunda consulta, igual en ambos gestores, muestra los periodos de inter\'es cuya fecha final cae dentro del periodo consultado. Cada periodo se identifica por su n\'umero (\nem{interes\us num}) dentro de la inversi\'on, que es tambi\'en la columna con la que se une a sus pagos. Para cada uno indica el n\'umero de pagos, la fecha del \'ultimo, el monto pagado y el saldo. Un periodo est\'a pagado cuando su saldo es cero; si el \'ultimo pago es posterior a su fecha final, se pag\'o con atraso, y si no tiene pagos y su fecha final ya pas\'o, est\'a vencido.

\begin{verbatim}
	SELECT t.interes_num AS periodo,
	       t.interes_fecha_inicial AS inicio, t.interes_fecha_final AS fin,
	       t.interes_monto AS interes,
	       COUNT(g.pago_fecha) AS pagos,
	       MAX(g.pago_fecha)   AS ultimo_pago,
	       COALESCE(SUM(g.pago_interes), 0) AS pagado,
	       t.interes_monto - COALESCE(SUM(g.pago_interes), 0) AS saldo
	FROM inversion i
	  JOIN interes t ON t.inversion_num = i.inversion_num
	  LEFT JOIN pago g
	    ON g.inversion_num = t.inversion_num
	   AND g.interes_num = t.interes_num
	WHERE i.pagare_num = '0018/2024'
	  AND t.interes_fecha_final BETWEEN '2025-06-01' AND '2026-03-31'
	GROUP BY t.inversion_num, t.interes_num, t.interes_fecha_inicial,
	         t.interes_fecha_final, t.interes_monto
	ORDER BY t.interes_num;
\end{verbatim}

El ejemplo corresponde al pagar\'e \nem{0018/2024} (inversi\'on 21), que no tiene aval y est\'a vencido. Entre junio de 2025 y marzo de 2026 tiene diez periodos. El primero, que vence el 22 de junio de 2025, se pag\'o en dos partes y el \'ultimo pago es posterior a su fecha final, por lo que se pag\'o con atraso. Los siete siguientes est\'an pagados, y los dos \'ultimos, vencidos sin pago.

\captura{objetivo4.png}{Formato de las respuestas del objetivo 4}{objetivo4}

\subsection{Objetivo 5}

\textit{Consultar las inversiones hechas en un periodo.}

La consulta lista los pagar\'es suscritos en el periodo, con la sucursal, el cliente, el capital, la tasa mensual, el plazo, la fecha de vencimiento y el estado de cada inversi\'on. Es igual en ambos gestores.

\paragraph{Primera instancia.} El plazo se obtiene contando los periodos de su tabla de inter\'es, ya que cada inversi\'on tiene un periodo por cada mes de plazo, y el estado se determina con un \nem{CASE}.

\begin{verbatim}
	SELECT i.inversion_num, i.pagare_num, i.pagare_fecha AS suscripcion,
	       s.sucursal_nombre,
	       CONCAT(c.cliente_nombre, ' ', c.cliente_apellido) AS cliente,
	       i.inversion_capital AS capital, i.inversion_tasa AS tasa,
	       (SELECT COUNT(*) FROM interes t
	        WHERE t.inversion_num = i.inversion_num) AS plazo_meses,
	       i.inversion_fecha_vencimiento AS vencimiento,
	       CASE WHEN l.inversion_num IS NOT NULL THEN 'Liquidada'
	            WHEN i.inversion_fecha_vencimiento < CURRENT_DATE THEN 'Vencida'
	            ELSE 'Vigente'
	       END AS estado
	FROM inversion i
	  JOIN sucursal s ON s.sucursal_num = i.sucursal_num
	  JOIN cliente c  ON c.cliente_num = i.cliente_num
	  LEFT JOIN liquidacion l ON l.inversion_num = i.inversion_num
	WHERE i.pagare_fecha BETWEEN '2026-01-01' AND '2026-10-28'
	ORDER BY i.pagare_fecha;
\end{verbatim}

\paragraph{Segunda instancia.} El plazo y el estado provienen de la vista \nem{inversion\us estado}.

\begin{verbatim}
	SELECT inversion_num, pagare_num, suscripcion, sucursal_nombre, cliente,
	       capital, tasa, plazo_meses, vencimiento, estado
	FROM inversion_estado
	WHERE suscripcion BETWEEN '2026-01-01' AND '2026-10-28'
	ORDER BY suscripcion;
\end{verbatim}

Del 1 de enero al 28 de octubre de 2026 se suscribieron cinco pagar\'es, las inversiones 36 a 40.

\captura{objetivo5.png}{Formato de la respuesta del objetivo 5}{objetivo5}

	\chapter{Configuraci\'on del entorno}
\label{cap:entorno}
El sistema se implementa en PostgreSQL 17.7 y MySQL 8.0.39. Este cap\'itulo describe la configuraci\'on de las instancias de cada gestor, incluidos los par\'ametros de registro de actividad que sustentan el registro del servidor.

Ambos gestores se instalan en un mismo equipo, cuyas caracter\'isticas se muestran en la tabla~\ref{tab:equipo-servidor}.

\begin{table}[H]
	\centering
	\small
	\caption{Caracter\'isticas del equipo servidor}
	\label{tab:equipo-servidor}
	\begin{tabular}{@{}L{4cm} L{11cm}@{}}
		\toprule
		\textbf{Componente} & \textbf{Caracter\'isticas} \\
		\midrule
		Sistema operativo & Windows 11 de 64 bits \\
		Procesador & AMD Ryzen 7 3700U a 2.30 GHz \\
		Memoria & 8 GB \\
		Almacenamiento & SSD de 447 GB \\
		Nombre del equipo & Armando \\
		Red & Hotspot m\'ovil con asignaci\'on din\'amica de direcciones y sin servidor DNS local \\
		\bottomrule
	\end{tabular}
\end{table}

Estas caracter\'isticas determinan los siguientes criterios que gu\'ian la configuraci\'on de ambos gestores:

\begin{itemize}[leftmargin=*]
	\item \textbf{Memoria limitada y compartida.} El equipo no es un servidor dedicado, por lo que se destinan como m\'aximo alrededor de 2~GB de memoria al conjunto de las cuatro instancias. Dentro de ese l\'imite, la memoria de cada instancia se determina a partir del volumen de datos del sistema.
	\item \textbf{Procesador de cuatro n\'ucleos.} Las cuatro instancias comparten cuatro n\'ucleos, por lo que se evita la ejecuci\'on de consultas en paralelo, que en estas condiciones compite por el procesador en lugar de acelerar el trabajo.
	\item \textbf{Red sin direcciones fijas.} El hotspot asigna las direcciones en cada sesi\'on y no cuenta con DNS, por lo que la configuraci\'on no depende de ninguna direcci\'on concreta ni de la resoluci\'on de nombres de los clientes.
\end{itemize}

La figura~\ref{fig:equipo} muestra las especificaciones del equipo, tal como las reporta el sistema operativo.

\captura[10cm]{equipo.png}{Especificaciones del equipo servidor}{equipo}

El volumen de datos del sistema, estimado a partir del diccionario de datos y del conjunto de datos de prueba, es del orden de decenas de megabytes, incluida la bit\'acora. La tabla~\ref{tab:presupuesto-memoria} muestra la distribuci\'on de memoria que resulta de estos criterios.

\begin{table}[H]
	\centering
	\small
	\caption{Distribuci\'on de memoria entre las instancias}
	\label{tab:presupuesto-memoria}
	\begin{tabular}{@{}L{4.5cm} L{5.5cm} L{4.5cm}@{}}
		\toprule
		\textbf{Instancias} & \textbf{Memoria compartida por instancia} & \textbf{Consumo estimado} \\
		\midrule
		PostgreSQL (2) & 128~MB (\nem{shared\us buffers}) & 550~MB \\
		MySQL (2) & 128~MB (\nem{innodb\us buffer\us pool\us size}) & 650~MB \\
		\midrule
		\textbf{Total} & 512~MB en total (256~MB por gestor) & \textbf{Alrededor de 1.2~GB} \\
		\bottomrule
	\end{tabular}
\end{table}

La memoria compartida corresponde a los valores configurados en cada gestor. El consumo total es una estimación que agrega la memoria de los procesos de fondo y de un número reducido de conexiones simultáneas; su valor real se verifica con el monitor de recursos del sistema operativo una vez que las instancias se encuentran en funcionamiento.

La figura~\ref{fig:memoria} muestra el consumo de memoria medido con las cuatro instancias en funcionamiento.

\captura{memoria.png}{Consumo de memoria de las instancias en el monitor de recursos}{memoria}

\paragraph{Configuraci\'on de la segunda instancia.} Los par\'ametros de este cap\'itulo se aplican con el mismo valor en la segunda instancia de cada gestor, salvo los que identifican a cada instancia, como el puerto y los directorios (cap\'itulo~\ref{cap:multiinstancias}).

\section{PostgreSQL}

\subsection{Directorios}

La primera instancia conserva las rutas que establece el instalador, mostradas en la tabla~\ref{tab:directorios-pg}. Al tratarse de la instancia de referencia, mantener la ubicaci\'on est\'andar facilita su localizaci\'on y conserva el servicio de Windows que el instalador registra. La segunda instancia utiliza directorios propios, ya que dos instancias no pueden compartir el mismo directorio de datos.

\begin{table}[H]
	\centering
	\small
	\caption{Directorios de la primera instancia de PostgreSQL}
	\label{tab:directorios-pg}
	\begin{tabular}{@{}L{2.3cm} L{7.2cm} L{5.6cm}@{}}
		\toprule
		\textbf{Directorio} & \textbf{Ruta} & \textbf{Contenido} \\
		\midrule
		Programas & \nem{C:\bs Program Files\bs PostgreSQL\bs 17\bs bin} & Servidor y herramientas cliente (\nem{psql}, \nem{pg\us dump}) \\
		Datos & \nem{C:\bs Program Files\bs PostgreSQL\bs 17\bs data} & Archivos de la base de datos y de configuraci\'on (\nem{postgresql.conf}, \nem{pg\us hba.conf}) \\
		Registros & \nem{C:\bs Program Files\bs PostgreSQL\bs 17\bs data\bs log} & Archivos de registro del servidor \\
		\bottomrule
	\end{tabular}
\end{table}

Los programas y los datos se encuentran en directorios separados, de modo que una actualizaci\'on del gestor no modifica los datos, y el respaldo de la instancia se limita al directorio de datos.

\subsection{Variables}

\paragraph{Variables de entorno.} Se agrega el directorio de programas a la variable \nem{PATH}, lo que permite ejecutar las herramientas cliente desde cualquier terminal. No se define la variable \nem{PGDATA}: con dos instancias en el mismo equipo, una variable que señala un único directorio de datos haría ambiguo a qué instancia se dirige cada comando. Por ello, el directorio de cada instancia se indica de forma explícita. La variable \nem{PATH} del sistema, con los directorios de programas de ambos gestores, se muestra en la figura~\ref{fig:path}.

\paragraph{Parámetros del servidor.} La tabla~\ref{tab:parametros-pg} muestra los parámetros de \nem{postgresql.conf} que se modifican.

\begin{table}[H]
	\centering
	\small
	\caption{Parámetros modificados en PostgreSQL}
	\label{tab:parametros-pg}
	\begin{tabular}{@{}L{3.4cm} L{1.9cm} L{2.5cm} L{7.6cm}@{}}
		\toprule
		\textbf{Parámetro} & \textbf{Por defecto} & \textbf{Configurado} & \textbf{Justificación} \\
		\midrule
		\nem{effective\us cache\us size} & 4GB & 512MB & No reserva memoria: indica al planificador la caché realmente disponible, para que no elija	planes de ejecución pensados para memoria que el equipo no tiene. \\
		\nem{max\us connections} & 100 & 30 & Cada conexión reserva memoria, y el número de usuarios del sistema no requiere más conexiones. \\
		\nem{max\us parallel\us workers\us per\us gather} & 2 & 0 & Evita la ejecución en paralelo, ya que los cuatro núcleos se comparten entre cuatro	instancias. \\
		\nem{idle\us in\us transaction\us session\us timeout} & 0 & 10min &	Cierra las sesiones que mantienen una transacción abierta sin actividad, que podrían bloquear tablas para el resto de los usuarios. \\
		\nem{timezone} & Del sistema & \nem{America/\allowbreak Mexico\us City} & Fija la hora local en el gestor, sin depender de la configuración del sistema operativo. \\
		\nem{datestyle} & Según el idioma del sistema & \nem{ISO, DMY} & Interpreta las fechas en el orden día, mes y año utilizado en el diccionario de datos (capítulo~\ref{cap:diccionario}). \\
		\bottomrule
	\end{tabular}
\end{table}

Los demás parámetros conservan su valor por defecto por considerarse adecuados para el sistema. Destacan \nem{shared\us buffers} (128~MB), que basta para mantener en memoria el volumen completo de datos; \nem{work\us mem} (4~MB) y \nem{maintenance\us work\us mem} (64~MB), suficientes para las operaciones de ordenamiento y mantenimiento sobre ese volumen; y \nem{autovacuum}, que se mantiene activo para conservar actualizadas las estadísticas y recuperar el espacio de las filas modificadas sin intervención manual. La instancia utiliza la codificación UTF-8, que admite los acentos y la letra ñ presentes en nombres y domicilios.

\paragraph{Parámetros de registro de actividad.} La tabla~\ref{tab:registro-pg} muestra los parámetros que determinan qué registra la instancia en sus archivos de registro, que en este manual se denominan registro del servidor.

\begin{table}[H]
	\centering
	\small
	\caption{Parámetros de registro de actividad en PostgreSQL}
	\label{tab:registro-pg}
	\begin{tabular}{@{}L{3.4cm} L{2.4cm} L{2.6cm} L{7cm}@{}}
		\toprule
		\textbf{Parámetro} & \textbf{Por defecto} & \textbf{Configurado} & \textbf{Justificación} \\
		\midrule
		\nem{logging\us collector} & \nem{on} (instalador) & \nem{on} & Guarda los mensajes del servidor en archivos; sin él, el registro del servidor no se conserva. Se mantiene el valor que establece el instalador. \\
		\nem{log\us directory} & \nem{log} & \nem{log} & La ruta es relativa al directorio de datos, por lo que cada instancia guarda sus registros en su propio directorio sin necesidad de cambiar el valor. \\
		\nem{log\us filename} & \nem{postgresql-\allowbreak\%Y-\%m-\%d\us \%H\%M\%S.log} & \nem{postgresql-\allowbreak\%Y-\%m-\%d.log} & Un archivo por día, identificado por su fecha, en lugar de uno por cada inicio del	servidor. \\
		\nem{log\us statement} & \nem{none} & \nem{mod} & Registra las sentencias que crean o modifican estructuras y datos (DDL y DML). Se excluyen las consultas, que aumentarían el volumen del registro sin	reflejar cambios. \\
		\nem{log\us connections} & \nem{off} & \nem{on} & Registra cada inicio de sesión, con el usuario y el origen. \\
		\nem{log\us disconnections} & \nem{off} & \nem{on} & Registra el cierre	de cada sesión y su duración. \\
		\nem{log\us line\us prefix} & \nem{'\%m [\%p] '} & \nem{'\%m [\%p] \%u@\%d \%r '} & Agrega a cada línea el usuario, la base de datos y el equipo de origen. Sin ellos, el registro indica qué se ejecutó, pero no quién ni desde dónde. \\
		\nem{log\us timezone} & Del sistema & \nem{America/\allowbreak Mexico\us City} & Las horas del registro coinciden con la zona horaria del gestor y con la de MySQL. \\
		\bottomrule
	\end{tabular}
\end{table}

El parámetro \nem{log\us hostname} conserva su valor por defecto (\nem{off}): el registro anota la dirección IP del cliente y no intenta resolver su nombre, lo que en una red sin DNS solo añadiría demoras.

\subsection{Conexi\'on}

La instancia escucha en todas las interfaces de red (\nem{listen\us addresses = '*'}), tanto IPv4 como IPv6, y en el puerto estándar 5432. No se fija una dirección concreta porque el hotspot la asigna en cada sesión. Los clientes se conectan utilizando el nombre del equipo servidor, que se resuelve en la red local sin necesidad de DNS; de esta forma, las conexiones no deben reconfigurarse cuando cambia la dirección. Si la resolución de nombres no está disponible, se utiliza la dirección IP del momento.

El acceso se controla mediante \nem{pg\us hba.conf}, cuyas reglas se evalúan en orden y se aplica la primera que coincide. La tabla~\ref{tab:pg-hba} muestra las reglas de la instancia.

\begin{table}[H]
	\centering
	\small
	\caption{Reglas de acceso de \nem{pg\us hba.conf}}
	\label{tab:pg-hba}
	\begin{tabular}{@{}l l l l l L{4.6cm}@{}}
		\toprule
		\textbf{Tipo} & \textbf{Base} & \textbf{Usuario} & \textbf{Dirección} &	\textbf{Método} & \textbf{Origen} \\
		\midrule
		\nem{host} & \nem{all} & \nem{all} & \nem{127.0.0.1/32} & \nem{scram-sha-256} & Instalador \\
		\nem{host} & \nem{all} & \nem{all} & \nem{::1/128} & \nem{scram-sha-256} & Instalador \\
		\nem{host} & \nem{all} & \nem{postgres} & \nem{samenet} & \nem{reject} & Agregada \\
		\nem{host} & \nem{inversiones} & \nem{all} & \nem{samenet} & \nem{scram-sha-256} & Agregada \\
		\bottomrule
	\end{tabular}
\end{table}

Las dos primeras reglas, creadas por el instalador, permiten las conexiones desde el propio servidor. La tercera impide que el superusuario se conecte desde la red, de modo que la administración de la instancia solo es posible desde el servidor. La cuarta permite que el resto de los usuarios se conecten a la base de datos del sistema desde cualquier equipo de la misma subred que el servidor. La palabra \nem{samenet} se adapta a la subred del momento, por lo que la regla no depende de las direcciones dinámicas del hotspot.

Las contraseñas se almacenan con \nem{scram-sha-256}, el valor por defecto de \nem{password\us encryption} y el método de autenticación más seguro que ofrece el gestor.	

La conexión desde un equipo cliente se realiza con el nombre del equipo servidor (figura~\ref{fig:conexion-remota}). La función \nem{inet\us client\us addr} muestra la dirección desde la que el servidor recibe la conexión:

\begin{verbatim}
	SELECT current_user, inet_client_addr();
\end{verbatim}

\section{MySQL}
\subsection{Directorios}

Al igual que en PostgreSQL, la primera instancia conserva las rutas del instalador, mostradas en la tabla~\ref{tab:directorios-mysql}, y la segunda instancia utiliza directorios propios.

\begin{table}[H]
	\centering
	\small
	\caption{Directorios de la primera instancia de MySQL}
	\label{tab:directorios-mysql}
	\begin{tabular}{@{}L{3.2cm} L{7.2cm} L{5.6cm}@{}}
		\toprule
		\textbf{Directorio} & \textbf{Ruta} & \textbf{Contenido} \\
		\midrule
		Programas & \nem{C:\bs Program Files\bs MySQL\bs MySQL Server 8.0\bs bin} &	Servidor y herramientas cliente (\nem{mysql}, \nem{mysqldump}) \\
		Configuración & \nem{C:\bs ProgramData\bs MySQL\bs MySQL Server 8.0\bs my.ini} & Parámetros del servidor \\
		Datos & \nem{C:\bs ProgramData\bs MySQL\bs MySQL Server 8.0\bs Data} & Archivos de la base de datos y registros del servidor \\
		\bottomrule
	\end{tabular}
\end{table}

En MySQL, el instalador separa los programas, ubicados en \nem{Program Files}, de los datos y la configuración, ubicados en \nem{ProgramData}. Con ello se obtiene la misma ventaja descrita para PostgreSQL: las actualizaciones no afectan los datos y el respaldo se limita a un solo directorio.

\subsection{Variables}

\paragraph{Variables de entorno.} Se agrega el directorio de programas a la variable \nem{PATH} para ejecutar las herramientas cliente desde cualquier terminal. El archivo de configuración de cada instancia se indica de forma explícita, por la misma razón expuesta para PostgreSQL. La figura~\ref{fig:path} muestra la variable \nem{PATH} con los directorios de programas de ambos gestores.

\captura{path.png}{Directorios de programas de ambos gestores en la variable PATH}{path}

\paragraph{Parámetros del servidor.} La tabla~\ref{tab:parametros-mysql} muestra los parámetros de \nem{my.ini} que se modifican.

\begin{table}[H]
	\centering
	\small
	\caption{Parámetros modificados en MySQL}
	\label{tab:parametros-mysql}
	\begin{tabular}{@{}L{3.4cm} L{1.9cm} L{2.5cm} L{7.6cm}@{}}
		\toprule
		\textbf{Parámetro} & \textbf{Por defecto} & \textbf{Configurado} & \textbf{Justificación} \\
		\midrule
		\nem{innodb\us buffer\us pool\us size} & Valor del instalador & 128M & Se fija explícitamente el valor por defecto del gestor, ya que el instalador puede asignar otro según el tipo de equipo elegido. Basta para mantener en memoria el volumen completo de datos. \\
		\nem{max\us connections} & 151 & 30 & Cada conexión reserva memoria, y el número de usuarios del sistema no requiere más conexiones. \\
		\nem{skip\us name\us resolve} & OFF & ON & Evita que el servidor intente resolver el nombre de cada cliente al conectarse. Sin DNS en la red, esa resolución vuelve lentas las conexiones o hace que fallen. Requiere cuentas de \nem{root} para las direcciones de loopback (sección~\ref{sec:conexion-mysql}). \\
		\nem{wait\us timeout} & 28800 (8 h) & 600 (10 min) & Cierra las	conexiones inactivas. Las conexiones a través del hotspot se interrumpen con facilidad y dejan sesiones abiertas que ocupan conexiones disponibles. \\
		\nem{default\us time\us zone} & \nem{SYSTEM} & \nem{'-06:00'} & Fija la	hora local en el gestor. Se utiliza el desplazamiento horario porque en	Windows el gestor no incluye las zonas horarias con nombre, y México ya no aplica horario de verano. \\
		\nem{innodb\us directories} & Vacío & \nem{C:/\allowbreak tablespaces/\allowbreak mysql1} & Declara la ubicación de los tablespaces fuera del directorio de datos; sin ella, el gestor rechaza crearlos en esa ruta (capítulo~\ref{cap:tablespaces}). \\
		\bottomrule
	\end{tabular}
\end{table}

Los demás parámetros conservan su valor por defecto. Destacan el modo SQL estricto (\nem{sql\us mode}), que rechaza los datos inválidos en lugar de truncarlos o ajustarlos, y \nem{character\us set\us server} con el valor \nem{utf8mb4}, que admite acentos y la letra ñ. El parámetro \nem{lower\us case\us table\us names} conserva el valor 1, propio de Windows; como solo puede definirse al inicializar la instancia y todos los nombres del esquema están en minúsculas, no requiere modificarse.

\paragraph{Parámetros de registro de actividad.} La tabla~\ref{tab:registro-mysql} muestra los parámetros que determinan qué registra la instancia en sus archivos de registro, que constituyen su registro del servidor.

\begin{table}[H]
	\centering
	\small
	\caption{Parámetros de registro de actividad en MySQL}
	\label{tab:registro-mysql}
	\begin{tabular}{@{}L{3.4cm} L{2.4cm} L{2.6cm} L{7cm}@{}}
		\toprule
		\textbf{Parámetro} & \textbf{Por defecto} & \textbf{Configurado} & \textbf{Justificación} \\
		\midrule
		\nem{general\us log} & \nem{OFF} & \nem{ON} & Activa el registro general, que anota cada conexión y cada sentencia recibida por el servidor. \\
		\nem{general\us log\us file} & Nombre del equipo & \nem{general.log} & Un nombre fijo y descriptivo. La ruta es relativa al directorio de datos, por lo que cada instancia guarda su propio archivo con el mismo valor. \\
		\nem{log\us output} & \nem{FILE} & \nem{FILE} & El registro se guarda en archivo, igual que en PostgreSQL, y no ocupa espacio dentro de la base de datos. \\
		\nem{log\us timestamps} & \nem{UTC} & \nem{SYSTEM} & Las horas del registro se expresan en la hora local, en lugar de la hora universal, para coincidir con PostgreSQL. \\
		\nem{log\us bin} & \nem{ON} & \nem{ON} & El registro binario conserva los cambios a los datos y permite reconstruirlos. Se mantiene activo, como establece el gestor por defecto. \\
		\nem{log\us bin\us trust\us function\us creators} & \nem{OFF} & \nem{ON} & Con el registro binario activo, el gestor solo permite crear triggers a cuentas con el privilegio \nem{SUPER}. Activarlo permite que el DBA, que no es superusuario, cree los triggers del esquema. \\
		\bottomrule
	\end{tabular}
\end{table}

A diferencia de PostgreSQL, el registro general de MySQL no permite filtrar las sentencias por tipo, por lo que también anota las consultas. Se acepta este volumen adicional porque el número de usuarios es reducido y el equipo cuenta con espacio de almacenamiento suficiente. El registro de errores (\nem{log\us error}) y la vigencia del registro binario (\nem{binlog\us expire\us logs\us seconds}, 30 días) conservan su valor por defecto.

\subsection{Conexi\'on}
\label{sec:conexion-mysql}

El servidor escucha en todas las interfaces de red, el valor por defecto de \nem{bind\us address} (\nem{*}), tanto en IPv4 como en IPv6, y en el puerto estándar 3306. Se conserva este valor porque el nombre del equipo servidor puede resolverse a una dirección IPv6, y escuchar únicamente en IPv4 haría que esas conexiones fallaran. Los clientes se conectan utilizando el nombre del equipo servidor y, si la resolución no está disponible, la dirección IP del momento.

En MySQL, el origen permitido forma parte de cada cuenta de usuario:

\begin{itemize}[leftmargin=*]
	\item \textbf{Superusuario.} Solo existen cuentas de \nem{root} para el propio servidor: \nem{'root'@'localhost'}, creada por el instalador, y \nem{'root'@'127.0.0.1'} y \nem{'root'@'::1'}, que se agregan para las direcciones de loopback de IPv4 e IPv6. De este modo, la administración de la instancia solo es posible desde el servidor.
	\item \textbf{Usuarios de operación.} Se definen con el origen \nem{'\%'}, es decir, desde cualquier dirección. MySQL no cuenta con un equivalente a \nem{samenet}, y la subred del hotspot cambia entre sesiones, por lo que restringir el origen a una subred concreta obligaría a modificar las cuentas en cada sesión. El acceso queda protegido por la autenticación de cada usuario y por los privilegios de su rol, descritos en el capítulo~\ref{cap:usuarios}.
\end{itemize}

\paragraph{Cuentas de loopback del superusuario.} Las cuentas adicionales de \nem{root} son consecuencia de \nem{skip\us name\us resolve}. En Windows, el cliente \nem{mysql} se conecta a \nem{localhost} por TCP/IP, porque la memoria compartida y las canalizaciones con nombre están desactivadas por defecto, de modo que el servidor recibe la conexión desde \nem{127.0.0.1} o \nem{::1}. Sin resolución de nombres, el servidor ya no asocia esas direcciones con el nombre \nem{localhost}, por lo que la cuenta del instalador deja de coincidir y rechaza al superusuario con el error 1130 (\nem{Host '127.0.0.1' is not allowed to connect}). Por ello, las cuentas se crean antes de activar el parámetro, con la misma contraseña y los mismos privilegios que \nem{'root'@'localhost'}:

\begin{verbatim}
	-- Superusuario, antes de activar skip_name_resolve
	CREATE USER 'root'@'127.0.0.1' IDENTIFIED BY '********';
	GRANT ALL PRIVILEGES ON *.* TO 'root'@'127.0.0.1' WITH GRANT OPTION;
	CREATE USER 'root'@'::1' IDENTIFIED BY '********';
	GRANT ALL PRIVILEGES ON *.* TO 'root'@'::1' WITH GRANT OPTION;
\end{verbatim}

\nem{ALL PRIVILEGES} sobre \nem{*.*} incluye los privilegios dinámicos, entre ellos \nem{SYSTEM\us USER}, así que las tres cuentas tienen las mismas facultades. Como ninguna admite conexiones desde otra dirección, \nem{root} sigue sin poder conectarse desde la red.

La autenticación utiliza \nem{caching\us sha2\us password}, el método por defecto del gestor y el más seguro que ofrece.

Además de las cuentas descritas, el gestor tiene sus cuentas internas (\nem{mysql.\allowbreak infoschema}, \nem{mysql.\allowbreak session} y \nem{mysql.\allowbreak sys}), que están bloqueadas y solo existen para \nem{localhost}.

La conexión desde un equipo cliente se realiza con el nombre del equipo servidor (figura~\ref{fig:conexion-remota}). \nem{USER()} indica el origen real de la conexión y \nem{CURRENT\us USER()}, la cuenta con la que se autenticó:

\begin{verbatim}
	SELECT USER(), CURRENT_USER();
\end{verbatim}

\captura{conexion-remota.png}{Conexión remota a ambos gestores por el nombre del servidor}{conexion-remota}
	\chapter{Modelo conceptual}

El modelo conceptual del sistema se compone de nueve entidades: \nem{EMPRESA}, \nem{SUCURSAL}, \nem{CLIENTE}, \nem{AVAL}, \nem{INVERSION}, \nem{INTERES}, \nem{PAGO}, \nem{LIQUIDACION} y \nem{BITACORA}.

\captura{diagrama_conceptual.jpg}{Diagrama conceptual del sistema}{diagrama-conceptual}

\begin{table}[H]
	\centering
	\small
	\caption{Entidades del modelo conceptual y sus llaves}
	\label{tab:entidades}
	\begin{tabular}{@{}l l@{}}
		\toprule
		\textbf{Entidad} & \textbf{Llave} \\
		\midrule
		\nem{EMPRESA}     & PK \nem{empresa\us rfc} \\
		\nem{SUCURSAL}    & PK \nem{sucursal\us num} \\
		\nem{CLIENTE}     & PK \nem{cliente\us num}; AK \nem{cliente\us rfc} \\
		\nem{AVAL}        & PK \nem{aval\us num}; AK \nem{aval\us rfc} \\
		\nem{INVERSION}   & PK \nem{inversion\us num}; AK \nem{pagare\us num} \\
		\nem{INTERES}     & PK (\nem{inversion\us num}, \nem{interes\us fecha\us inicial}) \\
		\nem{PAGO}        & PK (\nem{inversion\us num}, \nem{interes\us fecha\us inicial}, \nem{pago\us fecha}) \\
		\nem{LIQUIDACION} & PK \nem{inversion\us num} \\
		\nem{BITACORA}    & PK (\nem{bitacora\us num}, \nem{bitacora\us fecha}) \\
		\bottomrule
	\end{tabular}
\end{table}

\clearpage

Las relaciones entre las entidades son las siguientes:

\begin{itemize}[leftmargin=*]
	\item \nem{EMPRESA} $\rightarrow$ \nem{SUCURSAL}: cada sucursal pertenece a la empresa (\nem{empresa\us rfc}).
	\item \nem{SUCURSAL} $\rightarrow$ \nem{INVERSION}: cada inversi\'on registra la sucursal donde se suscribi\'o el pagar\'e (\nem{sucursal\us num}).
	\item \nem{CLIENTE} $\rightarrow$ \nem{INVERSION}: cada inversi\'on registra al cliente que invierte (\nem{cliente\us num}).
	\item \nem{AVAL} $\rightarrow$ \nem{INVERSION}: cada inversi\'on registra, si existe, al aval que garantiza el pago (\nem{aval\us num}).
	\item \nem{INVERSION} $\rightarrow$ \nem{INTERES}: cada periodo de inter\'es pertenece a una inversi\'on (\nem{inversion\us num}).
	\item \nem{INTERES} $\rightarrow$ \nem{PAGO}: cada pago cubre un periodo de inter\'es (\nem{inversion\us num}, \nem{interes\us fecha\us inicial}).
	\item \nem{INVERSION} $\rightarrow$ \nem{LIQUIDACION}: la liquidaci\'on registra la devoluci\'on del capital de una inversi\'on (\nem{inversion\us num}).
	\item \nem{BITACORA} no se relaciona con las dem\'as entidades.
\end{itemize}

\paragraph{Reglas de negocio.} Las siguientes reglas provienen del pagaré de inversión y de la operación de la empresa, y se reflejan en el modelo:

\begin{itemize}[leftmargin=*]
	\item El n\'umero de pagar\'e es \'unico y sigue el formato consecutivo/año.
	\item El lugar de suscripci\'on del pagar\'e es la ciudad de la sucursal donde se suscribe.
	\item El aval, cuando existe, garantiza el pago de las obligaciones de la empresa y debe ser una persona distinta del cliente de la misma inversi\'on.
	\item La fecha de vencimiento es posterior a la fecha de suscripci\'on del pagar\'e.
	\item La fecha de vencimiento coincide con el día de pago, de modo que el último pago de intereses se realiza junto con la devolución del capital.
	\item Las tasas ordinaria y moratoria se registran como tasas mensuales. Cuando el pagar\'e pacta una tasa anual, se registra dividida entre doce.
	\item El inter\'es ordinario se paga mensualmente en un d\'ia fijo, entre el 1 y 28 de cada mes, a una tasa fija. Como el capital se devuelve en una sola exhibici\'on, el saldo insoluto sobre el que se calcula el inter\'es es siempre el capital invertido.
	\item El inter\'es moratorio se aplica sobre los saldos vencidos y no pagados, tanto de inter\'es como de capital, a la tasa moratoria establecida en el pagar\'e.
	\item El capital se devuelve en una sola exhibici\'on, por el monto total invertido, en la fecha de vencimiento o despu\'es de ella. Una devoluci\'on posterior al vencimiento genera inter\'es moratorio.
	\item Una inversi\'on se encuentra en uno de tres estados: \textbf{vigente}, si no tiene liquidaci\'on registrada y su fecha de vencimiento no ha pasado; \textbf{vencida}, si no tiene liquidaci\'on registrada y su fecha de vencimiento ya pas\'o; y \textbf{liquidada}, si tiene liquidaci\'on registrada.
	\item La moneda del pagar\'e es, por defecto, el peso mexicano (MXN).
\end{itemize}

\paragraph{Decisiones de diseño.} El modelo adopta las siguientes decisiones:

\begin{itemize}[leftmargin=*]
	\item \textbf{Interés programado separado del pago realizado.} \nem{INTERES} contiene lo que la empresa debe pagar en cada periodo y \nem{PAGO} lo que efectivamente pagó. La diferencia entre ambos permite obtener los estados de cuenta y determinar la mora.
	\item \textbf{Liquidación como entidad propia.} La devolución del capital se registra en \nem{LIQUIDACION}, separada de los pagos de intereses, junto con el interés moratorio que genere su retraso.
	\item \textbf{Empresa a través de la sucursal.} La inversión se relaciona con la sucursal y, por medio de ella, con la empresa, de modo que los datos de la empresa no se repiten en cada inversión.
	\item \textbf{Tasa fija en la inversión.} Como el pagaré establece una tasa fija, las tasas ordinaria y moratoria se almacenan en \nem{INVERSION} en lugar de un catálogo de tasas independiente.
	\item \textbf{Llaves sustitutas.} \nem{SUCURSAL}, \nem{CLIENTE}, \nem{AVAL} e \nem{INVERSION} usan una clave numérica como llave primaria; el RFC y el número de pagaré se conservan como llaves alternas. \nem{EMPRESA} usa su RFC como llave natural.
	\item \textbf{Clave de partición en la llave primaria.} Las llaves primarias de \nem{PAGO} y \nem{BITACORA} incluyen su fecha, porque ambos gestores exigen que la columna de partición forme parte de la llave primaria en la instancia particionada.
	\item \textbf{Bitácora independiente.} \nem{BITACORA} registra los cambios de todas las tablas sin relacionarse con ellas, de modo que sus registros se conservan aunque los datos originales se modifiquen o eliminen.
\end{itemize}
	\chapter{Diccionario de datos}
\label{cap:diccionario}

El diccionario de datos lista los atributos del esquema organizados por tabla. La tabla~\ref{tab:convenciones} describe las convenciones utilizadas.

\begin{table}[H]
	\centering
	\small
	\caption{Convenciones del diccionario de datos}
	\label{tab:convenciones}
	\begin{tabular}{@{}l p{12.5cm}@{}}
		\toprule
		\textbf{Convención} & \textbf{Significado} \\
		\midrule
		Formato & A = letra, 9 = dígito, X = letra o dígito. Los guiones del RFC no se almacenan. \\
		Rango & Longitud máxima en texto, o precisión y escala en decimales. \\
		Llave & PK = primaria, AK = alterna, FK = foránea. \\
		Referencia & Tabla y columna referenciadas; entre paréntesis, la acción al eliminar y al actualizar el registro referenciado. \\
		\bottomrule
	\end{tabular}
\end{table}

\section{Tabla empresa}

\begin{diccionario}{empresa}
	\nem{empresa\us rfc} & RFC de la empresa emisora & VARCHAR & AAA(A)\allowbreak 999999XXX & 13 & No & PK natural & --- \\
	\nem{empresa\us nombre} & Nombre o razón social de la empresa & VARCHAR & AAAA & 100 & No & --- & --- \\
	\nem{empresa\us direccion} & Domicilio fiscal de la empresa & VARCHAR & AAAA & 150 & No & --- & --- \\
\end{diccionario}

\section{Tabla sucursal}

\begin{diccionario}{sucursal}
	\nem{sucursal\us num} & Clave de la sucursal & INT & $>0$, autoincremental & --- & No & PK sustituta & --- \\
	\nem{empresa\us rfc} & Empresa a la que pertenece la sucursal & VARCHAR & AAA(A)\allowbreak 999999XXX & 13 & No & FK & \nem{empresa.\allowbreak empresa\us rfc} (RESTRICT, CASCADE) \\
	\nem{sucursal\us nombre} & Nombre de la sucursal & VARCHAR & AAAA & 60 & No & --- & --- \\
	\nem{sucursal\us ciudad} & Ciudad de la sucursal; lugar de suscripción de sus pagarés & VARCHAR & AAAA & 100 & No & --- & --- \\
	\nem{sucursal\us direccion} & Domicilio de la sucursal & VARCHAR & AAAA & 150 & No & --- & --- \\
\end{diccionario}


\section{Tabla cliente}

\begin{diccionario}{cliente}
	\nem{cliente\us num} & Clave del cliente & INT & $>0$, autoincremental & --- & No & PK sustituta & --- \\
	\nem{cliente\us rfc} & RFC del cliente & CHAR & AAAA999999XXX & 13 & No & AK & --- \\
	\nem{cliente\us nombre} & Nombre(s) del cliente & VARCHAR & AAAA & 60 & No & --- & --- \\
	\nem{cliente\us apellido} & Apellidos del cliente & VARCHAR & AAAA & 80 & No & --- & --- \\
	\nem{cliente\us direccion} & Domicilio del cliente & VARCHAR & AAAA & 150 & No & --- & --- \\
	\nem{cliente\us telefono} & Teléfono del cliente & VARCHAR & 10 a 15 dígitos & 15 & No & --- & --- \\
\end{diccionario}

\section{Tabla aval}

\begin{diccionario}{aval}
	\nem{aval\us num} & Clave del aval & INT & $>0$, autoincremental & --- & No & PK sustituta & --- \\
	\nem{aval\us rfc} & RFC del aval & CHAR & AAAA999999XXX & 13 & No & AK & --- \\
	\nem{aval\us nombre} & Nombre(s) del aval & VARCHAR & AAAA & 60 & No & --- & --- \\
	\nem{aval\us apellido} & Apellidos del aval & VARCHAR & AAAA & 80 & No & --- & --- \\
	\nem{aval\us direccion} & Domicilio del aval & VARCHAR & AAAA & 150 & No & --- & --- \\
\end{diccionario}

\section{Tabla inversion}

\begin{diccionario}{inversion}
	\nem{inversion\us num} & Clave de la inversión & INT & $>0$, autoincremental & --- & No & PK sustituta & --- \\
	\nem{pagare\us num} & Número del pagaré & VARCHAR & 9999/9999 (consecutivo/año) & 30 & No & AK & --- \\
	\nem{sucursal\us num} & Sucursal donde se suscribió el pagaré & INT & $>0$ & --- & No & FK & \nem{sucursal.\allowbreak sucursal\us num} (RESTRICT, RESTRICT) \\
	\nem{cliente\us num} & Cliente que invierte (acreedor) & INT & $>0$ & --- & No & FK & \nem{cliente.\allowbreak cliente\us num} (RESTRICT, RESTRICT) \\
	\nem{aval\us num} & Aval que garantiza el pago de la empresa, si existe & INT & $>0$; RFC distinto del RFC del cliente & --- & Sí & FK & \nem{aval.\allowbreak aval\us num} (RESTRICT, RESTRICT) \\
	\nem{inversion\us capital} & Capital invertido (``bueno por'') & DECIMAL & $>0$, 0.00 & 14,2 & No & --- & --- \\
	\nem{inversion\us tasa} & Tasa mensual de interés ordinario & DECIMAL & 0.0000 a 99.9999 (porcentaje) & 6,4 & No & --- & --- \\
	\nem{inversion\us tasa\us moratoria} & Tasa mensual de interés moratorio & DECIMAL & 0.0000 a 99.9999 (porcentaje) & 6,4 & No & --- & --- \\
	\nem{inversion\us dia\us pago} & Día de pago de cada mes & SMALLINT & 1 a 28 & --- & No & --- & --- \\
	\nem{pagare\us fecha} & Fecha de suscripción del pagaré & DATE & DD/MM/YYYY & --- & No & --- & --- \\
	\nem{inversion\us moneda} & Moneda del pagaré & CHAR & AAA (ISO 4217), por defecto MXN & 3 & No & --- & --- \\
	\nem{pagare\us jurisdiccion} & Ciudad de los tribunales competentes en caso de disputa & VARCHAR & AAAA & 100 & No & --- & --- \\
	\nem{inversion\us fecha\us vencimiento} & Fecha de vencimiento del pagaré & DATE & DD/MM/YYYY, mayor a \nem{pagare\us fecha} & --- & No & --- & --- \\
\end{diccionario}

\section{Tabla interes}

\begin{diccionario}{interes}
	\nem{inversion\us num} & Inversión a la que pertenece el periodo & INT & $>0$ & --- & No & PK compuesta, FK & \nem{inversion.\allowbreak inversion\us num} (RESTRICT, RESTRICT) \\
	\nem{interes\us fecha\us inicial} & Fecha inicial del periodo de interés & DATE & DD/MM/YYYY & --- & No & PK compuesta & --- \\
	\nem{interes\us fecha\us final} & Fecha final del periodo (vence el pago) & DATE & DD/MM/YYYY, mayor a la inicial & --- & No & --- & --- \\
	\nem{interes\us monto} & Interés programado del periodo & DECIMAL & $\geq 0$, 0.00 & 14,2 & No & --- & --- \\
\end{diccionario}

\clearpage

\section{Tabla pago}

\begin{diccionario}{pago}
	\nem{inversion\us num} & Inversión pagada & INT & $>0$ & --- & No & PK compuesta, FK compuesta & \nem{interes} (\nem{inversion\us num}, \nem{interes\us fecha\us inicial}) (RESTRICT, RESTRICT) \\
	\nem{interes\us fecha\us inicial} & Periodo que cubre el pago & DATE & DD/MM/YYYY & --- & No & PK compuesta, FK compuesta & \nem{interes} (\nem{inversion\us num}, \nem{interes\us fecha\us inicial}) (RESTRICT, RESTRICT) \\
	\nem{pago\us fecha} & Fecha del pago (clave de partición) & DATE & DD/MM/YYYY & --- & No & PK compuesta & --- \\
	\nem{pago\us interes} & Monto pagado por interés ordinario & DECIMAL & $\geq 0$, 0.00 & 14,2 & No & --- & --- \\
	\nem{pago\us mora} & Monto pagado por interés moratorio & DECIMAL & $\geq 0$, 0.00, por defecto 0 & 14,2 & No & --- & --- \\
\end{diccionario}

\section{Tabla liquidacion}

\begin{diccionario}{liquidacion}
	\nem{inversion\us num} & Inversión liquidada & INT & $>0$ & --- & No & PK, FK & \nem{inversion.\allowbreak inversion\us num} (RESTRICT, RESTRICT) \\
	\nem{liquidacion\us fecha} & Fecha de devolución del capital & DATE & DD/MM/YYYY, mayor o igual a \nem{inversion\us fecha\us vencimiento} & --- & No & --- & --- \\
	\nem{liquidacion\us capital} & Capital devuelto & DECIMAL & $>0$, 0.00, igual a \nem{inversion\us capital} & 14,2 & No & --- & --- \\
	\nem{liquidacion\us mora} & Interés moratorio pagado por devolución tardía del capital & DECIMAL & $\geq 0$, 0.00, por defecto 0 & 14,2 & No & --- & --- \\
\end{diccionario}

\clearpage

\section{Tabla bitacora}

\begin{diccionario}{bitacora}
	\nem{bitacora\us num} & Número del registro de bitácora & BIGINT & $>0$, autoincremental & --- & No & PK compuesta & --- \\
	\nem{bitacora\us fecha} & Fecha y hora del evento (clave de partición) & TIMESTAMP & DD/MM/YYYY HH:MM:SS & --- & No & PK compuesta & --- \\
	\nem{bitacora\us usuario} & Usuario que inició la sesión (no el dueño ni el definer del trigger) & VARCHAR & AAAA & 100 & No & --- & --- \\
	\nem{bitacora\us tabla} & Tabla afectada por la operación & VARCHAR & AAAA & 63 & No & --- & --- \\
	\nem{bitacora\us operacion} & Tipo de operación ejecutada & VARCHAR & INSERT, UPDATE o DELETE & 6 & No & --- & --- \\
	\nem{bitacora\us ip} & Dirección IP de origen & VARCHAR & IPv4 / IPv6 & 45 & Sí & --- & --- \\
	\nem{bitacora\us sesion} & Identificador de sesión o transacción & BIGINT & $>0$ & --- & Sí & --- & --- \\
	\nem{bitacora\us query} & Sentencia ejecutada & TEXT & AAAA & --- & Sí & --- & --- \\
	\nem{bitacora\us dato\us anterior} & Valores antes del cambio & JSON & JSON & --- & Sí & --- & --- \\
	\nem{bitacora\us dato\us nuevo} & Valores después del cambio & JSON & JSON & --- & Sí & --- & --- \\
\end{diccionario}
	\chapter{Modelo relacional}
\label{cap:relacional}

El modelo relacional representa las tablas del esquema con sus atributos, tipos de dato, llaves y la cardinalidad de cada relaci\'on. Los tipos de dato corresponden al diccionario de datos; su equivalencia en cada gestor se detalla en el cap\'itulo~\ref{cap:implementacion}. El diagrama utiliza la notaci\'on de pata de gallo, cuyos s\'imbolos se describen en la tabla~\ref{tab:notacion-er}.

\captura{diagrama_er.jpg}{Diagrama relacional del sistema}{diagrama-er}

\begin{table}[H]
	\centering
	\small
	\caption{Notaci\'on del diagrama relacional}
	\label{tab:notacion-er}
	\begin{tabular}{@{}L{6.4cm} L{9.6cm}@{}}
		\toprule
		\textbf{Elemento} & \textbf{Significado} \\
		\midrule
		Dos barras & Exactamente uno (1) \\
		C\'irculo y barra & Cero o uno (0..1) \\
		Barra y pata de gallo & Uno o muchos (1..N) \\
		C\'irculo y pata de gallo & Cero o muchos (0..N) \\
		PK, AK, FK & Llave primaria, alterna y for\'anea; los atributos de la llave primaria se muestran en negritas \\
		\nem{NULL} & El atributo acepta valores nulos \\
		\bottomrule
	\end{tabular}
\end{table}

\clearpage

Cada s\'imbolo se lee junto a la tabla en la que se encuentra e indica cu\'antos registros de esa tabla corresponden a un registro de la tabla del otro extremo.

La tabla~\ref{tab:cardinalidades} describe la cardinalidad de cada relaci\'on. La primera cifra corresponde a la tabla referenciada y la segunda a la tabla que contiene la llave for\'anea.

\begin{table}[H]
	\centering
	\small
	\caption{Cardinalidad de las relaciones}
	\label{tab:cardinalidades}
	\begin{tabular}{@{}L{3.3cm} L{2.4cm} L{3.6cm} L{6.7cm}@{}}
		\toprule
		\textbf{Relación} & \textbf{Cardinalidad} & \textbf{Llave foránea} & \textbf{Interpretación} \\
		\midrule
		\nem{empresa} -- \nem{sucursal} & 1 : 1..N & \nem{sucursal.\allowbreak empresa\us rfc} & La empresa opera al menos una sucursal, y cada sucursal pertenece a la empresa. \\
		\nem{sucursal} -- \nem{inversion} & 1 : 0..N & \nem{inversion.\allowbreak sucursal\us num} & Cada inversión se suscribe en una sola sucursal; una sucursal puede no tener inversiones. \\
		\nem{cliente} -- \nem{inversion} & 1 : 0..N & \nem{inversion.\allowbreak cliente\us num} & Cada inversión pertenece a un solo cliente; un cliente puede tener varias inversiones o ninguna. \\
		\nem{aval} -- \nem{inversion} & 0..1 : 0..N & \nem{inversion.\allowbreak aval\us num} & Una inversión puede tener un aval o ninguno, por lo que la llave foránea acepta nulos; un aval puede garantizar varias inversiones. \\
		\nem{inversion} -- \nem{interes} & 1 : 1..N & \nem{interes.\allowbreak inversion\us num} & Cada inversión tiene un periodo de interés por cada mes de su plazo. \\
		\nem{interes} -- \nem{pago} & 1 : 0..N & \nem{pago} (\nem{inversion\us num}, \nem{interes\us num}) & Un periodo de interés puede estar pendiente de pago o cubrirse con uno o varios pagos. \\
		\nem{inversion} -- \nem{liquidacion} & 1 : 0..1 & \nem{liquidacion.\allowbreak inversion\us num} & Una inversión se liquida una sola vez, por lo que la llave foránea es también la llave primaria; mientras no se liquida, no tiene registro. \\
		\bottomrule
	\end{tabular}
\end{table}

La tabla \nem{bitacora} no se relaciona con las dem\'as tablas, ya que registra los cambios de todas ellas sin depender de su existencia.

Las llaves for\'aneas garantizan el m\'aximo de cada relaci\'on y que todo registro referenciado exista. Los m\'inimos de uno en el lado de la tabla que contiene la llave for\'anea, es decir, que la empresa tenga al menos una sucursal, y que cada inversi\'on tenga su tabla de inter\'es, se garantizan mediante la operaci\'on del sistema: el gerente registra las sucursales de la empresa, y el procedimiento \nem{generar\us tabla\us interes} crea los periodos de cada inversi\'on al darla de alta.
	\chapter{Usuarios y roles}
\recursoscapitulo{06-usuarios}

\section{Roles y matriz de privilegios}
\label{sec:roles-matriz}

Los roles del sistema corresponden a puestos reales que pueden existir en una empresa que emita pagar\'es de inversi\'on, en lugar de categorizarse \'unicamente por su capacidad global de consulta, inserci\'on o eliminaci\'on. Cada rol agrupa los privilegios de un puesto, y los usuarios individuales obtienen sus privilegios al recibir el rol correspondiente. Los clientes y los avales no cuentan con roles, ya que no acceden directamente a la base de datos.

\begin{table}[H]
	\centering
	\small
	\caption{Roles del sistema}
	\label{tab:roles}
	\begin{tabular}{@{}l L{3.4cm} L{8.2cm}@{}}
		\toprule
		\textbf{Rol} & \textbf{Puesto} & \textbf{Función en el sistema} \\
		\midrule
		\nem{rol\us admin\us bd}  & DBA & Responsable del esquema; crea y administra sus objetos y no opera el negocio. \\
		\nem{rol\us gerente}    & Gerente de inversiones & Supervisa, corrige inversiones y administra los datos de la empresa, sus sucursales y los avales. \\
		\nem{rol\us ejecutivo}  & Asesor o promotor & Da de alta clientes e inversiones y genera la tabla de interés. \\
		\nem{rol\us tesoreria}  & Cajero o tesorería & Registra los pagos mensuales y la liquidación del capital. \\
		\nem{rol\us auditor}    & Contralor o auditor & Solo lectura de todo, incluida la bitácora. \\
		\nem{rol\us respaldo}   & Cuenta técnica & Solo los privilegios que necesitan \nem{pg\us dump} y \nem{mysqldump}. \\
		\bottomrule
	\end{tabular}
\end{table}


La tabla~\ref{tab:matriz} muestra los privilegios de cada rol sobre las tablas del esquema (S = SELECT, I = INSERT, U = UPDATE, D = DELETE, X = EXECUTE).

\begin{table}[H]
	\centering
	\small
	\caption{Matriz de privilegios por rol y tabla}
	\label{tab:matriz}
	\begin{tabular}{@{}l *{9}{c}@{}}
		\toprule
		\textbf{Rol} &
		\nem{empresa} & \nem{sucursal} & \nem{cliente} & \nem{aval} & \nem{inversion} & \nem{interes} & \nem{pago} & \nem{liquidacion} & \nem{bitacora} \\
		\midrule
		\nem{rol\us admin\us bd}\textsuperscript{4} & Todos & Todos & Todos & Todos & Todos & Todos & Todos & Todos & Todos \\
		\nem{rol\us gerente}   & S U & S I U & S U & S I U & S U D\textsuperscript{1} & S & S & S & --- \\
		\nem{rol\us ejecutivo} & S & S & S I U\textsuperscript{2} & S & S I & S X\textsuperscript{3} & S & S & --- \\
		\nem{rol\us tesoreria} & S & S & S & S & S & S & S I & S I & --- \\
		\nem{rol\us auditor}   & S & S & S & S & S & S & S & S & S \\
		\nem{rol\us respaldo}\textsuperscript{5} & S & S & S & S & S & S & S & S & S \\
		\bottomrule
	\end{tabular}
\end{table}

\begin{small}
	\begin{enumerate}[leftmargin=*,label=\textsuperscript{\arabic*}]
		\item Solo puede eliminar inversiones sin intereses generados; la llave foránea con \nem{RESTRICT} lo impide en los demás casos.
		\item UPDATE por columna: únicamente teléfono y dirección. El RFC y el nombre no se modifican.
		\item El ejecutivo no inserta directamente en \nem{interes}: ejecuta el procedimiento \nem{generar\us tabla\us interes}, que recibe la clave de la inversión y calcula los periodos a partir del capital, la tasa, el día de pago y el vencimiento.
		\item Como responsable del esquema, tiene todos los privilegios sobre sus objetos, incluido DDL. La forma en que los obtiene en cada gestor se describe en la sección~\ref{sec:admin-instancia}.
		\item Además de SELECT, recibe los privilegios adicionales que requiere la herramienta de respaldo de cada gestor; se detallan en las secciones~\ref{sec:roles-pg} y~\ref{sec:roles-mysql}.
	\end{enumerate}
\end{small}

\paragraph{Bit\'acora.} Ning\'un rol de operaci\'on tiene privilegios de INSERT, UPDATE o DELETE sobre la tabla \nem{bitacora}; \'unicamente los triggers escriben en ella. El DBA, como responsable del esquema, conserva control t\'ecnico sobre la tabla; por ello, sus acciones quedan registradas adem\'as en la bit\'acora de instancia, cuya configuraci\'on solo puede modificar el superusuario.

\paragraph{Separaci\'on de funciones.} Quien da de alta la inversi\'on (\nem{rol\us ejecutivo}) no puede registrar pagos ni liquidaciones, y quien los registra (\nem{rol\us tesoreria}) no puede modificar las inversiones, por lo que el capital y las tasas quedan fuera de su alcance. Del mismo modo, el DBA administra la estructura de la base de datos, pero no participa en la operaci\'on del negocio.	

\paragraph{Usuarios.} El personal accede directamente a la base de datos, sin una aplicación intermedia. Por ello, cada integrante del equipo cuenta con un usuario propio en el gestor, que recibe el rol de su puesto; de esta forma, la bitácora identifica a la persona que realizó cada operación y no solo el puesto desde el que la realizó. Si una persona cambia de puesto, basta con cambiar su rol, sin crear otra cuenta ni perder su historial. La única excepción es el usuario de respaldo, que no corresponde a una persona sino a la herramienta que genera los respaldos. La tabla~\ref{tab:usuarios} muestra los usuarios del sistema, con el mismo nombre en ambos gestores.

\begin{table}[H]
	\centering
	\small
	\caption{Usuarios del sistema}
	\label{tab:usuarios}
	\begin{tabular}{@{}l l L{7cm}@{}}
		\toprule
		\textbf{Usuario} & \textbf{Rol} & \textbf{Titular} \\
		\midrule
		\nem{dlandero} & \nem{rol\us admin\us bd} & Daniel Landero Arias \\
		\nem{dmejia}   & \nem{rol\us gerente}   & Diego de Jesús Mejía Rodríguez \\
		\nem{pzuniga}  & \nem{rol\us ejecutivo} & Pedro Armando Zúñiga Durán \\
		\nem{xzaleta}  & \nem{rol\us tesoreria} & Ximena Zaleta Hernández \\
		\nem{enovoa}   & \nem{rol\us auditor}   & Eduardo Novoa Siles \\
		\nem{respaldo} & \nem{rol\us respaldo}  & Cuenta técnica \\
		\bottomrule
	\end{tabular}
\end{table}

Los nombres de usuario se forman con la inicial del nombre y el primer apellido, sin acentos ni la letra ñ, para que sean válidos y se escriban igual en ambos gestores.

\section{Administraci\'on de la instancia y de la base de datos}
\label{sec:admin-instancia}

La administraci\'on del sistema se divide en dos niveles. El superusuario de cada gestor (\nem{postgres} en PostgreSQL y \nem{root} en MySQL) administra la instancia, y el DBA (\nem{rol\us admin\us bd}) administra la base de datos y sus objetos. El superusuario no se utiliza para crear ni operar los objetos del esquema: su participaci\'on se limita a las tareas que requieran privilegios sobre la instancia completa.

\begin{table}[H]
	\centering
	\small
	\caption{Responsabilidades por nivel de administración}
	\label{tab:responsabilidades}
	\begin{tabular}{@{}L{3.6cm} L{6.2cm} L{6.2cm}@{}}
		\toprule
		& \textbf{Superusuario (instancia)} & \textbf{DBA (base de datos)} \\
		\midrule
		Cuenta & \nem{postgres} / \nem{root} & \nem{rol\us admin\us bd} \\
		Alcance & Toda la instancia y todas sus bases de datos & Únicamente la base de datos del sistema \\
		Instancias & Crea y configura las instancias, incluida la segunda instancia de cada gestor & --- \\
		Configuración & Parámetros del servidor y bitácora de instancia & --- \\
		Tablespaces & Crea los tablespaces y, en PostgreSQL, los asigna al DBA & Asigna las tablas a los tablespaces \\
		Base de datos & Crea el rol del DBA y le otorga control sobre la base de datos & Crea el esquema, las tablas, los índices, las particiones y la segmentación \\
		Objetos programados & --- & Crea los procedimientos, funciones, triggers y vistas \\
		Roles y usuarios & Crea únicamente al DBA & Crea los roles y usuarios de operación y les otorga sus privilegios \\
		Datos & --- & Carga los datos de prueba y los traslada a la segunda instancia \\
		\bottomrule
	\end{tabular}
\end{table}

\paragraph{Propiedad de los objetos.} Los dos gestores manejan la propiedad de los objetos de forma distinta. En PostgreSQL cada objeto tiene dueño, que es el rol que lo crea o al que se le asigna con \nem{OWNER}. En MySQL las tablas no tienen dueño: pertenecen a la base de datos, y el control sobre ellas se obtiene mediante privilegios. Solo los procedimientos, triggers y vistas tienen un equivalente, el \nem{DEFINER}, que es la cuenta con cuyos privilegios se ejecutan.

\begin{table}[H]
	\centering
	\small
	\caption{Propiedad de los objetos en cada gestor}
	\label{tab:propiedad}
	\begin{tabular}{@{}L{3.6cm} L{6.2cm} L{6.2cm}@{}}
		\toprule
		\textbf{Objeto} & \textbf{PostgreSQL} & \textbf{MySQL} \\
		\midrule
		Instancia & Cuenta del servicio y superusuario \nem{postgres} & Cuenta del servicio y superusuario \nem{root} \\
		Tablespaces & Creados por \nem{postgres}, con dueño \nem{rol\us admin\us bd} & Creados por \nem{root}; no tienen dueño \\
		Base de datos & Dueño \nem{rol\us admin\us bd} & Sin dueño; \nem{rol\us admin\us bd} tiene todos los privilegios sobre ella \\
		Esquema & Dueño \nem{rol\us admin\us bd} & Equivale a la base de datos \\
		Tablas & Dueño \nem{rol\us admin\us bd} & Sin dueño; controlados mediante los privilegios del DBA \\
		Procedimientos y funciones & Dueño \nem{rol\us admin\us bd} & \nem{DEFINER}: cuenta del DBA \\
		Triggers & Su función tiene dueño \nem{rol\us admin\us bd} & \nem{DEFINER}: cuenta del DBA \\
		Vistas & Dueño \nem{rol\us admin\us bd} & \nem{DEFINER}: cuenta del DBA \\
		Rol del DBA & Creado por \nem{postgres} & Creado por \nem{root} \\
		Roles y usuarios de operación & Creados por el DBA & Creados por el DBA \\
		\bottomrule
	\end{tabular}
\end{table}

En PostgreSQL, el usuario de la persona que funge como DBA trabaja con \nem{SET ROLE rol\us admin\us bd}, de modo que los objetos quedan a nombre del rol y no de la persona. Si el DBA cambia, basta con otorgar el rol al nuevo usuario sin transferir la propiedad de ning\'un objeto. En MySQL, el \nem{DEFINER} debe ser una cuenta de usuario y no un rol, por lo que los objetos programados se definen con la cuenta del DBA.

\section{Creaci\'on de roles y usuarios en PostgreSQL}
\label{sec:roles-pg}

En PostgreSQL, roles y usuarios son el mismo tipo de objeto; la diferencia es que un usuario tiene el atributo \nem{LOGIN}. Los roles de puesto se crean sin él, de modo que nadie puede iniciar sesión directamente como un puesto: solo a través de su usuario personal.

\paragraph{Ejecutado por el superusuario.} De acuerdo con la sección~\ref{sec:admin-instancia}, el superusuario crea únicamente el rol del DBA y su usuario. El rol recibe el atributo \nem{CREATEROLE}, que le permite crear los demás roles y usuarios sin ser superusuario.

\begin{verbatim}
	CREATE ROLE rol_admin_bd NOLOGIN CREATEROLE;
	CREATE USER dlandero WITH PASSWORD 'Dlandero#2026' IN ROLE rol_admin_bd;
\end{verbatim}

\paragraph{Ejecutado por el DBA.} Los atributos de un rol, como \nem{CREATEROLE}, no se heredan por pertenencia; se obtienen al asumir el rol con \nem{SET ROLE}. Esto coincide con la forma de trabajo descrita en la sección~\ref{sec:admin-instancia}: los roles creados quedan a nombre de \nem{rol\us admin\us bd} y no de la persona, y el gestor le concede automáticamente la facultad de asignarlos.

\begin{verbatim}
	SET ROLE rol_admin_bd;
	
	CREATE ROLE rol_gerente   NOLOGIN;
	CREATE ROLE rol_ejecutivo NOLOGIN;
	CREATE ROLE rol_tesoreria NOLOGIN;
	CREATE ROLE rol_auditor   NOLOGIN;
	CREATE ROLE rol_respaldo  NOLOGIN;
	
	CREATE USER dmejia   WITH PASSWORD 'Dmejia#2026'   IN ROLE rol_gerente;
	CREATE USER pzuniga  WITH PASSWORD 'Pzuniga#2026'  IN ROLE rol_ejecutivo;
	CREATE USER xzaleta  WITH PASSWORD 'Xzaleta#2026'  IN ROLE rol_tesoreria;
	CREATE USER enovoa   WITH PASSWORD 'Enovoa#2026'   IN ROLE rol_auditor;
	CREATE USER respaldo WITH PASSWORD 'Respaldo#2026' IN ROLE rol_respaldo;
	
	RESET ROLE;
\end{verbatim}

Los privilegios de cada rol sobre las tablas, incluidos los de \nem{rol\us respaldo}, se otorgan en la implementación física, una vez que las tablas existen. En PostgreSQL, \nem{pg\us dump} solo requiere lectura sobre las tablas y las secuencias del esquema, por lo que el usuario de respaldo no
necesita privilegios sobre la instancia.

Los roles, sus atributos y la pertenencia de cada usuario se consultan en el catálogo del gestor (figura~\ref{fig:pg-roles}):

\begin{verbatim}
	SELECT u.rolname AS usuario, r.rolname AS rol,
	       u.rolcanlogin AS inicia_sesion, r.rolcreaterole AS crea_roles
	FROM pg_auth_members m
	  JOIN pg_roles u ON u.oid = m.member
	  JOIN pg_roles r ON r.oid = m.roleid
	WHERE r.rolname LIKE 'rol\_%'
	ORDER BY r.rolname;
\end{verbatim}

\captura[12cm]{pg-roles.png}{Roles y usuarios creados en PostgreSQL}{pg-roles}

\section{Creaci\'on de roles y usuarios en MySQL}
\label{sec:roles-mysql}

En MySQL, los roles y los usuarios son objetos distintos, y los roles concedidos a un usuario permanecen inactivos hasta que se activan. Por ello, cada usuario se crea con su rol como rol por defecto (\nem{DEFAULT ROLE}), de modo que sus privilegios están disponibles desde el inicio de la sesión. Todos los usuarios se definen con el origen \nem{'\%'}, por las razones expuestas en la configuración del entorno.

\paragraph{Ejecutado por el superusuario.} El superusuario crea el rol del DBA, le otorga el control total sobre la base de datos \nem{inversiones}, que en MySQL equivale a ser su dueño, y los privilegios necesarios para administrar roles y usuarios. MySQL permite otorgar privilegios sobre una base de datos que aún no existe, por lo que este paso puede realizarse antes de crearla.

\begin{verbatim}
	CREATE ROLE rol_admin_bd;
	GRANT ALL PRIVILEGES ON inversiones.* TO rol_admin_bd WITH GRANT OPTION;
	GRANT CREATE ROLE, CREATE USER, ROLE_ADMIN ON *.* TO rol_admin_bd;

	CREATE USER 'dlandero'@'%' IDENTIFIED BY 'Dlandero#2026'
	  DEFAULT ROLE rol_admin_bd;
\end{verbatim}

El privilegio \nem{CREATE USER} permite modificar cualquier cuenta. Sin embargo, \nem{root} posee el privilegio \nem{SYSTEM\us USER}, y una cuenta con ese privilegio solo puede ser modificada por otra que también lo tenga; como el DBA no lo posee, no puede alterar la cuenta del superusuario.

\paragraph{Ejecutado por el DBA.} El DBA crea los roles de puesto y los
usuarios del personal:

\begin{verbatim}
	CREATE ROLE rol_gerente, rol_ejecutivo, rol_tesoreria,
	            rol_auditor, rol_respaldo;

	CREATE USER 'dmejia'@'%'   IDENTIFIED BY 'Dmejia#2026'
	  DEFAULT ROLE rol_gerente;
	CREATE USER 'pzuniga'@'%'  IDENTIFIED BY 'Pzuniga#2026'
	  DEFAULT ROLE rol_ejecutivo;
	CREATE USER 'xzaleta'@'%'  IDENTIFIED BY 'Xzaleta#2026'
	  DEFAULT ROLE rol_tesoreria;
	CREATE USER 'enovoa'@'%'   IDENTIFIED BY 'Enovoa#2026'
	  DEFAULT ROLE rol_auditor;
	CREATE USER 'respaldo'@'%' IDENTIFIED BY 'Respaldo#2026'
	  DEFAULT ROLE rol_respaldo;
\end{verbatim}

\paragraph{Privilegios de respaldo.} A diferencia de PostgreSQL, \nem{mysqldump} requiere privilegios adicionales a la lectura de datos. Los que corresponden a la base de datos los otorga el DBA; los que corresponden a toda la instancia los otorga el superusuario, de acuerdo con la división de la sección~\ref{sec:admin-instancia}:

\begin{verbatim}
	-- DBA: privilegios sobre la base de datos
	GRANT SELECT, SHOW VIEW, TRIGGER, LOCK TABLES
	  ON inversiones.* TO rol_respaldo;

	-- Superusuario: privilegios sobre la instancia
	GRANT PROCESS, SHOW_ROUTINE ON *.* TO rol_respaldo;
\end{verbatim}

\nem{SHOW VIEW} y \nem{TRIGGER} permiten respaldar la definición de las vistas y los triggers; \nem{LOCK TABLES}, bloquear las tablas para obtener una copia consistente; \nem{PROCESS}, leer la información de los tablespaces; y \nem{SHOW\us ROUTINE}, respaldar los procedimientos almacenados.

Los roles concedidos y los roles por defecto de cada usuario se consultan en las
tablas del sistema (figura~\ref{fig:mysql-roles}):

\begin{verbatim}
	SELECT TO_USER AS usuario, FROM_USER AS rol
	FROM mysql.role_edges
	ORDER BY FROM_USER;

	SELECT USER AS usuario, DEFAULT_ROLE_USER AS rol_por_defecto
	FROM mysql.default_roles;
\end{verbatim}

\captura{mysql-roles.png}{Roles y usuarios creados en MySQL}{mysql-roles}
	\chapter{Implementaci\'on f\'isica}
\recursoscapitulo{07-implementacion}
El esquema se implementa de la misma manera en ambos gestores: las mismas tablas, columnas, llaves y restricciones, con los mismos nombres. Las diferencias se limitan a la sintaxis y a los tipos de dato propios de cada gestor. Este cap\'itulo describe las decisiones comunes a ambas implementaciones.

\paragraph{Equivalencia de tipos.} El diccionario de datos utiliza tipos genéricos. La tabla~\ref{tab:equivalencia-tipos} muestra su equivalencia en cada gestor.

\begin{table}[H]
	\centering
	\small
	\caption{Equivalencia de tipos de dato entre gestores}
	\label{tab:equivalencia-tipos}
	\begin{tabular}{@{}L{3.3cm} L{3.9cm} L{3.1cm} L{5.5cm}@{}}
		\toprule
		\textbf{Diccionario} & \textbf{PostgreSQL} & \textbf{MySQL} & \textbf{Justificación} \\
		\midrule
		INT y BIGINT autoincremental & \nem{GENERATED BY DEFAULT AS IDENTITY} & \nem{AUTO\us INCREMENT} & El gestor genera la clave, pero admite insertar un valor explícito, lo que permite trasladar los datos a la segunda instancia conservando sus claves. \\
		DECIMAL(p,s) & \nem{NUMERIC(p,s)} & \nem{DECIMAL(p,s)} & Valores exactos, sin errores de redondeo en montos y tasas. \\
		TIMESTAMP & \nem{TIMESTAMP} & \nem{DATETIME} & Fecha y hora local sin conversión. En MySQL, \nem{TIMESTAMP} convierte los valores a UTC y solo admite fechas hasta 2038. \\
		JSON & \nem{JSONB} & \nem{JSON} & Ambos validan que el contenido sea JSON. En PostgreSQL, \nem{JSONB} lo almacena en formato binario, más eficiente para consultarlo. \\
		VARCHAR, CHAR, SMALLINT, DATE, TEXT & Igual & Igual & Tipos estándar con el mismo comportamiento en ambos gestores. \\
		\bottomrule
	\end{tabular}
\end{table}

\paragraph{Nomenclatura de restricciones.} Todas las restricciones se nombran explícitamente con el prefijo de su tipo y la tabla a la que pertenecen: \nem{pk\us} para llaves primarias, \nem{ak\us} para llaves alternas, \nem{fk\us} para llaves foráneas y \nem{ck\us} para reglas de validación. Así, cuando una operación viola una restricción, el mensaje de error identifica la tabla y la regla incumplida, y los nombres son los mismos en ambos gestores.

\paragraph{Validación de las reglas de negocio.} Las reglas que involucran columnas de una sola tabla se implementan como restricciones \nem{CHECK}, que el gestor evalúa en cada inserción y actualización sin código adicional. Las reglas que comparan datos de dos tablas no pueden expresarse con \nem{CHECK}, por lo que se implementan con triggers. La tabla~\ref{tab:validacion-reglas} muestra la distribución.

\begin{table}[H]
	\centering
	\small
	\caption{Mecanismo de validación de cada regla de negocio}
	\label{tab:validacion-reglas}
	\begin{tabular}{@{}L{2.2cm} L{13.6cm}@{}}
		\toprule
		\textbf{Mecanismo} & \textbf{Reglas} \\
		\midrule
		\nem{CHECK} & Capital mayor a cero; tasas no negativas; día de pago entre 1 y 28; vencimiento posterior a la suscripción; vencimiento en el día de pago; fecha final del periodo posterior a la inicial; montos de interés, pago, mora y liquidación no negativos; operación de bitácora válida. \\
		Trigger & RFC del aval distinto del RFC del cliente; capital liquidado igual al invertido; liquidación en la fecha de vencimiento o después. \\
		\bottomrule
	\end{tabular}
\end{table}

El límite superior de las tasas (99.9999) no requiere una restricción, ya que lo impone la precisión del tipo \nem{NUMERIC(6,4)}.

\section{PostgreSQL}

\subsection{Base de datos}

La base de datos la crea el superusuario y la asigna al DBA como dueño, de acuerdo con la sección~\ref{sec:admin-instancia}. Se crea a partir de \nem{template0}, la plantilla sin modificaciones, para garantizar que la codificación UTF-8 se aplique sin heredar objetos de otra plantilla.

\begin{verbatim}
	-- Superusuario
	CREATE DATABASE inversiones
	  OWNER rol_admin_bd
	  ENCODING 'UTF8'
	  TEMPLATE template0;
\end{verbatim}

Dentro de la base de datos, el DBA crea el esquema \nem{inversiones}, donde residen todos los objetos del sistema, y lo establece como ruta de búsqueda predeterminada para que las sentencias no requieran el prefijo del esquema. El esquema \nem{public}, que el gestor crea por defecto, se elimina para que ningún objeto pueda crearse por error fuera del esquema del sistema.

\begin{verbatim}
	-- DBA, conectado a la base de datos inversiones
	SET ROLE rol_admin_bd;
	
	CREATE SCHEMA inversiones AUTHORIZATION rol_admin_bd;
	ALTER DATABASE inversiones SET search_path TO inversiones;
	DROP SCHEMA public;
\end{verbatim}

\begin{verbatim}
	\l inversiones
	\dn
\end{verbatim}

\captura{pg-base-datos.png}{Base de datos y esquema del sistema en PostgreSQL}{pg-base-datos}

\subsection{Tablas}
Las tablas se crean en el orden de sus dependencias: primero las tablas referenciadas y después las que contienen las llaves foráneas. El DBA las crea con el rol \nem{rol\us admin\us bd} activo, de modo que el rol es su dueño.

\begin{verbatim}
	SET ROLE rol_admin_bd;

	CREATE TABLE empresa (
	  empresa_rfc        VARCHAR(13)  NOT NULL,
	  empresa_nombre     VARCHAR(100) NOT NULL,
	  empresa_direccion  VARCHAR(150) NOT NULL,
	  CONSTRAINT pk_empresa PRIMARY KEY (empresa_rfc)
	);

	CREATE TABLE sucursal (
	  sucursal_num        INT GENERATED BY DEFAULT AS IDENTITY,
	  empresa_rfc         VARCHAR(13)  NOT NULL,
	  sucursal_nombre     VARCHAR(60)  NOT NULL,
	  sucursal_ciudad     VARCHAR(100) NOT NULL,
	  sucursal_direccion  VARCHAR(150) NOT NULL,
	  CONSTRAINT pk_sucursal PRIMARY KEY (sucursal_num),
	  CONSTRAINT fk_sucursal_empresa FOREIGN KEY (empresa_rfc)
	    REFERENCES empresa (empresa_rfc)
	    ON DELETE RESTRICT ON UPDATE CASCADE
	);

	CREATE TABLE cliente (
	  cliente_num        INT GENERATED BY DEFAULT AS IDENTITY,
	  cliente_rfc        CHAR(13)     NOT NULL,
	  cliente_nombre     VARCHAR(60)  NOT NULL,
	  cliente_apellido   VARCHAR(80)  NOT NULL,
	  cliente_direccion  VARCHAR(150) NOT NULL,
	  cliente_telefono   VARCHAR(15)  NOT NULL,
	  CONSTRAINT pk_cliente PRIMARY KEY (cliente_num),
	  CONSTRAINT ak_cliente_rfc UNIQUE (cliente_rfc)
	);

	CREATE TABLE aval (
	  aval_num        INT GENERATED BY DEFAULT AS IDENTITY,
	  aval_rfc        CHAR(13)     NOT NULL,
	  aval_nombre     VARCHAR(60)  NOT NULL,
	  aval_apellido   VARCHAR(80)  NOT NULL,
	  aval_direccion  VARCHAR(150) NOT NULL,
	  CONSTRAINT pk_aval PRIMARY KEY (aval_num),
	  CONSTRAINT ak_aval_rfc UNIQUE (aval_rfc)
	);

	CREATE TABLE inversion (
	  inversion_num                INT GENERATED BY DEFAULT AS IDENTITY,
	  pagare_num                   VARCHAR(30)   NOT NULL,
	  sucursal_num                 INT           NOT NULL,
	  cliente_num                  INT           NOT NULL,
	  aval_num                     INT           NULL,
	  inversion_capital            NUMERIC(14,2) NOT NULL,
	  inversion_tasa               NUMERIC(6,4)  NOT NULL,
	  inversion_tasa_moratoria     NUMERIC(6,4)  NOT NULL,
	  inversion_dia_pago           SMALLINT      NOT NULL,
	  pagare_fecha                 DATE          NOT NULL,
	  inversion_moneda             CHAR(3)       NOT NULL DEFAULT 'MXN',
	  pagare_jurisdiccion          VARCHAR(100)  NOT NULL,
	  inversion_fecha_vencimiento  DATE          NOT NULL,
	  CONSTRAINT pk_inversion PRIMARY KEY (inversion_num),
	  CONSTRAINT ak_inversion_pagare UNIQUE (pagare_num),
	  CONSTRAINT fk_inversion_sucursal FOREIGN KEY (sucursal_num)
	    REFERENCES sucursal (sucursal_num)
	    ON DELETE RESTRICT ON UPDATE RESTRICT,
	  CONSTRAINT fk_inversion_cliente FOREIGN KEY (cliente_num)
	    REFERENCES cliente (cliente_num)
	    ON DELETE RESTRICT ON UPDATE RESTRICT,
	  CONSTRAINT fk_inversion_aval FOREIGN KEY (aval_num)
	    REFERENCES aval (aval_num)
	    ON DELETE RESTRICT ON UPDATE RESTRICT,
	  CONSTRAINT ck_inversion_capital CHECK (inversion_capital > 0),
	  CONSTRAINT ck_inversion_tasa CHECK (inversion_tasa >= 0),
	  CONSTRAINT ck_inversion_tasa_moratoria
	    CHECK (inversion_tasa_moratoria >= 0),
	  CONSTRAINT ck_inversion_dia_pago
	    CHECK (inversion_dia_pago BETWEEN 1 AND 28),
	  CONSTRAINT ck_inversion_vencimiento
	    CHECK (inversion_fecha_vencimiento > pagare_fecha),
	  CONSTRAINT ck_inversion_vencimiento_dia
	    CHECK (EXTRACT(DAY FROM inversion_fecha_vencimiento)
	           = inversion_dia_pago)
	);

	CREATE TABLE interes (
	  inversion_num          INT           NOT NULL,
	  interes_fecha_inicial  DATE          NOT NULL,
	  interes_fecha_final    DATE          NOT NULL,
	  interes_monto          NUMERIC(14,2) NOT NULL,
	  CONSTRAINT pk_interes PRIMARY KEY (inversion_num, interes_fecha_inicial),
	  CONSTRAINT fk_interes_inversion FOREIGN KEY (inversion_num)
	    REFERENCES inversion (inversion_num)
	    ON DELETE RESTRICT ON UPDATE RESTRICT,
	  CONSTRAINT ck_interes_fechas
	    CHECK (interes_fecha_final > interes_fecha_inicial),
	  CONSTRAINT ck_interes_monto CHECK (interes_monto >= 0)
	);

	CREATE TABLE pago (
	  inversion_num          INT           NOT NULL,
	  interes_fecha_inicial  DATE          NOT NULL,
	  pago_fecha             DATE          NOT NULL,
	  pago_interes           NUMERIC(14,2) NOT NULL,
	  pago_mora              NUMERIC(14,2) NOT NULL DEFAULT 0,
	  CONSTRAINT pk_pago
	    PRIMARY KEY (inversion_num, interes_fecha_inicial, pago_fecha),
	  CONSTRAINT fk_pago_interes FOREIGN KEY (inversion_num, interes_fecha_inicial)
	    REFERENCES interes (inversion_num, interes_fecha_inicial)
	    ON DELETE RESTRICT ON UPDATE RESTRICT,
	  CONSTRAINT ck_pago_montos CHECK (pago_interes >= 0 AND pago_mora >= 0)
	);

	CREATE TABLE liquidacion (
	  inversion_num        INT           NOT NULL,
	  liquidacion_fecha    DATE          NOT NULL,
	  liquidacion_capital  NUMERIC(14,2) NOT NULL,
	  liquidacion_mora     NUMERIC(14,2) NOT NULL DEFAULT 0,
	  CONSTRAINT pk_liquidacion PRIMARY KEY (inversion_num),
	  CONSTRAINT fk_liquidacion_inversion FOREIGN KEY (inversion_num)
	    REFERENCES inversion (inversion_num)
	    ON DELETE RESTRICT ON UPDATE RESTRICT,
	  CONSTRAINT ck_liquidacion_capital CHECK (liquidacion_capital > 0),
	  CONSTRAINT ck_liquidacion_mora CHECK (liquidacion_mora >= 0)
	);

	CREATE TABLE bitacora (
	  bitacora_num            BIGINT GENERATED BY DEFAULT AS IDENTITY,
	  bitacora_fecha          TIMESTAMP    NOT NULL DEFAULT LOCALTIMESTAMP,
	  bitacora_usuario        VARCHAR(100) NOT NULL,
	  bitacora_tabla          VARCHAR(63)  NOT NULL,
	  bitacora_operacion      VARCHAR(6)   NOT NULL,
	  bitacora_ip             VARCHAR(45)  NULL,
	  bitacora_sesion         BIGINT       NULL,
	  bitacora_query          TEXT         NULL,
	  bitacora_dato_anterior  JSONB        NULL,
	  bitacora_dato_nuevo     JSONB        NULL,
	  CONSTRAINT pk_bitacora PRIMARY KEY (bitacora_num, bitacora_fecha),
	  CONSTRAINT ck_bitacora_operacion
	    CHECK (bitacora_operacion IN ('INSERT', 'UPDATE', 'DELETE'))
	);
\end{verbatim}

La fecha de la bitácora se asigna con \nem{LOCALTIMESTAMP}, que devuelve la hora local sin zona horaria, en lugar de \nem{CURRENT\us TIMESTAMP}, que la incluye y obligaría a convertirla para almacenarla en una columna \nem{TIMESTAMP}.

Las tablas y sus restricciones se consultan en el catálogo del gestor (figura~\ref{fig:pg-tablas}):

\begin{verbatim}
	SELECT conrelid::regclass AS tabla, conname AS restriccion,
	       CASE contype WHEN 'p' THEN 'PK' WHEN 'u' THEN 'AK'
	                    WHEN 'f' THEN 'FK' WHEN 'c' THEN 'CHECK' END AS tipo
	FROM pg_constraint
	WHERE connamespace = 'inversiones'::regnamespace
	ORDER BY 1, 3, 2;
\end{verbatim}

\captura{pg-tablas.png}{Tablas y restricciones del esquema en PostgreSQL}{pg-tablas}

\subsection{Objetos programados}

\paragraph{Procedimiento \nem{generar\us tabla\us interes}.} Genera los periodos de interés de una inversión a partir de las condiciones de su pagaré. Se implementa como procedimiento almacenado, y no como un cálculo que realice cada usuario, para que todas las tablas de interés se generen con la misma regla:

\begin{itemize}[leftmargin=*]
	\item El monto de cada periodo es fijo e igual a capital $\times$ tasa / 100, ya que el pagaré establece un importe mensual fijo.
	\item El primer periodo inicia en la fecha de suscripción y termina en el día de pago del mes siguiente; los demás van de un día de pago al siguiente.
	\item El último periodo termina en la fecha de vencimiento, que por regla coincide con el día de pago.
\end{itemize}

El procedimiento rechaza una inversión que ya tenga tabla de interés, para no duplicar periodos, y una cuyo vencimiento sea anterior al primer día de pago, que no tendría ningún periodo completo.

Se define con \nem{SECURITY DEFINER}, es decir, se ejecuta con los privilegios de su dueño. De esta forma, el ejecutivo puede generar los periodos sin tener privilegio de inserción sobre \nem{interes}. Por seguridad, el procedimiento fija su propia ruta de búsqueda, de modo que no puede ser engañado con objetos del mismo nombre creados en otro esquema.

\begin{verbatim}
	CREATE PROCEDURE generar_tabla_interes(p_inversion_num INT)
	LANGUAGE plpgsql
	SECURITY DEFINER
	SET search_path = inversiones
	AS $$
	DECLARE
	  v        inversion%ROWTYPE;
	  v_monto  NUMERIC(14,2);
	  v_inicio DATE;
	  v_fin    DATE;
	BEGIN
	  SELECT * INTO v FROM inversion WHERE inversion_num = p_inversion_num;
	  IF NOT FOUND THEN
	    RAISE EXCEPTION 'La inversión % no existe', p_inversion_num;
	  END IF;

	  IF EXISTS (SELECT 1 FROM interes WHERE inversion_num = p_inversion_num) THEN
	    RAISE EXCEPTION 'La inversión % ya tiene tabla de interés',
	      p_inversion_num;
	  END IF;

	  v_monto  := round(v.inversion_capital * v.inversion_tasa / 100, 2);
	  v_inicio := v.pagare_fecha;
	  v_fin    := (date_trunc('month', v.pagare_fecha) + INTERVAL '1 month'
	               + (v.inversion_dia_pago - 1) * INTERVAL '1 day')::DATE;

	  IF v_fin > v.inversion_fecha_vencimiento THEN
	    RAISE EXCEPTION 'El vencimiento es anterior al primer día de pago';
	  END IF;

	  WHILE v_fin <= v.inversion_fecha_vencimiento LOOP
	    INSERT INTO interes (inversion_num, interes_fecha_inicial,
	                         interes_fecha_final, interes_monto)
	    VALUES (p_inversion_num, v_inicio, v_fin, v_monto);
	    v_inicio := v_fin;
	    v_fin    := (v_fin + INTERVAL '1 month')::DATE;
	  END LOOP;
	END;
	$$;
\end{verbatim}

Como el día de pago es como máximo 28, sumar un mes a la fecha nunca produce un día inexistente, como el 30 de febrero.

\paragraph{Triggers de validación.} Las reglas que comparan dos tablas se validan antes de cada inserción o actualización. Si la regla no se cumple, el trigger cancela la operación con un mensaje que explica el motivo. A diferencia del procedimiento, estas funciones se ejecutan con los privilegios del usuario que realiza la operación, ya que solo consultan tablas que el ejecutivo y la tesorería pueden leer.

\begin{verbatim}
	CREATE FUNCTION fn_validar_aval() RETURNS trigger
	LANGUAGE plpgsql
	SET search_path = inversiones
	AS $$
	BEGIN
	  IF NEW.aval_num IS NOT NULL AND
	     (SELECT aval_rfc FROM aval WHERE aval_num = NEW.aval_num) =
	     (SELECT cliente_rfc FROM cliente WHERE cliente_num = NEW.cliente_num)
	  THEN
	    RAISE EXCEPTION 'El aval no puede ser el cliente de la inversión';
	  END IF;
	  RETURN NEW;
	END;
	$$;

	CREATE TRIGGER trg_inversion_aval
	  BEFORE INSERT OR UPDATE OF aval_num, cliente_num ON inversion
	  FOR EACH ROW EXECUTE FUNCTION fn_validar_aval();

	CREATE FUNCTION fn_validar_liquidacion() RETURNS trigger
	LANGUAGE plpgsql
	SET search_path = inversiones
	AS $$
	DECLARE
	  v_capital      NUMERIC(14,2);
	  v_vencimiento  DATE;
	BEGIN
	  SELECT inversion_capital, inversion_fecha_vencimiento
	    INTO v_capital, v_vencimiento
	  FROM inversion
	  WHERE inversion_num = NEW.inversion_num;

	  IF NEW.liquidacion_capital <> v_capital THEN
	    RAISE EXCEPTION 'El capital liquidado debe ser igual al invertido (%)',
	      v_capital;
	  END IF;

	  IF NEW.liquidacion_fecha < v_vencimiento THEN
	    RAISE EXCEPTION 'La liquidación no puede ser anterior al vencimiento (%)',
	      v_vencimiento;
	  END IF;
	  RETURN NEW;
	END;
	$$;

	CREATE TRIGGER trg_liquidacion_validar
	  BEFORE INSERT OR UPDATE ON liquidacion
	  FOR EACH ROW EXECUTE FUNCTION fn_validar_liquidacion();
\end{verbatim}

\begin{verbatim}
	SELECT p.proname AS objeto,
	       CASE p.prokind WHEN 'p' THEN 'Procedimiento' ELSE 'Función' END AS tipo,
	       pg_get_userbyid(p.proowner) AS dueno,
	       p.prosecdef AS security_definer
	FROM pg_proc p
	WHERE p.pronamespace = 'inversiones'::regnamespace;

	SELECT tgname AS trigger, tgrelid::regclass AS tabla
	FROM pg_trigger
	WHERE NOT tgisinternal;
\end{verbatim}

\captura{pg-objetos.png}{Procedimiento, funciones y triggers en PostgreSQL}{pg-objetos}

\subsection{Permisos}

Los privilegios de la matriz de la tabla~\ref{tab:matriz} se otorgan a los roles, nunca a los usuarios, de modo que un cambio de puesto solo requiere cambiar el rol de la persona. Antes de otorgarlos, se retiran los privilegios que PostgreSQL concede por defecto a todos los usuarios (\nem{PUBLIC}): la conexión a la base de datos y la ejecución de funciones y procedimientos. Sin este paso, cualquier usuario de la instancia podría conectarse a la base de datos o ejecutar el procedimiento, aunque no tuviera un rol del sistema.

\begin{verbatim}
	SET ROLE rol_admin_bd;

	-- Privilegios por defecto de PUBLIC
	REVOKE CONNECT, TEMPORARY ON DATABASE inversiones FROM PUBLIC;
	REVOKE EXECUTE ON PROCEDURE generar_tabla_interes(INT) FROM PUBLIC;
	REVOKE EXECUTE ON FUNCTION fn_validar_aval(), fn_validar_liquidacion()
	  FROM PUBLIC;

	-- Acceso a la base de datos y al esquema
	GRANT CONNECT ON DATABASE inversiones
	  TO rol_gerente, rol_ejecutivo, rol_tesoreria, rol_auditor, rol_respaldo;
	GRANT USAGE ON SCHEMA inversiones
	  TO rol_gerente, rol_ejecutivo, rol_tesoreria, rol_auditor, rol_respaldo;

	-- Gerente
	GRANT SELECT, UPDATE         ON empresa, cliente           TO rol_gerente;
	GRANT SELECT, INSERT, UPDATE ON sucursal, aval             TO rol_gerente;
	GRANT SELECT, UPDATE, DELETE ON inversion                  TO rol_gerente;
	GRANT SELECT                 ON interes, pago, liquidacion TO rol_gerente;
	GRANT USAGE ON SEQUENCE sucursal_sucursal_num_seq, aval_aval_num_seq
	  TO rol_gerente;

	-- Ejecutivo
	GRANT SELECT ON empresa, sucursal, aval, interes, pago, liquidacion
	  TO rol_ejecutivo;
	GRANT SELECT, INSERT ON cliente, inversion TO rol_ejecutivo;
	GRANT UPDATE (cliente_telefono, cliente_direccion) ON cliente
	  TO rol_ejecutivo;
	GRANT USAGE ON SEQUENCE cliente_cliente_num_seq, inversion_inversion_num_seq
	  TO rol_ejecutivo;
	GRANT EXECUTE ON PROCEDURE generar_tabla_interes(INT) TO rol_ejecutivo;

	-- Tesorería
	GRANT SELECT ON empresa, sucursal, cliente, aval, inversion, interes
	  TO rol_tesoreria;
	GRANT SELECT, INSERT ON pago, liquidacion TO rol_tesoreria;

	-- Auditor y respaldo
	GRANT SELECT ON ALL TABLES IN SCHEMA inversiones TO rol_auditor, rol_respaldo;
	GRANT SELECT ON ALL SEQUENCES IN SCHEMA inversiones TO rol_respaldo;

	-- Tablas creadas posteriormente
	ALTER DEFAULT PRIVILEGES FOR ROLE rol_admin_bd IN SCHEMA inversiones
	  GRANT SELECT ON TABLES TO rol_auditor, rol_respaldo;
\end{verbatim}

El privilegio \nem{USAGE} sobre las secuencias de las columnas de identidad permite a los roles que insertan registros generar sus claves. El último bloque establece privilegios por defecto: toda tabla que el DBA cree después, como las particiones de la segunda instancia, será legible automáticamente por el auditor y por el usuario de respaldo, sin necesidad de otorgar el privilegio en cada ocasión.

Los privilegios otorgados se consultan con el comando \nem{\textbackslash dp} (figura~\ref{fig:pg-permisos}):

\begin{verbatim}
	\dp inversiones.*
\end{verbatim}

\captura{pg-permisos.png}{Privilegios sobre las tablas en PostgreSQL}{pg-permisos}

\section{MySQL}

\subsection{Base de datos}

En MySQL, la base de datos y el esquema son el mismo objeto, por lo que no se crea un esquema adicional. El DBA crea la base de datos, ya que recibió todos los privilegios sobre \nem{inversiones} desde antes de que existiera (sección~\ref{sec:roles-mysql}).

La base de datos utiliza la codificación \nem{utf8mb4} y la intercalación \nem{utf8mb4\us 0900\us ai\us ci}, que compara los textos sin distinguir mayúsculas, minúsculas ni acentos. Así, una búsqueda por nombre encuentra a ``Jose'' aunque esté registrado como ``José''. PostgreSQL, en cambio, distingue mayúsculas y acentos; las consultas por nombre de ambos gestores se escriben de modo que devuelvan los mismos resultados (sección~\ref{sec:consultas-objetivos}).

\begin{verbatim}
	-- DBA
	CREATE DATABASE inversiones
	  CHARACTER SET utf8mb4
	  COLLATE utf8mb4_0900_ai_ci;
\end{verbatim}

\begin{verbatim}
	SELECT SCHEMA_NAME, DEFAULT_CHARACTER_SET_NAME, DEFAULT_COLLATION_NAME
	FROM information_schema.SCHEMATA
	WHERE SCHEMA_NAME = 'inversiones';
\end{verbatim}

\captura{mysql-base-datos.png}{Base de datos del sistema en MySQL}{mysql-base-datos}

\subsection{Tablas}

Las tablas tienen las mismas columnas, llaves y restricciones que en PostgreSQL, con los tipos equivalentes de la tabla~\ref{tab:equivalencia-tipos}. Todas utilizan el motor InnoDB, el único de MySQL que admite llaves foráneas y transacciones; aunque es el motor por defecto, se indica de forma explícita para que la definición no dependa de la configuración del servidor.

\begin{verbatim}
	USE inversiones;

	CREATE TABLE empresa (
	  empresa_rfc        VARCHAR(13)  NOT NULL,
	  empresa_nombre     VARCHAR(100) NOT NULL,
	  empresa_direccion  VARCHAR(150) NOT NULL,
	  CONSTRAINT pk_empresa PRIMARY KEY (empresa_rfc)
	) ENGINE = InnoDB;

	CREATE TABLE sucursal (
	  sucursal_num        INT          NOT NULL AUTO_INCREMENT,
	  empresa_rfc         VARCHAR(13)  NOT NULL,
	  sucursal_nombre     VARCHAR(60)  NOT NULL,
	  sucursal_ciudad     VARCHAR(100) NOT NULL,
	  sucursal_direccion  VARCHAR(150) NOT NULL,
	  CONSTRAINT pk_sucursal PRIMARY KEY (sucursal_num),
	  CONSTRAINT fk_sucursal_empresa FOREIGN KEY (empresa_rfc)
	    REFERENCES empresa (empresa_rfc)
	    ON DELETE RESTRICT ON UPDATE CASCADE
	) ENGINE = InnoDB;

	CREATE TABLE cliente (
	  cliente_num        INT          NOT NULL AUTO_INCREMENT,
	  cliente_rfc        CHAR(13)     NOT NULL,
	  cliente_nombre     VARCHAR(60)  NOT NULL,
	  cliente_apellido   VARCHAR(80)  NOT NULL,
	  cliente_direccion  VARCHAR(150) NOT NULL,
	  cliente_telefono   VARCHAR(15)  NOT NULL,
	  CONSTRAINT pk_cliente PRIMARY KEY (cliente_num),
	  CONSTRAINT ak_cliente_rfc UNIQUE (cliente_rfc)
	) ENGINE = InnoDB;

	CREATE TABLE aval (
	  aval_num        INT          NOT NULL AUTO_INCREMENT,
	  aval_rfc        CHAR(13)     NOT NULL,
	  aval_nombre     VARCHAR(60)  NOT NULL,
	  aval_apellido   VARCHAR(80)  NOT NULL,
	  aval_direccion  VARCHAR(150) NOT NULL,
	  CONSTRAINT pk_aval PRIMARY KEY (aval_num),
	  CONSTRAINT ak_aval_rfc UNIQUE (aval_rfc)
	) ENGINE = InnoDB;

	CREATE TABLE inversion (
	  inversion_num                INT           NOT NULL AUTO_INCREMENT,
	  pagare_num                   VARCHAR(30)   NOT NULL,
	  sucursal_num                 INT           NOT NULL,
	  cliente_num                  INT           NOT NULL,
	  aval_num                     INT           NULL,
	  inversion_capital            DECIMAL(14,2) NOT NULL,
	  inversion_tasa               DECIMAL(6,4)  NOT NULL,
	  inversion_tasa_moratoria     DECIMAL(6,4)  NOT NULL,
	  inversion_dia_pago           SMALLINT      NOT NULL,
	  pagare_fecha                 DATE          NOT NULL,
	  inversion_moneda             CHAR(3)       NOT NULL DEFAULT 'MXN',
	  pagare_jurisdiccion          VARCHAR(100)  NOT NULL,
	  inversion_fecha_vencimiento  DATE          NOT NULL,
	  CONSTRAINT pk_inversion PRIMARY KEY (inversion_num),
	  CONSTRAINT ak_inversion_pagare UNIQUE (pagare_num),
	  CONSTRAINT fk_inversion_sucursal FOREIGN KEY (sucursal_num)
	    REFERENCES sucursal (sucursal_num)
	    ON DELETE RESTRICT ON UPDATE RESTRICT,
	  CONSTRAINT fk_inversion_cliente FOREIGN KEY (cliente_num)
	    REFERENCES cliente (cliente_num)
	    ON DELETE RESTRICT ON UPDATE RESTRICT,
	  CONSTRAINT fk_inversion_aval FOREIGN KEY (aval_num)
	    REFERENCES aval (aval_num)
	    ON DELETE RESTRICT ON UPDATE RESTRICT,
	  CONSTRAINT ck_inversion_capital CHECK (inversion_capital > 0),
	  CONSTRAINT ck_inversion_tasa CHECK (inversion_tasa >= 0),
	  CONSTRAINT ck_inversion_tasa_moratoria
	    CHECK (inversion_tasa_moratoria >= 0),
	  CONSTRAINT ck_inversion_dia_pago
	    CHECK (inversion_dia_pago BETWEEN 1 AND 28),
	  CONSTRAINT ck_inversion_vencimiento
	    CHECK (inversion_fecha_vencimiento > pagare_fecha),
	  CONSTRAINT ck_inversion_vencimiento_dia
	    CHECK (DAY(inversion_fecha_vencimiento) = inversion_dia_pago)
	) ENGINE = InnoDB;

	CREATE TABLE interes (
	  inversion_num          INT           NOT NULL,
	  interes_fecha_inicial  DATE          NOT NULL,
	  interes_fecha_final    DATE          NOT NULL,
	  interes_monto          DECIMAL(14,2) NOT NULL,
	  CONSTRAINT pk_interes PRIMARY KEY (inversion_num, interes_fecha_inicial),
	  CONSTRAINT fk_interes_inversion FOREIGN KEY (inversion_num)
	    REFERENCES inversion (inversion_num)
	    ON DELETE RESTRICT ON UPDATE RESTRICT,
	  CONSTRAINT ck_interes_fechas
	    CHECK (interes_fecha_final > interes_fecha_inicial),
	  CONSTRAINT ck_interes_monto CHECK (interes_monto >= 0)
	) ENGINE = InnoDB;

	CREATE TABLE pago (
	  inversion_num          INT           NOT NULL,
	  interes_fecha_inicial  DATE          NOT NULL,
	  pago_fecha             DATE          NOT NULL,
	  pago_interes           DECIMAL(14,2) NOT NULL,
	  pago_mora              DECIMAL(14,2) NOT NULL DEFAULT 0,
	  CONSTRAINT pk_pago
	    PRIMARY KEY (inversion_num, interes_fecha_inicial, pago_fecha),
	  CONSTRAINT fk_pago_interes FOREIGN KEY (inversion_num, interes_fecha_inicial)
	    REFERENCES interes (inversion_num, interes_fecha_inicial)
	    ON DELETE RESTRICT ON UPDATE RESTRICT,
	  CONSTRAINT ck_pago_montos CHECK (pago_interes >= 0 AND pago_mora >= 0)
	) ENGINE = InnoDB;

	CREATE TABLE liquidacion (
	  inversion_num        INT           NOT NULL,
	  liquidacion_fecha    DATE          NOT NULL,
	  liquidacion_capital  DECIMAL(14,2) NOT NULL,
	  liquidacion_mora     DECIMAL(14,2) NOT NULL DEFAULT 0,
	  CONSTRAINT pk_liquidacion PRIMARY KEY (inversion_num),
	  CONSTRAINT fk_liquidacion_inversion FOREIGN KEY (inversion_num)
	    REFERENCES inversion (inversion_num)
	    ON DELETE RESTRICT ON UPDATE RESTRICT,
	  CONSTRAINT ck_liquidacion_capital CHECK (liquidacion_capital > 0),
	  CONSTRAINT ck_liquidacion_mora CHECK (liquidacion_mora >= 0)
	) ENGINE = InnoDB;

	CREATE TABLE bitacora (
	  bitacora_num            BIGINT       NOT NULL AUTO_INCREMENT,
	  bitacora_fecha          DATETIME     NOT NULL DEFAULT CURRENT_TIMESTAMP,
	  bitacora_usuario        VARCHAR(100) NOT NULL,
	  bitacora_tabla          VARCHAR(63)  NOT NULL,
	  bitacora_operacion      VARCHAR(6)   NOT NULL,
	  bitacora_ip             VARCHAR(45)  NULL,
	  bitacora_sesion         BIGINT       NULL,
	  bitacora_query          TEXT         NULL,
	  bitacora_dato_anterior  JSON         NULL,
	  bitacora_dato_nuevo     JSON         NULL,
	  CONSTRAINT pk_bitacora PRIMARY KEY (bitacora_num, bitacora_fecha),
	  CONSTRAINT ck_bitacora_operacion
	    CHECK (bitacora_operacion IN ('INSERT', 'UPDATE', 'DELETE'))
	) ENGINE = InnoDB;
\end{verbatim}

En MySQL, \nem{CURRENT\us TIMESTAMP} asignado a una columna \nem{DATETIME} devuelve la hora local según la zona horaria del servidor (\nem{default\us time\us zone}), por lo que cumple la misma función que \nem{LOCALTIMESTAMP} en PostgreSQL.

Las restricciones de cada tabla se consultan en el diccionario de datos del gestor (figura~\ref{fig:mysql-tablas}):

\begin{verbatim}
	SELECT TABLE_NAME, CONSTRAINT_NAME, CONSTRAINT_TYPE
	FROM information_schema.TABLE_CONSTRAINTS
	WHERE CONSTRAINT_SCHEMA = 'inversiones'
	ORDER BY TABLE_NAME, CONSTRAINT_TYPE, CONSTRAINT_NAME;
\end{verbatim}

\captura{mysql-tablas.png}{Tablas y restricciones del esquema en MySQL}{mysql-tablas}

\subsection{Objetos programados}

\paragraph{Procedimiento \nem{generar\us tabla\us interes}.} Aplica la misma regla que en PostgreSQL. Se define con \nem{SQL SECURITY DEFINER}, equivalente a \nem{SECURITY DEFINER}, y su \nem{DEFINER} es la cuenta del DBA que lo crea, de acuerdo con la sección~\ref{sec:admin-instancia}.

\begin{verbatim}
	DELIMITER //

	CREATE DEFINER = CURRENT_USER PROCEDURE generar_tabla_interes(
	  IN p_inversion_num INT)
	SQL SECURITY DEFINER
	MODIFIES SQL DATA
	BEGIN
	  DECLARE v_capital     DECIMAL(14,2);
	  DECLARE v_tasa        DECIMAL(6,4);
	  DECLARE v_dia         SMALLINT;
	  DECLARE v_suscripcion DATE;
	  DECLARE v_vencimiento DATE;
	  DECLARE v_monto       DECIMAL(14,2);
	  DECLARE v_inicio      DATE;
	  DECLARE v_fin         DATE;

	  SELECT inversion_capital, inversion_tasa, inversion_dia_pago,
	         pagare_fecha, inversion_fecha_vencimiento
	    INTO v_capital, v_tasa, v_dia, v_suscripcion, v_vencimiento
	  FROM inversion
	  WHERE inversion_num = p_inversion_num;

	  IF v_capital IS NULL THEN
	    SIGNAL SQLSTATE '45000'
	      SET MESSAGE_TEXT = 'La inversión no existe';
	  END IF;

	  IF EXISTS (SELECT 1 FROM interes
	             WHERE inversion_num = p_inversion_num) THEN
	    SIGNAL SQLSTATE '45000'
	      SET MESSAGE_TEXT = 'La inversión ya tiene tabla de interés';
	  END IF;

	  SET v_monto  = ROUND(v_capital * v_tasa / 100, 2);
	  SET v_inicio = v_suscripcion;
	  -- Primer día del mes siguiente, más los días hasta el día de pago
	  SET v_fin = DATE_ADD(
	                DATE_ADD(DATE_SUB(v_suscripcion,
	                                  INTERVAL DAY(v_suscripcion) - 1 DAY),
	                         INTERVAL 1 MONTH),
	                INTERVAL v_dia - 1 DAY);

	  IF v_fin > v_vencimiento THEN
	    SIGNAL SQLSTATE '45000'
	      SET MESSAGE_TEXT = 'El vencimiento es anterior al primer día de pago';
	  END IF;

	  WHILE v_fin <= v_vencimiento DO
	    INSERT INTO interes (inversion_num, interes_fecha_inicial,
	                         interes_fecha_final, interes_monto)
	    VALUES (p_inversion_num, v_inicio, v_fin, v_monto);
	    SET v_inicio = v_fin;
	    SET v_fin    = DATE_ADD(v_fin, INTERVAL 1 MONTH);
	  END WHILE;
	END //

	DELIMITER ;
\end{verbatim}

\paragraph{Validaciones entre tablas.} MySQL no permite que un trigger responda a más de un tipo de operación, por lo que cada regla requiere un trigger para la inserción y otro para la actualización. Para no duplicar la lógica, cada regla se escribe una sola vez en un procedimiento, y ambos triggers lo invocan. A diferencia de PostgreSQL, los triggers de MySQL siempre se ejecutan con los privilegios de su \nem{DEFINER}; como estos procedimientos solo consultan datos, esto no otorga a los usuarios ninguna capacidad adicional.

\begin{verbatim}
	DELIMITER //

	CREATE DEFINER = CURRENT_USER PROCEDURE validar_aval(
	  IN p_cliente_num INT, IN p_aval_num INT)
	READS SQL DATA
	BEGIN
	  IF p_aval_num IS NOT NULL AND
	     (SELECT aval_rfc FROM aval WHERE aval_num = p_aval_num) =
	     (SELECT cliente_rfc FROM cliente WHERE cliente_num = p_cliente_num)
	  THEN
	    SIGNAL SQLSTATE '45000'
	      SET MESSAGE_TEXT = 'El aval no puede ser el cliente de la inversión';
	  END IF;
	END //

	CREATE DEFINER = CURRENT_USER PROCEDURE validar_liquidacion(
	  IN p_inversion_num INT, IN p_fecha DATE, IN p_capital DECIMAL(14,2))
	READS SQL DATA
	BEGIN
	  DECLARE v_capital     DECIMAL(14,2);
	  DECLARE v_vencimiento DATE;

	  SELECT inversion_capital, inversion_fecha_vencimiento
	    INTO v_capital, v_vencimiento
	  FROM inversion
	  WHERE inversion_num = p_inversion_num;

	  IF p_capital <> v_capital THEN
	    SIGNAL SQLSTATE '45000'
	      SET MESSAGE_TEXT = 'El capital liquidado debe ser igual al invertido';
	  END IF;

	  IF p_fecha < v_vencimiento THEN
	    SIGNAL SQLSTATE '45000'
	      SET MESSAGE_TEXT = 'La liquidación no puede ser anterior al vencimiento';
	  END IF;
	END //

	CREATE DEFINER = CURRENT_USER TRIGGER trg_inversion_aval_ins
	  BEFORE INSERT ON inversion FOR EACH ROW
	BEGIN
	  CALL validar_aval(NEW.cliente_num, NEW.aval_num);
	END //

	CREATE DEFINER = CURRENT_USER TRIGGER trg_inversion_aval_upd
	  BEFORE UPDATE ON inversion FOR EACH ROW
	BEGIN
	  CALL validar_aval(NEW.cliente_num, NEW.aval_num);
	END //

	CREATE DEFINER = CURRENT_USER TRIGGER trg_liquidacion_validar_ins
	  BEFORE INSERT ON liquidacion FOR EACH ROW
	BEGIN
	  CALL validar_liquidacion(NEW.inversion_num, NEW.liquidacion_fecha,
	                           NEW.liquidacion_capital);
	END //

	CREATE DEFINER = CURRENT_USER TRIGGER trg_liquidacion_validar_upd
	  BEFORE UPDATE ON liquidacion FOR EACH ROW
	BEGIN
	  CALL validar_liquidacion(NEW.inversion_num, NEW.liquidacion_fecha,
	                           NEW.liquidacion_capital);
	END //

	DELIMITER ;
\end{verbatim}

\begin{verbatim}
	SELECT ROUTINE_NAME, ROUTINE_TYPE, DEFINER, SECURITY_TYPE
	FROM information_schema.ROUTINES
	WHERE ROUTINE_SCHEMA = 'inversiones';
	
	SELECT TRIGGER_NAME, EVENT_MANIPULATION, EVENT_OBJECT_TABLE, DEFINER
	FROM information_schema.TRIGGERS
	WHERE TRIGGER_SCHEMA = 'inversiones';
\end{verbatim}

\captura{mysql-objetos.png}{Procedimientos y triggers en MySQL}{mysql-objetos}

\subsection{Permisos}

Los privilegios se otorgan a los roles, con la misma matriz que en PostgreSQL. Hay dos diferencias que simplifican la implementación:

\begin{itemize}[leftmargin=*]
	\item MySQL no concede privilegios por defecto a todos los usuarios, por lo que no es necesario retirar privilegios antes de otorgar los de cada rol. Un usuario sin privilegios sobre \nem{inversiones} puede iniciar sesión en el servidor, pero no puede ver ni usar la base de datos.
	\item No existen secuencias: \nem{AUTO\us INCREMENT} genera las claves sin requerir un privilegio adicional.
\end{itemize}

Los privilegios del auditor se otorgan sobre la base de datos completa (\nem{inversiones.*}) y no tabla por tabla, con lo que se aplican también a las tablas que se creen después. Este es el equivalente de los privilegios por defecto de PostgreSQL. El usuario de respaldo ya cuenta con sus privilegios sobre la base de datos (sección~\ref{sec:roles-mysql}).

\begin{verbatim}
	-- Gerente
	GRANT SELECT, UPDATE         ON inversiones.empresa     TO rol_gerente;
	GRANT SELECT, INSERT, UPDATE ON inversiones.sucursal    TO rol_gerente;
	GRANT SELECT, UPDATE         ON inversiones.cliente     TO rol_gerente;
	GRANT SELECT, INSERT, UPDATE ON inversiones.aval        TO rol_gerente;
	GRANT SELECT, UPDATE, DELETE ON inversiones.inversion   TO rol_gerente;
	GRANT SELECT                 ON inversiones.interes     TO rol_gerente;
	GRANT SELECT                 ON inversiones.pago        TO rol_gerente;
	GRANT SELECT                 ON inversiones.liquidacion TO rol_gerente;

	-- Ejecutivo
	GRANT SELECT         ON inversiones.empresa     TO rol_ejecutivo;
	GRANT SELECT         ON inversiones.sucursal    TO rol_ejecutivo;
	GRANT SELECT, INSERT ON inversiones.cliente     TO rol_ejecutivo;
	GRANT UPDATE (cliente_telefono, cliente_direccion)
	  ON inversiones.cliente TO rol_ejecutivo;
	GRANT SELECT         ON inversiones.aval        TO rol_ejecutivo;
	GRANT SELECT, INSERT ON inversiones.inversion   TO rol_ejecutivo;
	GRANT SELECT         ON inversiones.interes     TO rol_ejecutivo;
	GRANT SELECT         ON inversiones.pago        TO rol_ejecutivo;
	GRANT SELECT         ON inversiones.liquidacion TO rol_ejecutivo;
	GRANT EXECUTE ON PROCEDURE inversiones.generar_tabla_interes
	  TO rol_ejecutivo;

	-- Tesorería
	GRANT SELECT         ON inversiones.empresa     TO rol_tesoreria;
	GRANT SELECT         ON inversiones.sucursal    TO rol_tesoreria;
	GRANT SELECT         ON inversiones.cliente     TO rol_tesoreria;
	GRANT SELECT         ON inversiones.aval        TO rol_tesoreria;
	GRANT SELECT         ON inversiones.inversion   TO rol_tesoreria;
	GRANT SELECT         ON inversiones.interes     TO rol_tesoreria;
	GRANT SELECT, INSERT ON inversiones.pago        TO rol_tesoreria;
	GRANT SELECT, INSERT ON inversiones.liquidacion TO rol_tesoreria;

	-- Auditor
	GRANT SELECT ON inversiones.* TO rol_auditor;
\end{verbatim}

Los privilegios de cada rol se consultan con \nem{SHOW GRANTS} (figura~\ref{fig:mysql-permisos}):

\begin{verbatim}
	SHOW GRANTS FOR rol_gerente;
	SHOW GRANTS FOR rol_ejecutivo;
	SHOW GRANTS FOR rol_tesoreria;
	SHOW GRANTS FOR rol_auditor;
	SHOW GRANTS FOR rol_respaldo;
\end{verbatim}

\captura{mysql-permisos.png}{Privilegios de los roles en MySQL}{mysql-permisos}

\section{Datos de prueba}

Los datos de prueba son ficticios y representan la operación de la empresa entre enero de 2024 y la fecha de carga, el 28 de octubre de 2026. El volumen supera el mínimo de quince registros por tabla por dos razones. Primero, cada estado de una inversión y cada forma de pago debe tener varios registros, para que las consultas de los objetivos devuelvan resultados significativos. Segundo, las tablas que se particionan en la segunda instancia deben tener registros en cada rango. Aun así, el volumen se mantiene pequeño para que la carga ocupe una fracción mínima de la memoria asignada a cada instancia.

La tabla~\ref{tab:escenarios-prueba} describe los escenarios que cubren los datos y la tabla~\ref{tab:volumen-prueba} el número de registros por tabla.

\begin{table}[H]
	\centering
	\small
	\caption{Escenarios cubiertos por los datos de prueba}
	\label{tab:escenarios-prueba}
	\begin{tabular}{@{}L{4.2cm} L{3.4cm} L{8.4cm}@{}}
		\toprule
		\textbf{Escenario} & \textbf{Registros} & \textbf{Propósito} \\
		\midrule
		Inversiones liquidadas & Inversiones 1 a 15 & Capital devuelto; las
		inversiones 5, 10 y 15 se liquidan 20 días tarde y pagan interés
		moratorio. \\
		Inversiones vencidas sin liquidar & Inversiones 16 a 25 & Los dos
		últimos periodos de cada una quedan sin pago (20 periodos), lo que
		genera adeudos de interés y de capital. \\
		Inversiones vigentes & Inversiones 26 a 40 & Periodos ya pagados y 144
		periodos que aún no vencen. \\
		Inversiones sin aval & 13 inversiones & Llave foránea
		\nem{aval\us num} nula. \\
		Avales sin inversiones & 5 avales & Mínimo cero de la relación
		\nem{aval} -- \nem{inversion}. \\
		Clientes con varias inversiones & 15 clientes & Consultas agrupadas
		por cliente. \\
		Pagos tardíos & 37 pagos de 10 inversiones & Pagos realizados 10 días
		después de la fecha de pago, con interés moratorio. \\
		Pagos parciales & 3 periodos (inversiones 14, 21 y 35) & Un periodo
		cubierto con dos pagos: 60\,\% en la fecha de pago y el resto 15 días
		después, con interés moratorio. \\
		Plazos y montos & 40 inversiones & Plazos de 6, 12, 15, 18 y 24
		meses; capitales de \$50{,}000 a \$500{,}000; tasas de 0.85\,\% a
		1.25\,\% mensual y moratorias de 2.5\,\% y 3.0\,\%. \\
		Distribución en el tiempo & 2024 a 2026 & Suscripciones: 18 en 2024,
		17 en 2025 y 5 en 2026. Pagos: 74 en 2024, 237 en 2025 y 167 en
		2026. \\
		\bottomrule
	\end{tabular}
\end{table}

\begin{table}[H]
	\centering
	\small
	\caption{Volumen de los datos de prueba}
	\label{tab:volumen-prueba}
	\begin{tabular}{@{}l r L{9.5cm}@{}}
		\toprule
		\textbf{Tabla} & \textbf{Registros} & \textbf{Origen} \\
		\midrule
		\nem{empresa} & 1 & Inserción directa \\
		\nem{sucursal} & 15 & Inserción directa \\
		\nem{cliente} & 25 & Inserción directa \\
		\nem{aval} & 15 & Inserción directa \\
		\nem{inversion} & 40 & Inserción directa \\
		\nem{interes} & 639 & Procedimiento \nem{generar\us tabla\us interes},
		un periodo por cada mes de plazo \\
		\nem{liquidacion} & 15 & Derivado de \nem{inversion} \\
		\nem{pago} & 478 & Derivado de \nem{interes}: 435 completos, 37
		tardíos y 6 de periodos parciales \\
		\nem{bitacora} & 0 & La carga inicial no se registra \\
		\bottomrule
	\end{tabular}
\end{table}

\paragraph{Forma de la carga.} El DBA realiza la carga sobre las tablas vacías y antes de crear los triggers de bitácora del capítulo~\ref{cap:bitacora}. La carga inicial no es una operación de los usuarios del sistema, por lo que no se audita. Como las tablas están vacías, las llaves sustitutas se generan del 1 al $N$ en el orden de inserción, y las inversiones referencian a sus sucursales, clientes y avales por esos números. Las inserciones de las tablas independientes y de las inversiones son idénticas en ambos gestores. La moneda toma el valor por defecto (\nem{MXN}).

Los periodos de interés no se insertan como valores literales: se generan con el procedimiento \nem{generar\us tabla\us interes}, el mismo que usa el ejecutivo en la operación, de modo que la carga sigue las reglas del sistema. Los pagos y las liquidaciones tampoco se capturan uno por uno. Se derivan de los periodos y de las inversiones mediante \nem{INSERT \ldots{} SELECT} con reglas fijas, lo que mantiene los montos consistentes y la carga reproducible. Esta es la única parte que difiere entre gestores, por la aritmética de fechas. Las liquidaciones se insertan antes que los pagos, porque la regla de los periodos sin pago identifica a las inversiones vencidas como aquellas que no tienen liquidación.

\subsection{Datos comunes}

\begin{verbatim}
	-- Ejecutado por dlandero en la base de datos inversiones
	INSERT INTO empresa (empresa_rfc, empresa_nombre, empresa_direccion)
	VALUES ('IGO150312AB1', 'Inversiones del Golfo, S.A. de C.V.',
	        'Av. Independencia 1250, Centro, Veracruz, Ver.');

	INSERT INTO sucursal (empresa_rfc, sucursal_nombre, sucursal_ciudad,
	                      sucursal_direccion)
	VALUES
	  ('IGO150312AB1', 'Sucursal Centro', 'Veracruz',
	   'Av. Independencia 1250, Veracruz'),
	  ('IGO150312AB1', 'Sucursal Costa Verde', 'Boca del Río',
	   'Blvd. Manuel Ávila Camacho 3010, Boca del Río'),
	  ('IGO150312AB1', 'Sucursal Xalapa Centro', 'Xalapa',
	   'Av. Xalapa 455, Xalapa'),
	  ('IGO150312AB1', 'Sucursal Córdoba', 'Córdoba',
	   'Calle 3 número 210, Córdoba'),
	  ('IGO150312AB1', 'Sucursal Orizaba', 'Orizaba',
	   'Av. Oriente 6 número 880, Orizaba'),
	  ('IGO150312AB1', 'Sucursal Coatzacoalcos', 'Coatzacoalcos',
	   'Av. Zaragoza 512, Coatzacoalcos'),
	  ('IGO150312AB1', 'Sucursal Poza Rica', 'Poza Rica',
	   'Blvd. Adolfo Ruiz Cortines 1401, Poza Rica'),
	  ('IGO150312AB1', 'Sucursal Tuxpan', 'Tuxpan',
	   'Av. Juárez 37, Tuxpan'),
	  ('IGO150312AB1', 'Sucursal Puebla Angelópolis', 'Puebla',
	   'Blvd. Atlixcáyotl 5320, Puebla'),
	  ('IGO150312AB1', 'Sucursal Ciudad de México Reforma', 'Ciudad de México',
	   'Paseo de la Reforma 222, Ciudad de México'),
	  ('IGO150312AB1', 'Sucursal Monterrey San Pedro', 'Monterrey',
	   'Av. Vasconcelos 150, Monterrey'),
	  ('IGO150312AB1', 'Sucursal Guadalajara Providencia', 'Guadalajara',
	   'Av. Providencia 2577, Guadalajara'),
	  ('IGO150312AB1', 'Sucursal Mérida Norte', 'Mérida',
	   'Calle 60 número 299, Mérida'),
	  ('IGO150312AB1', 'Sucursal Oaxaca Centro', 'Oaxaca',
	   'Av. Juárez 703, Oaxaca'),
	  ('IGO150312AB1', 'Sucursal Villahermosa Tabasco 2000', 'Villahermosa',
	   'Prol. Paseo Usumacinta 1001, Villahermosa');

	INSERT INTO cliente (cliente_rfc, cliente_nombre, cliente_apellido,
	                     cliente_direccion, cliente_telefono)
	VALUES
	  ('HEGJ720604HM3', 'José Luis', 'Hernández García',
	   'Calle Hidalgo 100, Veracruz', '2291000000'),
	  ('LOMM791107RZ6', 'María Fernanda', 'López Martínez',
	   'Calle Morelos 113, Boca del Río', '2291037291'),
	  ('RASC860410YL9', 'Carlos', 'Ramírez Sánchez',
	   'Calle Allende 126, Xalapa', '2291074582'),
	  ('TOGA680913FY2', 'Ana Sofía', 'Torres Gómez',
	   'Calle Juárez 139, Córdoba', '2291111873'),
	  ('FODJ750216NK5', 'Jorge', 'Flores Díaz',
	   'Calle Madero 152, Orizaba', '2291149164'),
	  ('MOCL820719WX8', 'Lucía', 'Morales Cruz',
	   'Calle Hidalgo 165, Coatzacoalcos', '2291186455'),
	  ('REOM891222DJ1', 'Miguel Ángel', 'Reyes Ortiz',
	   'Calle Morelos 178, Poza Rica', '2291223746'),
	  ('JIRD710525LW4', 'Daniela', 'Jiménez Ruiz',
	   'Calle Allende 191, Tuxpan', '2291261037'),
	  ('VACR781001UH7', 'Roberto', 'Vázquez Castillo',
	   'Calle Juárez 204, Puebla', '2291298328'),
	  ('MEAP850304BV0', 'Patricia', 'Mendoza Aguilar',
	   'Calle Madero 217, Ciudad de México', '2291335619'),
	  ('GURF670807JG3', 'Fernando', 'Gutiérrez Romero',
	   'Calle Hidalgo 230, Monterrey', '2291372910'),
	  ('CAHG740110SU6', 'Gabriela', 'Castro Herrera',
	   'Calle Morelos 243, Guadalajara', '2291410201'),
	  ('CAMR810613ZF9', 'Ricardo', 'Chávez Medina',
	   'Calle Allende 256, Mérida', '2291447492'),
	  ('VARA881116GT2', 'Alejandra', 'Vargas Ríos',
	   'Calle Juárez 269, Oaxaca', '2291484783'),
	  ('DOPE700419PE5', 'Eduardo', 'Domínguez Peña',
	   'Calle Madero 282, Villahermosa', '2291522074'),
	  ('SAIV770922XS8', 'Verónica', 'Salazar Ibarra',
	   'Calle Hidalgo 295, Veracruz', '2291559365'),
	  ('NACA840225ED1', 'Arturo', 'Navarro Cortés',
	   'Calle Morelos 308, Boca del Río', '2291596656'),
	  ('ROFC660701MR4', 'Claudia', 'Rojas Fuentes',
	   'Calle Allende 321, Xalapa', '2291633947'),
	  ('GULH731204VC7', 'Héctor', 'Guzmán Luna',
	   'Calle Juárez 334, Córdoba', '2291671238'),
	  ('EIVM800507CP0', 'Mónica', 'Espinoza Vega',
	   'Calle Madero 347, Orizaba', '2291708529'),
	  ('SAPF871010KB3', 'Francisco', 'Sandoval Pacheco',
	   'Calle Hidalgo 360, Coatzacoalcos', '2291745820'),
	  ('CAST690313TN6', 'Teresa', 'Cabrera Soto',
	   'Calle Morelos 373, Poza Rica', '2291783111'),
	  ('COLR760816AA9', 'Raúl', 'Contreras Lara',
	   'Calle Allende 386, Tuxpan', '2291820402'),
	  ('MORS830119HM2', 'Silvia', 'Montes Rosales',
	   'Calle Juárez 399, Puebla', '2291857693'),
	  ('NUVA650622RZ5', 'Andrés', 'Núñez Valdez',
	   'Calle Madero 412, Ciudad de México', '2291894984');

	INSERT INTO aval (aval_rfc, aval_nombre, aval_apellido, aval_direccion)
	VALUES
	  ('PEZM630913ED0', 'Manuel', 'Pérez Zamora',
	   'Av. Revolución 200, Córdoba'),
	  ('DECL700216MR3', 'Laura', 'Delgado Campos',
	   'Av. Insurgentes 217, Orizaba'),
	  ('AOBE770719VC6', 'Ernesto', 'Acosta Bravo',
	   'Av. Constitución 234, Coatzacoalcos'),
	  ('POGR591222CP9', 'Rosa María', 'Ponce Galván',
	   'Av. Reforma 251, Poza Rica'),
	  ('OOTG660525KB2', 'Guillermo', 'Ochoa Trejo',
	   'Av. Universidad 268, Tuxpan'),
	  ('SOMB731001TN5', 'Beatriz', 'Solís Mejía',
	   'Av. Revolución 285, Puebla'),
	  ('CAQA800304AA8', 'Alberto', 'Carrillo Quiroz',
	   'Av. Insurgentes 302, Ciudad de México'),
	  ('BALL620807HM1', 'Leticia', 'Bautista Lozano',
	   'Av. Constitución 319, Monterrey'),
	  ('MOAO690110RZ4', 'Óscar', 'Molina Arias',
	   'Av. Reforma 336, Guadalajara'),
	  ('FITE760613YL7', 'Elena', 'Figueroa Tapia',
	   'Av. Universidad 353, Mérida'),
	  ('VARS581116FY0', 'Sergio', 'Valencia Robles',
	   'Av. Revolución 370, Oaxaca'),
	  ('EUDA650419NK3', 'Adriana', 'Esquivel Duarte',
	   'Av. Insurgentes 387, Villahermosa'),
	  ('CEMJ720922WX6', 'Javier', 'Cervantes Mora',
	   'Av. Constitución 404, Veracruz'),
	  ('BECN790225DJ9', 'Norma', 'Beltrán Cuevas',
	   'Av. Reforma 421, Boca del Río'),
	  ('AESR610701LW2', 'Rubén', 'Avendaño Silva',
	   'Av. Universidad 438, Xalapa');

	INSERT INTO inversion (pagare_num, sucursal_num, cliente_num, aval_num,
	                       inversion_capital, inversion_tasa,
	                       inversion_tasa_moratoria, inversion_dia_pago,
	                       pagare_fecha, pagare_jurisdiccion,
	                       inversion_fecha_vencimiento)
	VALUES
	  ('0001/2024',  1,  1,    1,  50000.00, 0.8500, 3.0000, 10,
	   '2024-01-01', 'Veracruz', '2024-07-10'),
	  ('0002/2024',  2,  2,    2,  75000.00, 0.9500, 2.5000, 19,
	   '2024-02-06', 'Boca del Río', '2025-02-19'),
	  ('0003/2024',  3,  3, NULL, 100000.00, 1.0500, 3.0000, 28,
	   '2024-03-11', 'Xalapa', '2025-09-28'),
	  ('0004/2024',  4,  4,    4, 150000.00, 1.1500, 2.5000,  9,
	   '2024-04-16', 'Córdoba', '2024-10-09'),
	  ('0005/2024',  5,  5,    5, 200000.00, 1.2500, 3.0000, 18,
	   '2024-05-21', 'Orizaba', '2025-05-18'),
	  ('0006/2024',  6,  6, NULL, 250000.00, 0.8500, 2.5000, 27,
	   '2024-06-26', 'Coatzacoalcos', '2025-12-27'),
	  ('0008/2024',  7,  7,    7, 300000.00, 0.9500, 3.0000,  8,
	   '2024-07-04', 'Poza Rica', '2025-01-08'),
	  ('0009/2024',  8,  8,    8, 500000.00, 1.0500, 2.5000, 17,
	   '2024-08-09', 'Tuxpan', '2025-08-17'),
	  ('0011/2024',  9,  9, NULL,  50000.00, 1.1500, 3.0000, 26,
	   '2024-09-14', 'Puebla', '2026-03-26'),
	  ('0013/2024', 10, 10,   10,  75000.00, 1.2500, 2.5000,  7,
	   '2024-10-19', 'Ciudad de México', '2025-04-07'),
	  ('0016/2024', 11, 11,   11, 100000.00, 0.8500, 3.0000, 16,
	   '2024-11-24', 'Monterrey', '2025-11-16'),
	  ('0017/2024', 12, 12, NULL, 150000.00, 0.9500, 2.5000, 25,
	   '2024-12-02', 'Guadalajara', '2026-06-25'),
	  ('0001/2025', 13, 13,   13, 200000.00, 1.0500, 3.0000,  6,
	   '2025-01-07', 'Mérida', '2025-07-06'),
	  ('0003/2025', 14, 14,   14, 250000.00, 1.1500, 2.5000, 15,
	   '2025-02-12', 'Oaxaca', '2026-02-15'),
	  ('0007/2025', 15, 15, NULL, 300000.00, 1.2500, 3.0000, 24,
	   '2025-03-17', 'Villahermosa', '2026-09-24'),
	  ('0007/2024',  1, 16,    1, 500000.00, 0.8500, 2.5000,  5,
	   '2024-07-02', 'Veracruz', '2025-07-05'),
	  ('0010/2024',  2, 17,    2,  50000.00, 0.9500, 3.0000, 14,
	   '2024-08-09', 'Boca del Río', '2025-11-14'),
	  ('0012/2024',  3, 18, NULL,  75000.00, 1.0500, 2.5000, 23,
	   '2024-09-16', 'Xalapa', '2025-09-23'),
	  ('0014/2024',  4, 19,    4, 100000.00, 1.1500, 3.0000,  4,
	   '2024-10-23', 'Córdoba', '2026-01-04'),
	  ('0015/2024',  5, 20,    5, 150000.00, 1.2500, 2.5000, 13,
	   '2024-11-03', 'Orizaba', '2025-11-13'),
	  ('0018/2024',  6, 21, NULL, 200000.00, 0.8500, 3.0000, 22,
	   '2024-12-10', 'Coatzacoalcos', '2026-03-22'),
	  ('0002/2025',  7, 22,    7, 250000.00, 0.9500, 2.5000,  3,
	   '2025-01-17', 'Poza Rica', '2026-01-03'),
	  ('0004/2025',  8, 23,    8, 300000.00, 1.0500, 3.0000, 12,
	   '2025-02-24', 'Tuxpan', '2026-05-12'),
	  ('0006/2025',  9, 24, NULL, 500000.00, 1.1500, 2.5000, 21,
	   '2025-03-04', 'Puebla', '2026-03-21'),
	  ('0009/2025', 10, 25,   10,  50000.00, 1.2500, 3.0000,  2,
	   '2025-04-11', 'Ciudad de México', '2026-07-02'),
	  ('0005/2025', 11,  1,   11,  75000.00, 0.8500, 2.5000, 11,
	   '2025-03-01', 'Monterrey', '2027-03-11'),
	  ('0008/2025', 12,  2, NULL, 100000.00, 0.9500, 3.0000, 20,
	   '2025-04-05', 'Guadalajara', '2027-04-20'),
	  ('0010/2025', 13,  3,   13, 150000.00, 1.0500, 2.5000,  1,
	   '2025-05-09', 'Mérida', '2027-05-01'),
	  ('0011/2025', 14,  4,   14, 200000.00, 1.1500, 3.0000, 10,
	   '2025-06-13', 'Oaxaca', '2027-06-10'),
	  ('0012/2025', 15,  5, NULL, 250000.00, 1.2500, 2.5000, 19,
	   '2025-07-17', 'Villahermosa', '2027-01-19'),
	  ('0013/2025',  1,  6,    1, 300000.00, 0.8500, 3.0000, 28,
	   '2025-08-21', 'Veracruz', '2027-08-28'),
	  ('0014/2025',  2,  7,    2, 500000.00, 0.9500, 2.5000,  9,
	   '2025-09-25', 'Boca del Río', '2027-03-09'),
	  ('0015/2025',  3,  8, NULL,  50000.00, 1.0500, 3.0000, 18,
	   '2025-10-02', 'Xalapa', '2027-10-18'),
	  ('0016/2025',  4,  9,    4,  75000.00, 1.1500, 2.5000, 27,
	   '2025-11-06', 'Córdoba', '2027-05-27'),
	  ('0017/2025',  5, 10,    5, 100000.00, 1.2500, 3.0000,  8,
	   '2025-12-10', 'Orizaba', '2027-12-08'),
	  ('0001/2026',  6, 11, NULL, 150000.00, 0.8500, 2.5000, 17,
	   '2026-01-14', 'Coatzacoalcos', '2027-07-17'),
	  ('0002/2026',  7, 12,    7, 200000.00, 0.9500, 3.0000, 26,
	   '2026-02-18', 'Poza Rica', '2028-02-26'),
	  ('0003/2026',  8, 13,    8, 250000.00, 1.0500, 2.5000,  7,
	   '2026-03-22', 'Tuxpan', '2027-09-07'),
	  ('0004/2026',  9, 14, NULL, 300000.00, 1.1500, 3.0000, 16,
	   '2026-04-26', 'Puebla', '2028-04-16'),
	  ('0005/2026', 10, 15,   10, 500000.00, 1.2500, 2.5000, 25,
	   '2026-05-03', 'Ciudad de México', '2027-11-25');

	-- Tabla de interés de cada inversión
	CALL generar_tabla_interes(1);   CALL generar_tabla_interes(2);
	CALL generar_tabla_interes(3);   CALL generar_tabla_interes(4);
	CALL generar_tabla_interes(5);   CALL generar_tabla_interes(6);
	CALL generar_tabla_interes(7);   CALL generar_tabla_interes(8);
	CALL generar_tabla_interes(9);   CALL generar_tabla_interes(10);
	CALL generar_tabla_interes(11);  CALL generar_tabla_interes(12);
	CALL generar_tabla_interes(13);  CALL generar_tabla_interes(14);
	CALL generar_tabla_interes(15);  CALL generar_tabla_interes(16);
	CALL generar_tabla_interes(17);  CALL generar_tabla_interes(18);
	CALL generar_tabla_interes(19);  CALL generar_tabla_interes(20);
	CALL generar_tabla_interes(21);  CALL generar_tabla_interes(22);
	CALL generar_tabla_interes(23);  CALL generar_tabla_interes(24);
	CALL generar_tabla_interes(25);  CALL generar_tabla_interes(26);
	CALL generar_tabla_interes(27);  CALL generar_tabla_interes(28);
	CALL generar_tabla_interes(29);  CALL generar_tabla_interes(30);
	CALL generar_tabla_interes(31);  CALL generar_tabla_interes(32);
	CALL generar_tabla_interes(33);  CALL generar_tabla_interes(34);
	CALL generar_tabla_interes(35);  CALL generar_tabla_interes(36);
	CALL generar_tabla_interes(37);  CALL generar_tabla_interes(38);
	CALL generar_tabla_interes(39);  CALL generar_tabla_interes(40);
\end{verbatim}

\subsection{Liquidaciones y pagos}

Las reglas para derivar las liquidaciones y los pagos son las mismas en ambos gestores:

\begin{itemize}
	\item Las inversiones 1 a 15 se liquidan en su fecha de vencimiento. Las que tienen un número múltiplo de 5 se liquidan 20 días después y pagan interés moratorio sobre el capital por esos 20 días.
	\item Solo se pagan los periodos cuya fecha final es anterior o igual a la fecha de carga. En las inversiones vencidas sin liquidar, los periodos	que terminan en sus dos últimos meses quedan sin pago.
	\item En las inversiones con número múltiplo de 4, los periodos que	terminan en marzo, junio, septiembre o diciembre se pagan 10 días tarde,	con interés moratorio sobre el interés del periodo.
	\item En las inversiones con número múltiplo de 7, el periodo que termina en junio, si no es tardío, se paga en dos partes: 60\,\% en la fecha de pago y el resto 15 días después, con interés moratorio sobre el resto.
	\item Los demás periodos se pagan completos en su fecha final.
\end{itemize}

El interés moratorio se calcula de forma proporcional a los días de atraso, sobre la tasa moratoria mensual de la inversión y meses de 30 días.

\subsubsection{PostgreSQL}

\begin{verbatim}
	INSERT INTO liquidacion (inversion_num, liquidacion_fecha,
	                         liquidacion_capital, liquidacion_mora)
	SELECT inversion_num,
	       CASE WHEN inversion_num % 5 = 0
	            THEN inversion_fecha_vencimiento + 20
	            ELSE inversion_fecha_vencimiento
	       END,
	       inversion_capital,
	       CASE WHEN inversion_num % 5 = 0
	            THEN ROUND(inversion_capital * inversion_tasa_moratoria
	                       / 100 * 20 / 30, 2)
	            ELSE 0
	       END
	FROM inversion
	WHERE inversion_num BETWEEN 1 AND 15;

	INSERT INTO pago (inversion_num, interes_fecha_inicial, pago_fecha,
	                  pago_interes, pago_mora)
	WITH periodo AS (
	  SELECT i.inversion_num, i.interes_fecha_inicial,
	         i.interes_fecha_final AS fin, i.interes_monto AS monto,
	         v.inversion_tasa_moratoria AS tasa_mor,
	         (v.inversion_num % 4 = 0 AND
	          EXTRACT(MONTH FROM i.interes_fecha_final) IN (3, 6, 9, 12))
	           AS tardio,
	         (v.inversion_num % 7 = 0 AND
	          EXTRACT(MONTH FROM i.interes_fecha_final) = 6) AS parcial
	  FROM interes i
	    JOIN inversion v ON v.inversion_num = i.inversion_num
	  WHERE i.interes_fecha_final <= DATE '2026-10-28'
	    AND NOT (v.inversion_fecha_vencimiento < DATE '2026-10-28'
	             AND NOT EXISTS (SELECT 1 FROM liquidacion l
	                             WHERE l.inversion_num = v.inversion_num)
	             AND i.interes_fecha_final >
	                 v.inversion_fecha_vencimiento - INTERVAL '2 months')
	)
	-- Pago completo en la fecha de pago
	SELECT inversion_num, interes_fecha_inicial, fin, monto, 0
	FROM periodo
	WHERE NOT tardio AND NOT parcial
	UNION ALL
	-- Pago tardío: 10 días de atraso
	SELECT inversion_num, interes_fecha_inicial, fin + 10, monto,
	       ROUND(monto * tasa_mor / 100 * 10 / 30, 2)
	FROM periodo
	WHERE tardio
	UNION ALL
	-- Pago parcial: 60 % en la fecha de pago
	SELECT inversion_num, interes_fecha_inicial, fin,
	       ROUND(monto * 0.6, 2), 0
	FROM periodo
	WHERE parcial AND NOT tardio
	UNION ALL
	-- Resto del pago parcial: 15 días de atraso
	SELECT inversion_num, interes_fecha_inicial, fin + 15,
	       monto - ROUND(monto * 0.6, 2),
	       ROUND((monto - ROUND(monto * 0.6, 2)) * tasa_mor / 100 * 15 / 30, 2)
	FROM periodo
	WHERE parcial AND NOT tardio;
\end{verbatim}

\subsubsection{MySQL}

\begin{verbatim}
	INSERT INTO liquidacion (inversion_num, liquidacion_fecha,
	                         liquidacion_capital, liquidacion_mora)
	SELECT inversion_num,
	       CASE WHEN MOD(inversion_num, 5) = 0
	            THEN DATE_ADD(inversion_fecha_vencimiento, INTERVAL 20 DAY)
	            ELSE inversion_fecha_vencimiento
	       END,
	       inversion_capital,
	       CASE WHEN MOD(inversion_num, 5) = 0
	            THEN ROUND(inversion_capital * inversion_tasa_moratoria
	                       / 100 * 20 / 30, 2)
	            ELSE 0
	       END
	FROM inversion
	WHERE inversion_num BETWEEN 1 AND 15;

	INSERT INTO pago (inversion_num, interes_fecha_inicial, pago_fecha,
	                  pago_interes, pago_mora)
	WITH periodo AS (
	  SELECT i.inversion_num, i.interes_fecha_inicial,
	         i.interes_fecha_final AS fin, i.interes_monto AS monto,
	         v.inversion_tasa_moratoria AS tasa_mor,
	         (MOD(v.inversion_num, 4) = 0 AND
	          MONTH(i.interes_fecha_final) IN (3, 6, 9, 12)) AS tardio,
	         (MOD(v.inversion_num, 7) = 0 AND
	          MONTH(i.interes_fecha_final) = 6) AS parcial
	  FROM interes i
	    JOIN inversion v ON v.inversion_num = i.inversion_num
	  WHERE i.interes_fecha_final <= DATE '2026-10-28'
	    AND NOT (v.inversion_fecha_vencimiento < DATE '2026-10-28'
	             AND NOT EXISTS (SELECT 1 FROM liquidacion l
	                             WHERE l.inversion_num = v.inversion_num)
	             AND i.interes_fecha_final >
	                 DATE_SUB(v.inversion_fecha_vencimiento, INTERVAL 2 MONTH))
	)
	-- Pago completo en la fecha de pago
	SELECT inversion_num, interes_fecha_inicial, fin, monto, 0
	FROM periodo
	WHERE NOT tardio AND NOT parcial
	UNION ALL
	-- Pago tardío: 10 días de atraso
	SELECT inversion_num, interes_fecha_inicial,
	       DATE_ADD(fin, INTERVAL 10 DAY), monto,
	       ROUND(monto * tasa_mor / 100 * 10 / 30, 2)
	FROM periodo
	WHERE tardio
	UNION ALL
	-- Pago parcial: 60 % en la fecha de pago
	SELECT inversion_num, interes_fecha_inicial, fin,
	       ROUND(monto * 0.6, 2), 0
	FROM periodo
	WHERE parcial AND NOT tardio
	UNION ALL
	-- Resto del pago parcial: 15 días de atraso
	SELECT inversion_num, interes_fecha_inicial,
	       DATE_ADD(fin, INTERVAL 15 DAY),
	       monto - ROUND(monto * 0.6, 2),
	       ROUND((monto - ROUND(monto * 0.6, 2)) * tasa_mor / 100 * 15 / 30, 2)
	FROM periodo
	WHERE parcial AND NOT tardio;
\end{verbatim}

\subsection{Resultado de la carga}

La consulta siguiente, igual en ambos gestores, cuenta los registros de cada tabla para compararlos con la tabla~\ref{tab:volumen-prueba}.

\begin{verbatim}
	SELECT 'empresa' AS tabla, COUNT(*) AS registros FROM empresa
	UNION ALL SELECT 'sucursal',    COUNT(*) FROM sucursal
	UNION ALL SELECT 'cliente',     COUNT(*) FROM cliente
	UNION ALL SELECT 'aval',        COUNT(*) FROM aval
	UNION ALL SELECT 'inversion',   COUNT(*) FROM inversion
	UNION ALL SELECT 'interes',     COUNT(*) FROM interes
	UNION ALL SELECT 'liquidacion', COUNT(*) FROM liquidacion
	UNION ALL SELECT 'pago',        COUNT(*) FROM pago
	UNION ALL SELECT 'bitacora',    COUNT(*) FROM bitacora;
\end{verbatim}

\captura{pg-datos-conteo.png}{Registros cargados en PostgreSQL}{pg-datos-conteo}
\captura{mysql-datos-conteo.png}{Registros cargados en MySQL}{mysql-datos-conteo}

\section{Pruebas de alcance}
\label{sec:pruebas-alcance}
Las pruebas se ejecutan después de crear los triggers de bitácora del capítulo~\ref{cap:bitacora}, para que las operaciones permitidas queden auditadas.

Las pruebas de alcance comprueban que cada rol puede ejecutar las operaciones que le asigna la matriz de permisos de la sección~\ref{sec:roles-matriz} y que el gestor rechaza las demás. Cada usuario ejecuta al menos una operación permitida y una denegada. Las operaciones son las que realizaría el puesto en la operación diaria, de modo que las pruebas permitidas quedan como registros reales del sistema: un cliente, una inversión con su tabla de interés y un pago. Las sentencias son iguales en ambos gestores; solo cambian la conexión y el texto del error. La tabla~\ref{tab:pruebas-alcance} resume las pruebas.

\begin{table}[H]
	\centering
	\small
	\caption{Pruebas de alcance por rol}
	\label{tab:pruebas-alcance}
	\begin{tabular}{@{}L{3.2cm} L{8.2cm} L{4.6cm}@{}}
		\toprule
		\textbf{Usuario (rol)} & \textbf{Operación} & \textbf{Resultado	esperado} \\
		\midrule
		\nem{dmejia} \newline (\nem{rol\us gerente}) & Actualizar la dirección de una sucursal & Permitida \\
		& Registrar un pago & Denegada: sin privilegio \\
		& Eliminar una inversión con periodos de interés & Denegada: llave foránea (nota 1 de la matriz) \\
		\midrule
		\nem{pzuniga} \newline (\nem{rol\us ejecutivo}) & Registrar un cliente y una inversión, generar su tabla de interés y actualizar el teléfono del cliente & Permitida \\
		& Modificar el RFC del cliente & Denegada: privilegio por columna \\
		& Insertar un periodo de interés directamente & Denegada: sin privilegio \\
		\midrule
		\nem{xzaleta} \newline (\nem{rol\us tesoreria}) & Registrar el pago	tardío de un periodo vencido & Permitida \\
		& Modificar el capital de una inversión & Denegada: sin privilegio \\
		\midrule
		\nem{enovoa} \newline (\nem{rol\us auditor}) & Consultar los pagos con interés moratorio & Permitida \\
		& Eliminar registros de la bitácora & Denegada: sin privilegio \\
		\midrule
		\nem{respaldo} \newline (\nem{rol\us respaldo}) & Generar un respaldo lógico de la base de datos & Permitida \\
		& Eliminar pagos & Denegada: sin privilegio \\
		\bottomrule
	\end{tabular}
\end{table}

\subsection{Sentencias de prueba}

\paragraph{Gerente.} La eliminación de la inversión 40 es una operación que el gerente tiene permitida, pero el gestor la rechaza porque la inversión ya tiene periodos de interés.

\begin{verbatim}
	-- Permitida
	UPDATE sucursal
	SET sucursal_direccion = 'Av. Paseo Tabasco 1203, Villahermosa'
	WHERE sucursal_num = 15;

	-- Denegada: sin privilegio
	INSERT INTO pago (inversion_num, interes_fecha_inicial, pago_fecha,
	                  pago_interes)
	VALUES (16, '2025-05-05', '2026-10-28', 4250.00);

	-- Denegada: llave foránea
	DELETE FROM inversion WHERE inversion_num = 40;
\end{verbatim}

\paragraph{Ejecutivo.} El ejecutivo registra una nueva clienta y su inversión, que recibe el número 41, y genera su tabla de interés de 12 periodos mediante el procedimiento.

\begin{verbatim}
	-- Permitidas
	INSERT INTO cliente (cliente_rfc, cliente_nombre, cliente_apellido,
	                     cliente_direccion, cliente_telefono)
	VALUES ('OESM920314HT5', 'Mariana', 'Ortega Salinas',
	        'Calle Enríquez 58, Xalapa', '2281234567');

	INSERT INTO inversion (pagare_num, sucursal_num, cliente_num, aval_num,
	                       inversion_capital, inversion_tasa,
	                       inversion_tasa_moratoria, inversion_dia_pago,
	                       pagare_fecha, pagare_jurisdiccion,
	                       inversion_fecha_vencimiento)
	VALUES ('0006/2026', 3, 26, NULL, 120000.00, 1.0000, 2.5000, 15,
	        '2026-10-28', 'Xalapa', '2027-10-15');

	CALL generar_tabla_interes(41);

	UPDATE cliente SET cliente_telefono = '2287654321'
	WHERE cliente_num = 26;

	-- Denegada: privilegio por columna
	UPDATE cliente SET cliente_rfc = 'OESM920314HT6'
	WHERE cliente_num = 26;

	-- Denegada: sin privilegio
	INSERT INTO interes (inversion_num, interes_fecha_inicial,
	                     interes_fecha_final, interes_monto)
	VALUES (41, '2027-10-15', '2027-11-15', 1200.00);
\end{verbatim}

\paragraph{Tesorería.} La tesorería paga el periodo del 5 de mayo al 5 de junio de 2025 de la inversión 16, que estaba vencido. El interés moratorio corresponde a 510 días de atraso: $4{,}250.00 \times 2.5\mkern 3mu\% \times 510/30 = 1{,}806.25$.


\begin{verbatim}
	-- Permitida
	INSERT INTO pago (inversion_num, interes_fecha_inicial, pago_fecha,
	                  pago_interes, pago_mora)
	VALUES (16, '2025-05-05', '2026-10-28', 4250.00, 1806.25);

	-- Denegada: sin privilegio
	UPDATE inversion SET inversion_capital = 600000.00
	WHERE inversion_num = 16;
\end{verbatim}

\paragraph{Auditor.}

\begin{verbatim}
	-- Permitida
	SELECT inversion_num, COUNT(*) AS pagos, SUM(pago_mora) AS mora
	FROM pago
	WHERE pago_mora > 0
	GROUP BY inversion_num
	ORDER BY inversion_num;

	-- Denegada: sin privilegio
	DELETE FROM bitacora;
\end{verbatim}

\paragraph{Respaldo.} La operación permitida es el respaldo lógico, que se ejecuta con la herramienta de cada gestor (ver las secciones siguientes). La operación denegada es:

\begin{verbatim}
	-- Denegada: sin privilegio
	DELETE FROM pago WHERE inversion_num = 1;
\end{verbatim}

\subsection{PostgreSQL}

Cada usuario se conecta con su propia cuenta y ejecuta sus sentencias:

\begin{verbatim}
	psql -h localhost -p 5432 -U dmejia -d inversiones
	psql -h localhost -p 5432 -U pzuniga -d inversiones
	psql -h localhost -p 5432 -U xzaleta -d inversiones
	psql -h localhost -p 5432 -U enovoa -d inversiones
	psql -h localhost -p 5432 -U respaldo -d inversiones
	
	pg_dump -h localhost -p 5432 -U respaldo -d inversiones -n inversiones -F c -f inversiones.dump
\end{verbatim}

Las operaciones permitidas responden con \nem{UPDATE 1}, \nem{INSERT 0 1}, \nem{CALL} o el resultado de la consulta. Las denegadas por falta de privilegio responden con \nem{ERROR: permission denied for table}, seguido del nombre de la tabla, incluso en la actualización del RFC, porque el ejecutivo no tiene privilegio sobre esa columna. La eliminación de la inversión responde con \nem{ERROR: update or delete on table "inversion" violates foreign key constraint \ldots{} on table "interes"}.

\captura{pg-alcance-gerente.png}{Pruebas de alcance del gerente en PostgreSQL}{pg-alcance-gerente}
\captura{pg-alcance-ejecutivo.png}{Pruebas de alcance del ejecutivo en PostgreSQL}{pg-alcance-ejecutivo}
\captura{pg-alcance-tesoreria.png}{Pruebas de alcance de tesorería en PostgreSQL}{pg-alcance-tesoreria}
\captura{pg-alcance-auditor.png}{Pruebas de alcance del auditor en PostgreSQL}{pg-alcance-auditor}
\captura{pg-alcance-respaldo.png}{Pruebas de alcance del respaldo en PostgreSQL}{pg-alcance-respaldo}

\subsection{MySQL}

\begin{verbatim}
	mysql -h localhost -P 3306 -u dmejia -p inversiones
	mysql -h localhost -P 3306 -u pzuniga -p inversiones
	mysql -h localhost -P 3306 -u xzaleta -p inversiones
	mysql -h localhost -P 3306 -u enovoa -p inversiones
	mysql -h localhost -P 3306 -u respaldo -p inversiones
	
	mysqldump -h localhost -P 3306 -u respaldo -p --single-transaction --routines --triggers inversiones > inversiones.sql
\end{verbatim}

Las operaciones permitidas responden con \nem{Query OK}. Las denegadas por falta de privilegio sobre la tabla responden con el error 1142: \nem{INSERT command denied to user \ldots{} for table 'pago'}. La actualización del RFC responde con el error 1143, propio de los privilegios por columna: \nem{UPDATE command denied to user \ldots{} for column 'cliente\us rfc' in table 'cliente'}. La eliminación de la inversión responde con el error 1451: \nem{Cannot delete or update a parent row: a foreign key constraint fails}.

\captura{mysql-alcance-gerente.png}{Pruebas de alcance del gerente en MySQL}{mysql-alcance-gerente}
\captura{mysql-alcance-ejecutivo.png}{Pruebas de alcance del ejecutivo en MySQL}{mysql-alcance-ejecutivo}
\captura{mysql-alcance-tesoreria.png}{Pruebas de alcance de tesorería en MySQL}{mysql-alcance-tesoreria}
\captura{mysql-alcance-auditor.png}{Pruebas de alcance del auditor en MySQL}{mysql-alcance-auditor}
\captura{mysql-alcance-respaldo.png}{Pruebas de alcance del respaldo en MySQL}{mysql-alcance-respaldo}
	\chapter{Bitácora}
\label{cap:bitacora}
\recursoscapitulo{08-bitacora}

La bitácora registra cada cambio efectivo en los datos del sistema: qué usuario lo hizo, cuándo, desde qué equipo, en qué tabla, con qué operación y cuáles eran los valores antes y después del cambio. Cubre las ocho tablas de datos del esquema; la propia \nem{bitacora} no se audita a sí misma. Su propósito es que el auditor pueda reconstruir la historia de cualquier registro, en particular de los montos de inversiones, pagos y liquidaciones, sin depender de la palabra de quien hizo el cambio.

\section{Diseño de la bitácora}

\paragraph{Registro por renglón.} Los triggers son \nem{FOR EACH ROW} y se ejecutan después de la operación (\nem{AFTER}). Así, una sentencia que afecta varios renglones deja un registro por cada uno, con sus valores propios. Por ejemplo, la generación de una tabla de interés de 12 periodos deja 12 registros. Ejecutarse después de la operación garantiza que solo se registren los cambios que pasaron todas las restricciones y validaciones.

\paragraph{Misma transacción.} El registro de bitácora forma parte de la transacción que hizo el cambio. Si la transacción se revierte, el registro también desaparece, de modo que la bitácora no contiene cambios que nunca ocurrieron. Los intentos rechazados quedan en el registro del servidor configurado en el capítulo~\ref{cap:entorno} (ver tabla~\ref{tab:bitacora-registro}).

\paragraph{Usuario de la sesión.} Los triggers se ejecutan con los permisos del dueño de los objetos (\nem{SECURITY DEFINER} en PostgreSQL y \nem{DEFINER} en MySQL). Por eso el usuario se obtiene de la sesión, y no del contexto de ejecución, que siempre sería el del DBA. Lo mismo ocurre con el procedimiento \nem{generar\us tabla\us interes}: los periodos que crea quedan registrados a nombre del ejecutivo que lo llamó.

\paragraph{Valores en JSON.} Los valores anteriores y nuevos se guardan como un objeto JSON con todas las columnas del renglón. Una sola tabla de bitácora sirve así para las ocho tablas, sin una columna por cada atributo auditado. En una inserción solo hay valor nuevo y en una eliminación solo valor anterior.

La tabla~\ref{tab:bitacora-campos} indica de dónde obtiene cada gestor el valor de cada campo.

\begin{table}[H]
	\centering
	\small
	\caption{Origen de los campos de la bitácora}
	\label{tab:bitacora-campos}
	\begin{tabular}{@{}L{4.2cm} L{5.6cm} L{6.2cm}@{}}
		\toprule
		\textbf{Campo} & \textbf{PostgreSQL} & \textbf{MySQL} \\
		\midrule
		\nem{bitacora\us num} & Identidad & \nem{AUTO\us INCREMENT} \\
		\nem{bitacora\us fecha} & \nem{LOCALTIMESTAMP} & \nem{NOW()} \\
		\nem{bitacora\us usuario} & \nem{session\us user} & Parte del usuario de \nem{USER()} \\
		\nem{bitacora\us tabla} & \nem{TG\us TABLE\us NAME} & Literal en cada trigger \\
		\nem{bitacora\us operacion} & \nem{TG\us OP} & Literal en cada trigger \\
		\nem{bitacora\us ip} & \nem{inet\us client\us addr()} & Parte del equipo de \nem{USER()} \\
		\nem{bitacora\us sesion} & \nem{pg\us backend\us pid()} & \nem{CONNECTION\us ID()} \\
		\nem{bitacora\us query} & \nem{current\us query()} & \nem{PROCESSLIST\us INFO} de \nem{performance\us schema.\allowbreak threads} \\
		\nem{bitacora\us dato\us anterior} & \nem{to\us jsonb(OLD)} & \nem{JSON\us OBJECT} con las columnas de \nem{OLD} \\
		\nem{bitacora\us dato\us nuevo} & \nem{to\us jsonb(NEW)} & \nem{JSON\us OBJECT} con las columnas de \nem{NEW} \\
		\bottomrule
	\end{tabular}
\end{table}

En MySQL, \nem{USER()} devuelve el usuario y el equipo con los que el cliente se autenticó, a diferencia de \nem{CURRENT\us USER()}, que dentro de un trigger devuelve el \nem{DEFINER}. Como el servidor no resuelve nombres (\nem{skip\us name\us resolve}), el equipo es la dirección IP del cliente.

\paragraph{Bitácora y registro del servidor.} La bitácora no sustituye al registro del servidor configurado en el capítulo~\ref{cap:entorno}; se complementan, como muestra la tabla~\ref{tab:bitacora-registro}.

\begin{table}[H]
	\centering
	\small
	\caption{Bitácora y registro del servidor}
	\label{tab:bitacora-registro}
	\begin{tabular}{@{}L{5.4cm} L{5.3cm} L{5.3cm}@{}}
		\toprule
		\textbf{Evento} & \textbf{Bitácora} & \textbf{Registro del
			servidor} \\
		\midrule
		Cambio de datos confirmado & Sí, con valores anteriores y nuevos & Sí,
		solo la sentencia \\
		Cambio revertido o rechazado & No & Sí \\
		Conexiones y desconexiones & No & Sí \\
		Cambios a la estructura (DDL) & No & Sí \\
		Consulta desde SQL & Sí, el auditor la consulta como una tabla & No, se revisa en los
		archivos del servidor \\
		\bottomrule
	\end{tabular}
\end{table}

\paragraph{Protección.} Ningún rol tiene privilegios de inserción, modificación o eliminación sobre \nem{bitacora}. Solo los triggers escriben en ella, con los permisos del dueño, y el auditor y la cuenta de respaldo solo pueden consultarla. Un trigger adicional rechaza cualquier \nem{UPDATE} o \nem{DELETE} sobre la bitácora, incluso del DBA, con lo que la tabla solo admite inserciones. El DBA podría desactivar ese trigger, pero hacerlo es un cambio de estructura que queda en el registro del servidor, que administra el superusuario. Las inserciones directas tampoco son posibles para los usuarios, porque ningún rol tiene \nem{INSERT} sobre la bitácora ni puede ejecutar la función o el procedimiento que la escribe.

\section{PostgreSQL}

Una sola función de trigger sirve para las ocho tablas, porque PostgreSQL expone el nombre de la tabla y la operación en las variables \nem{TG\us TABLE\us NAME} y \nem{TG\us OP}, y convierte el renglón completo a JSON con \nem{to\us jsonb}. La función es \nem{SECURITY DEFINER} para que inserte en la bitácora con los permisos de \nem{rol\us admin\us bd}. Las funciones de trigger no pueden llamarse directamente, por lo que nadie puede usarla para insertar registros falsos.

\nem{session\us user} es el usuario que inició la sesión y no cambia con \nem{SET ROLE} ni dentro de una función \nem{SECURITY DEFINER}, a diferencia de \nem{current\us user}. \nem{current\us query()} devuelve la sentencia que envió el cliente; en los periodos generados por el procedimiento, es la llamada \nem{CALL generar\us tabla\us interes(\ldots)}.

La sentencia \nem{TRUNCATE} no activa los triggers por renglón, así que la bitácora también la rechaza con un trigger por sentencia. Las demás tablas no lo necesitan: ningún rol tiene el privilegio \nem{TRUNCATE}.

\begin{verbatim}
	-- Ejecutado por dlandero en la base de datos inversiones
	SET ROLE rol_admin_bd;

	CREATE FUNCTION fn_bitacora()
	RETURNS trigger
	LANGUAGE plpgsql
	SECURITY DEFINER
	SET search_path = inversiones
	AS $$
	BEGIN
	  INSERT INTO bitacora (bitacora_fecha, bitacora_usuario, bitacora_tabla,
	                        bitacora_operacion, bitacora_ip, bitacora_sesion,
	                        bitacora_query, bitacora_dato_anterior,
	                        bitacora_dato_nuevo)
	  VALUES (LOCALTIMESTAMP, session_user, TG_TABLE_NAME, TG_OP,
	          host(inet_client_addr()), pg_backend_pid(), current_query(),
	          CASE WHEN TG_OP IN ('UPDATE', 'DELETE') THEN to_jsonb(OLD) END,
	          CASE WHEN TG_OP IN ('INSERT', 'UPDATE') THEN to_jsonb(NEW) END);
	  RETURN NULL;
	END;
	$$;

	CREATE TRIGGER trg_empresa_bitacora
	  AFTER INSERT OR UPDATE OR DELETE ON empresa
	  FOR EACH ROW EXECUTE FUNCTION fn_bitacora();

	CREATE TRIGGER trg_sucursal_bitacora
	  AFTER INSERT OR UPDATE OR DELETE ON sucursal
	  FOR EACH ROW EXECUTE FUNCTION fn_bitacora();

	CREATE TRIGGER trg_cliente_bitacora
	  AFTER INSERT OR UPDATE OR DELETE ON cliente
	  FOR EACH ROW EXECUTE FUNCTION fn_bitacora();

	CREATE TRIGGER trg_aval_bitacora
	  AFTER INSERT OR UPDATE OR DELETE ON aval
	  FOR EACH ROW EXECUTE FUNCTION fn_bitacora();

	CREATE TRIGGER trg_inversion_bitacora
	  AFTER INSERT OR UPDATE OR DELETE ON inversion
	  FOR EACH ROW EXECUTE FUNCTION fn_bitacora();

	CREATE TRIGGER trg_interes_bitacora
	  AFTER INSERT OR UPDATE OR DELETE ON interes
	  FOR EACH ROW EXECUTE FUNCTION fn_bitacora();

	CREATE TRIGGER trg_pago_bitacora
	  AFTER INSERT OR UPDATE OR DELETE ON pago
	  FOR EACH ROW EXECUTE FUNCTION fn_bitacora();

	CREATE TRIGGER trg_liquidacion_bitacora
	  AFTER INSERT OR UPDATE OR DELETE ON liquidacion
	  FOR EACH ROW EXECUTE FUNCTION fn_bitacora();

	-- Protección de la bitácora
	CREATE FUNCTION fn_bitacora_proteger()
	RETURNS trigger
	LANGUAGE plpgsql
	AS $$
	BEGIN
	  RAISE EXCEPTION 'La bitácora no admite la operación %', TG_OP;
	END;
	$$;

	CREATE TRIGGER trg_bitacora_proteger
	  BEFORE UPDATE OR DELETE ON bitacora
	  FOR EACH ROW EXECUTE FUNCTION fn_bitacora_proteger();

	CREATE TRIGGER trg_bitacora_truncate
	  BEFORE TRUNCATE ON bitacora
	  FOR EACH STATEMENT EXECUTE FUNCTION fn_bitacora_proteger();
\end{verbatim}

La consulta siguiente lista los triggers del esquema. Es igual en ambos
gestores.

\begin{verbatim}
	SELECT event_object_table, trigger_name, action_timing,
	       event_manipulation
	FROM information_schema.triggers
	WHERE trigger_schema = 'inversiones'
	ORDER BY event_object_table, trigger_name, event_manipulation;
\end{verbatim}

\captura[12cm]{pg-bitacora-triggers.png}{Triggers de bitácora en PostgreSQL}{pg-bitacora-triggers}

\section{MySQL}

MySQL requiere un trigger por cada operación y no puede convertir un renglón completo a JSON, por lo que cada trigger enumera las columnas de su tabla con \nem{JSON\us OBJECT}. Para no repetir en los 24 triggers la obtención del usuario, el equipo, la sesión y la sentencia, los triggers llaman a un procedimiento, \nem{registrar\us bitacora}, que hace la inserción. Los triggers y el procedimiento se ejecutan con los permisos de su \nem{DEFINER}, el DBA. Ningún rol tiene \nem{EXECUTE} sobre \nem{registrar\us bitacora}, de modo que los usuarios no pueden llamarlo directamente.

MySQL no tiene una función equivalente a \nem{current\us query()}. La sentencia que ejecuta cada conexión se obtiene de la columna \nem{PROCESSLIST\us INFO} de la tabla \nem{performance\us schema.\allowbreak threads}, buscando la conexión actual con \nem{CONNECTION\us ID()}. Esta tabla está disponible porque \nem{performance\us schema} se mantiene activado, que es su valor por defecto. Ese esquema está fuera de la base de datos \nem{inversiones}, así que el permiso de consulta lo otorga el superusuario:

\begin{verbatim}
	-- Ejecutado por root
	GRANT SELECT ON performance_schema.threads TO rol_admin_bd;
\end{verbatim}

Hay dos diferencias con PostgreSQL. Primero, las acciones en cascada de las llaves foráneas no activan triggers en MySQL. Si cambia el RFC de la empresa, la bitácora registra la actualización de \nem{empresa}, pero no la de cada sucursal que se actualiza en cascada. Segundo, la sentencia \nem{TRUNCATE} requiere el privilegio \nem{DROP}, que solo tiene el DBA, y queda en el registro general del servidor.

\begin{verbatim}
	-- Ejecutado por dlandero en la base de datos inversiones
	CREATE DEFINER = CURRENT_USER PROCEDURE registrar_bitacora(
	  IN p_tabla     VARCHAR(63),
	  IN p_operacion VARCHAR(6),
	  IN p_anterior  JSON,
	  IN p_nuevo     JSON)
	SQL SECURITY DEFINER
	INSERT INTO bitacora (bitacora_fecha, bitacora_usuario, bitacora_tabla,
	                      bitacora_operacion, bitacora_ip, bitacora_sesion,
	                      bitacora_query, bitacora_dato_anterior,
	                      bitacora_dato_nuevo)
	VALUES (NOW(), SUBSTRING_INDEX(USER(), '@', 1), p_tabla, p_operacion,
	        SUBSTRING_INDEX(USER(), '@', -1), CONNECTION_ID(),
	        (SELECT PROCESSLIST_INFO
	         FROM performance_schema.threads
	         WHERE PROCESSLIST_ID = CONNECTION_ID()),
	        p_anterior, p_nuevo);

	-- empresa
	CREATE TRIGGER trg_empresa_bitacora_ins AFTER INSERT ON empresa
	FOR EACH ROW CALL registrar_bitacora('empresa', 'INSERT',
	  NULL,
	  JSON_OBJECT('empresa_rfc', NEW.empresa_rfc,
	              'empresa_nombre', NEW.empresa_nombre,
	              'empresa_direccion', NEW.empresa_direccion));

	CREATE TRIGGER trg_empresa_bitacora_upd AFTER UPDATE ON empresa
	FOR EACH ROW CALL registrar_bitacora('empresa', 'UPDATE',
	  JSON_OBJECT('empresa_rfc', OLD.empresa_rfc,
	              'empresa_nombre', OLD.empresa_nombre,
	              'empresa_direccion', OLD.empresa_direccion),
	  JSON_OBJECT('empresa_rfc', NEW.empresa_rfc,
	              'empresa_nombre', NEW.empresa_nombre,
	              'empresa_direccion', NEW.empresa_direccion));

	CREATE TRIGGER trg_empresa_bitacora_del AFTER DELETE ON empresa
	FOR EACH ROW CALL registrar_bitacora('empresa', 'DELETE',
	  JSON_OBJECT('empresa_rfc', OLD.empresa_rfc,
	              'empresa_nombre', OLD.empresa_nombre,
	              'empresa_direccion', OLD.empresa_direccion),
	  NULL);

	-- sucursal
	CREATE TRIGGER trg_sucursal_bitacora_ins AFTER INSERT ON sucursal
	FOR EACH ROW CALL registrar_bitacora('sucursal', 'INSERT',
	  NULL,
	  JSON_OBJECT('sucursal_num', NEW.sucursal_num,
	              'empresa_rfc', NEW.empresa_rfc,
	              'sucursal_nombre', NEW.sucursal_nombre,
	              'sucursal_ciudad', NEW.sucursal_ciudad,
	              'sucursal_direccion', NEW.sucursal_direccion));

	CREATE TRIGGER trg_sucursal_bitacora_upd AFTER UPDATE ON sucursal
	FOR EACH ROW CALL registrar_bitacora('sucursal', 'UPDATE',
	  JSON_OBJECT('sucursal_num', OLD.sucursal_num,
	              'empresa_rfc', OLD.empresa_rfc,
	              'sucursal_nombre', OLD.sucursal_nombre,
	              'sucursal_ciudad', OLD.sucursal_ciudad,
	              'sucursal_direccion', OLD.sucursal_direccion),
	  JSON_OBJECT('sucursal_num', NEW.sucursal_num,
	              'empresa_rfc', NEW.empresa_rfc,
	              'sucursal_nombre', NEW.sucursal_nombre,
	              'sucursal_ciudad', NEW.sucursal_ciudad,
	              'sucursal_direccion', NEW.sucursal_direccion));

	CREATE TRIGGER trg_sucursal_bitacora_del AFTER DELETE ON sucursal
	FOR EACH ROW CALL registrar_bitacora('sucursal', 'DELETE',
	  JSON_OBJECT('sucursal_num', OLD.sucursal_num,
	              'empresa_rfc', OLD.empresa_rfc,
	              'sucursal_nombre', OLD.sucursal_nombre,
	              'sucursal_ciudad', OLD.sucursal_ciudad,
	              'sucursal_direccion', OLD.sucursal_direccion),
	  NULL);

	-- cliente
	CREATE TRIGGER trg_cliente_bitacora_ins AFTER INSERT ON cliente
	FOR EACH ROW CALL registrar_bitacora('cliente', 'INSERT',
	  NULL,
	  JSON_OBJECT('cliente_num', NEW.cliente_num,
	              'cliente_rfc', NEW.cliente_rfc,
	              'cliente_nombre', NEW.cliente_nombre,
	              'cliente_apellido', NEW.cliente_apellido,
	              'cliente_direccion', NEW.cliente_direccion,
	              'cliente_telefono', NEW.cliente_telefono));

	CREATE TRIGGER trg_cliente_bitacora_upd AFTER UPDATE ON cliente
	FOR EACH ROW CALL registrar_bitacora('cliente', 'UPDATE',
	  JSON_OBJECT('cliente_num', OLD.cliente_num,
	              'cliente_rfc', OLD.cliente_rfc,
	              'cliente_nombre', OLD.cliente_nombre,
	              'cliente_apellido', OLD.cliente_apellido,
	              'cliente_direccion', OLD.cliente_direccion,
	              'cliente_telefono', OLD.cliente_telefono),
	  JSON_OBJECT('cliente_num', NEW.cliente_num,
	              'cliente_rfc', NEW.cliente_rfc,
	              'cliente_nombre', NEW.cliente_nombre,
	              'cliente_apellido', NEW.cliente_apellido,
	              'cliente_direccion', NEW.cliente_direccion,
	              'cliente_telefono', NEW.cliente_telefono));

	CREATE TRIGGER trg_cliente_bitacora_del AFTER DELETE ON cliente
	FOR EACH ROW CALL registrar_bitacora('cliente', 'DELETE',
	  JSON_OBJECT('cliente_num', OLD.cliente_num,
	              'cliente_rfc', OLD.cliente_rfc,
	              'cliente_nombre', OLD.cliente_nombre,
	              'cliente_apellido', OLD.cliente_apellido,
	              'cliente_direccion', OLD.cliente_direccion,
	              'cliente_telefono', OLD.cliente_telefono),
	  NULL);

	-- aval
	CREATE TRIGGER trg_aval_bitacora_ins AFTER INSERT ON aval
	FOR EACH ROW CALL registrar_bitacora('aval', 'INSERT',
	  NULL,
	  JSON_OBJECT('aval_num', NEW.aval_num,
	              'aval_rfc', NEW.aval_rfc,
	              'aval_nombre', NEW.aval_nombre,
	              'aval_apellido', NEW.aval_apellido,
	              'aval_direccion', NEW.aval_direccion));

	CREATE TRIGGER trg_aval_bitacora_upd AFTER UPDATE ON aval
	FOR EACH ROW CALL registrar_bitacora('aval', 'UPDATE',
	  JSON_OBJECT('aval_num', OLD.aval_num,
	              'aval_rfc', OLD.aval_rfc,
	              'aval_nombre', OLD.aval_nombre,
	              'aval_apellido', OLD.aval_apellido,
	              'aval_direccion', OLD.aval_direccion),
	  JSON_OBJECT('aval_num', NEW.aval_num,
	              'aval_rfc', NEW.aval_rfc,
	              'aval_nombre', NEW.aval_nombre,
	              'aval_apellido', NEW.aval_apellido,
	              'aval_direccion', NEW.aval_direccion));

	CREATE TRIGGER trg_aval_bitacora_del AFTER DELETE ON aval
	FOR EACH ROW CALL registrar_bitacora('aval', 'DELETE',
	  JSON_OBJECT('aval_num', OLD.aval_num,
	              'aval_rfc', OLD.aval_rfc,
	              'aval_nombre', OLD.aval_nombre,
	              'aval_apellido', OLD.aval_apellido,
	              'aval_direccion', OLD.aval_direccion),
	  NULL);

	-- inversion
	CREATE TRIGGER trg_inversion_bitacora_ins AFTER INSERT ON inversion
	FOR EACH ROW CALL registrar_bitacora('inversion', 'INSERT',
	  NULL,
	  JSON_OBJECT('inversion_num', NEW.inversion_num,
	              'pagare_num', NEW.pagare_num,
	              'sucursal_num', NEW.sucursal_num,
	              'cliente_num', NEW.cliente_num,
	              'aval_num', NEW.aval_num,
	              'inversion_capital', NEW.inversion_capital,
	              'inversion_tasa', NEW.inversion_tasa,
	              'inversion_tasa_moratoria', NEW.inversion_tasa_moratoria,
	              'inversion_dia_pago', NEW.inversion_dia_pago,
	              'pagare_fecha', NEW.pagare_fecha,
	              'inversion_moneda', NEW.inversion_moneda,
	              'pagare_jurisdiccion', NEW.pagare_jurisdiccion,
	              'inversion_fecha_vencimiento', NEW.inversion_fecha_vencimiento));

	CREATE TRIGGER trg_inversion_bitacora_upd AFTER UPDATE ON inversion
	FOR EACH ROW CALL registrar_bitacora('inversion', 'UPDATE',
	  JSON_OBJECT('inversion_num', OLD.inversion_num,
	              'pagare_num', OLD.pagare_num,
	              'sucursal_num', OLD.sucursal_num,
	              'cliente_num', OLD.cliente_num,
	              'aval_num', OLD.aval_num,
	              'inversion_capital', OLD.inversion_capital,
	              'inversion_tasa', OLD.inversion_tasa,
	              'inversion_tasa_moratoria', OLD.inversion_tasa_moratoria,
	              'inversion_dia_pago', OLD.inversion_dia_pago,
	              'pagare_fecha', OLD.pagare_fecha,
	              'inversion_moneda', OLD.inversion_moneda,
	              'pagare_jurisdiccion', OLD.pagare_jurisdiccion,
	              'inversion_fecha_vencimiento', OLD.inversion_fecha_vencimiento),
	  JSON_OBJECT('inversion_num', NEW.inversion_num,
	              'pagare_num', NEW.pagare_num,
	              'sucursal_num', NEW.sucursal_num,
	              'cliente_num', NEW.cliente_num,
	              'aval_num', NEW.aval_num,
	              'inversion_capital', NEW.inversion_capital,
	              'inversion_tasa', NEW.inversion_tasa,
	              'inversion_tasa_moratoria', NEW.inversion_tasa_moratoria,
	              'inversion_dia_pago', NEW.inversion_dia_pago,
	              'pagare_fecha', NEW.pagare_fecha,
	              'inversion_moneda', NEW.inversion_moneda,
	              'pagare_jurisdiccion', NEW.pagare_jurisdiccion,
	              'inversion_fecha_vencimiento', NEW.inversion_fecha_vencimiento));

	CREATE TRIGGER trg_inversion_bitacora_del AFTER DELETE ON inversion
	FOR EACH ROW CALL registrar_bitacora('inversion', 'DELETE',
	  JSON_OBJECT('inversion_num', OLD.inversion_num,
	              'pagare_num', OLD.pagare_num,
	              'sucursal_num', OLD.sucursal_num,
	              'cliente_num', OLD.cliente_num,
	              'aval_num', OLD.aval_num,
	              'inversion_capital', OLD.inversion_capital,
	              'inversion_tasa', OLD.inversion_tasa,
	              'inversion_tasa_moratoria', OLD.inversion_tasa_moratoria,
	              'inversion_dia_pago', OLD.inversion_dia_pago,
	              'pagare_fecha', OLD.pagare_fecha,
	              'inversion_moneda', OLD.inversion_moneda,
	              'pagare_jurisdiccion', OLD.pagare_jurisdiccion,
	              'inversion_fecha_vencimiento', OLD.inversion_fecha_vencimiento),
	  NULL);

	-- interes
	CREATE TRIGGER trg_interes_bitacora_ins AFTER INSERT ON interes
	FOR EACH ROW CALL registrar_bitacora('interes', 'INSERT',
	  NULL,
	  JSON_OBJECT('inversion_num', NEW.inversion_num,
	              'interes_fecha_inicial', NEW.interes_fecha_inicial,
	              'interes_fecha_final', NEW.interes_fecha_final,
	              'interes_monto', NEW.interes_monto));

	CREATE TRIGGER trg_interes_bitacora_upd AFTER UPDATE ON interes
	FOR EACH ROW CALL registrar_bitacora('interes', 'UPDATE',
	  JSON_OBJECT('inversion_num', OLD.inversion_num,
	              'interes_fecha_inicial', OLD.interes_fecha_inicial,
	              'interes_fecha_final', OLD.interes_fecha_final,
	              'interes_monto', OLD.interes_monto),
	  JSON_OBJECT('inversion_num', NEW.inversion_num,
	              'interes_fecha_inicial', NEW.interes_fecha_inicial,
	              'interes_fecha_final', NEW.interes_fecha_final,
	              'interes_monto', NEW.interes_monto));

	CREATE TRIGGER trg_interes_bitacora_del AFTER DELETE ON interes
	FOR EACH ROW CALL registrar_bitacora('interes', 'DELETE',
	  JSON_OBJECT('inversion_num', OLD.inversion_num,
	              'interes_fecha_inicial', OLD.interes_fecha_inicial,
	              'interes_fecha_final', OLD.interes_fecha_final,
	              'interes_monto', OLD.interes_monto),
	  NULL);

	-- pago
	CREATE TRIGGER trg_pago_bitacora_ins AFTER INSERT ON pago
	FOR EACH ROW CALL registrar_bitacora('pago', 'INSERT',
	  NULL,
	  JSON_OBJECT('inversion_num', NEW.inversion_num,
	              'interes_fecha_inicial', NEW.interes_fecha_inicial,
	              'pago_fecha', NEW.pago_fecha,
	              'pago_interes', NEW.pago_interes,
	              'pago_mora', NEW.pago_mora));

	CREATE TRIGGER trg_pago_bitacora_upd AFTER UPDATE ON pago
	FOR EACH ROW CALL registrar_bitacora('pago', 'UPDATE',
	  JSON_OBJECT('inversion_num', OLD.inversion_num,
	              'interes_fecha_inicial', OLD.interes_fecha_inicial,
	              'pago_fecha', OLD.pago_fecha,
	              'pago_interes', OLD.pago_interes,
	              'pago_mora', OLD.pago_mora),
	  JSON_OBJECT('inversion_num', NEW.inversion_num,
	              'interes_fecha_inicial', NEW.interes_fecha_inicial,
	              'pago_fecha', NEW.pago_fecha,
	              'pago_interes', NEW.pago_interes,
	              'pago_mora', NEW.pago_mora));

	CREATE TRIGGER trg_pago_bitacora_del AFTER DELETE ON pago
	FOR EACH ROW CALL registrar_bitacora('pago', 'DELETE',
	  JSON_OBJECT('inversion_num', OLD.inversion_num,
	              'interes_fecha_inicial', OLD.interes_fecha_inicial,
	              'pago_fecha', OLD.pago_fecha,
	              'pago_interes', OLD.pago_interes,
	              'pago_mora', OLD.pago_mora),
	  NULL);

	-- liquidacion
	CREATE TRIGGER trg_liquidacion_bitacora_ins AFTER INSERT ON liquidacion
	FOR EACH ROW CALL registrar_bitacora('liquidacion', 'INSERT',
	  NULL,
	  JSON_OBJECT('inversion_num', NEW.inversion_num,
	              'liquidacion_fecha', NEW.liquidacion_fecha,
	              'liquidacion_capital', NEW.liquidacion_capital,
	              'liquidacion_mora', NEW.liquidacion_mora));

	CREATE TRIGGER trg_liquidacion_bitacora_upd AFTER UPDATE ON liquidacion
	FOR EACH ROW CALL registrar_bitacora('liquidacion', 'UPDATE',
	  JSON_OBJECT('inversion_num', OLD.inversion_num,
	              'liquidacion_fecha', OLD.liquidacion_fecha,
	              'liquidacion_capital', OLD.liquidacion_capital,
	              'liquidacion_mora', OLD.liquidacion_mora),
	  JSON_OBJECT('inversion_num', NEW.inversion_num,
	              'liquidacion_fecha', NEW.liquidacion_fecha,
	              'liquidacion_capital', NEW.liquidacion_capital,
	              'liquidacion_mora', NEW.liquidacion_mora));

	CREATE TRIGGER trg_liquidacion_bitacora_del AFTER DELETE ON liquidacion
	FOR EACH ROW CALL registrar_bitacora('liquidacion', 'DELETE',
	  JSON_OBJECT('inversion_num', OLD.inversion_num,
	              'liquidacion_fecha', OLD.liquidacion_fecha,
	              'liquidacion_capital', OLD.liquidacion_capital,
	              'liquidacion_mora', OLD.liquidacion_mora),
	  NULL);

	-- Protección de la bitácora
	CREATE TRIGGER trg_bitacora_proteger_upd BEFORE UPDATE ON bitacora
	FOR EACH ROW SIGNAL SQLSTATE '45000'
	  SET MESSAGE_TEXT = 'La bitácora no admite la operación UPDATE';

	CREATE TRIGGER trg_bitacora_proteger_del BEFORE DELETE ON bitacora
	FOR EACH ROW SIGNAL SQLSTATE '45000'
	  SET MESSAGE_TEXT = 'La bitácora no admite la operación DELETE';
\end{verbatim}

\captura{mysql-bitacora-triggers.png}{Triggers de bitácora en MySQL}{mysql-bitacora-triggers}

\section{Registro de las pruebas de alcance}

Las pruebas de alcance de la sección~\ref{sec:pruebas-alcance} se ejecutan después de crear la bitácora, de modo que sus operaciones permitidas quedan registradas con el usuario de cada prueba. Las operaciones denegadas no dejan registro. La tabla~\ref{tab:bitacora-pruebas} muestra los 17 registros esperados. La numeración puede tener huecos, porque una inserción revertida también consume un número de \nem{bitacora\us num}.

\begin{table}[H]
	\centering
	\small
	\caption{Registros de bitácora esperados tras las pruebas de alcance}
	\label{tab:bitacora-pruebas}
	\begin{tabular}{@{}l l l r L{6.2cm}@{}}
		\toprule
		\textbf{Usuario} & \textbf{Tabla} & \textbf{Operación} &
		\textbf{Registros} & \textbf{Origen} \\
		\midrule
		\nem{dmejia} & \nem{sucursal} & \nem{UPDATE} & 1 & Dirección de la
		sucursal 15 \\
		\nem{pzuniga} & \nem{cliente} & \nem{INSERT} & 1 & Cliente 26 \\
		\nem{pzuniga} & \nem{inversion} & \nem{INSERT} & 1 & Inversión 41 \\
		\nem{pzuniga} & \nem{interes} & \nem{INSERT} & 12 & Periodos generados
		por el procedimiento \\
		\nem{pzuniga} & \nem{cliente} & \nem{UPDATE} & 1 & Teléfono del
		cliente 26 \\
		\nem{xzaleta} & \nem{pago} & \nem{INSERT} & 1 & Pago tardío de la
		inversión 16 \\
		\bottomrule
	\end{tabular}
\end{table}

El auditor consulta los registros con:

\begin{verbatim}
	SELECT bitacora_num, bitacora_fecha, bitacora_usuario, bitacora_ip,
	       bitacora_sesion, bitacora_tabla, bitacora_operacion
	FROM bitacora
	ORDER BY bitacora_num;

	SELECT bitacora_dato_anterior, bitacora_dato_nuevo
	FROM bitacora
	WHERE bitacora_tabla = 'cliente' AND bitacora_operacion = 'UPDATE';
\end{verbatim}

\captura{pg-bitacora-registros.png}{Registros de bitácora en PostgreSQL}{pg-bitacora-registros}
\captura{mysql-bitacora-registros.png}{Registros de bitácora en MySQL}{mysql-bitacora-registros}
	\chapter{Tablespaces}
\label{cap:tablespaces}

Un tablespace es una ubicación en disco donde el gestor guarda los datos de los objetos que se le asignan. En PostgreSQL, un tablespace es un directorio; en MySQL, un tablespace general es un archivo \nem{.ibd} que comparten varias tablas. Si no se indica otra cosa, PostgreSQL guarda todos los objetos en el tablespace \nem{pg\us default}, dentro del directorio de datos, y MySQL guarda cada tabla en su propio archivo, también dentro del directorio de datos.

\section{Distribuci\'on de las tablas}

\paragraph{Criterio.} Las tablas se agrupan en tres tablespaces según su uso y su ritmo de crecimiento (tabla~\ref{tab:distribucion-tablespaces}). La distribución es la misma en ambos gestores.

\begin{table}[H]
	\centering
	\small
	\caption{Distribución de las tablas en los tablespaces}
	\label{tab:distribucion-tablespaces}
	\begin{tabular}{@{}L{2.6cm} L{5cm} L{8.2cm}@{}}
		\toprule
		\textbf{Tablespace} & \textbf{Tablas} & \textbf{Criterio} \\
		\midrule
		\nem{ts\us catalogos} & \nem{empresa}, \nem{sucursal}, \nem{cliente}, \nem{aval} & Datos maestros: pocos registros que cambian ocasionalmente. \\
		\nem{ts\us operacion} & \nem{inversion}, \nem{interes}, \nem{pago}, \nem{liquidacion} & Datos de la operación: crecen con cada inversión y con cada mes de pagos. \\
		\nem{ts\us bitacora} & \nem{bitacora} & Registro de auditoría: crece con cada cambio en los datos y sus registros nunca se modifican ni se eliminan. \\
		\bottomrule
	\end{tabular}
\end{table}

\paragraph{Justificación.} Los tres tablespaces se ubican en el mismo disco que el directorio de datos, por lo que separarlos no mejora el rendimiento: todas las lecturas y escrituras compiten por el mismo dispositivo. El beneficio es administrativo:

\begin{itemize}[leftmargin=*]
	\item El tamaño de cada tablespace indica cuánto crece cada grupo de tablas. En particular, permite vigilar la bitácora, que es la tabla que crece de forma continua.
	\item Los archivos de cada grupo se encuentran en un directorio conocido, fuera del directorio de datos.
	\item Si en el futuro el equipo cuenta con otro disco, un grupo de tablas puede trasladarse a él cambiando solo la ubicación de su tablespace, sin modificar el esquema ni las consultas.
\end{itemize}

\paragraph{Índices.} Cada índice se ubica en el mismo tablespace que su tabla. En MySQL no existe otra opción, ya que InnoDB guarda los índices de una tabla en el mismo tablespace que sus datos. PostgreSQL sí permite separarlos, pero se mantienen junto a su tabla para que la distribución sea igual en ambos gestores.

\paragraph{Responsabilidades.} De acuerdo con la sección~\ref{sec:admin-instancia}, el superusuario prepara los directorios y crea los tablespaces, porque ambos gestores reservan esa operación a cuentas con privilegios sobre toda la instancia. El DBA asigna las tablas a los tablespaces. En PostgreSQL, para poder hacerlo, el DBA debe ser dueño de los tablespaces; en MySQL, los tablespaces no tienen dueño y basta con los privilegios del DBA sobre la base de datos.

\paragraph{Momento de la asignación.} Las tablas se crean como se describe en la implementación física, en la ubicación por defecto, y después el DBA las traslada a su tablespace. Así, los scripts de creación son los mismos con o sin tablespaces, y en ambas instancias. El traslado copia los datos de la tabla y la bloquea mientras dura la copia; con el volumen de este sistema es inmediato, pero en una base de datos en operación con tablas grandes conviene crear las tablas directamente en su tablespace.

\paragraph{Segunda instancia.} La segunda instancia de cada gestor tiene sus tablespaces en un directorio propio. Ahí, las tablas \nem{pago} y \nem{bitacora} están particionadas, lo que produce la única diferencia entre gestores (tabla~\ref{tab:tablespaces-instancias}):

\begin{itemize}[leftmargin=*]
	\item En PostgreSQL, cada partición puede ubicarse en un tablespace, por lo que las particiones de \nem{pago} se ubican en \nem{ts\us operacion} y las de \nem{bitacora} en \nem{ts\us bitacora}. Su creación se describe en el capítulo de particionamiento.
	\item En MySQL, una tabla particionada no puede ubicarse en un tablespace general: el gestor rechaza la operación con el error 1478 (\nem{A partitioned table is not allowed in a shared tablespace}). La única alternativa es indicar un directorio para cada partición con la opción \nem{DATA DIRECTORY}, que requiere el privilegio global \nem{FILE}. Ese privilegio permite leer y escribir cualquier archivo al que tenga acceso el servidor, por lo que no se otorga al DBA; y que el superusuario creara las tablas contradiría la división de responsabilidades. Por ello, en la segunda instancia de MySQL, \nem{pago} y \nem{bitacora} permanecen en archivos individuales dentro del directorio de datos, y el tablespace \nem{ts\us bitacora} no se crea, ya que no contendría ninguna tabla.
\end{itemize}

\begin{table}[H]
	\centering
	\small
	\caption{Ubicación de las tablas en cada instancia}
	\label{tab:tablespaces-instancias}
	\begin{tabular}{@{}L{3.6cm} L{3.4cm} L{3.4cm} L{4.6cm}@{}}
		\toprule
		\textbf{Tablas} & \textbf{Primera instancia (ambos gestores)} & \textbf{Segunda instancia de PostgreSQL} & \textbf{Segunda instancia de MySQL} \\
		\midrule
		\nem{empresa}, \nem{sucursal}, \nem{cliente}, \nem{aval} & \nem{ts\us catalogos} & \nem{ts\us catalogos} & \nem{ts\us catalogos} \\
		\nem{inversion}, \nem{interes}, \nem{liquidacion} & \nem{ts\us operacion} & \nem{ts\us operacion} & \nem{ts\us operacion} \\
		\nem{pago} & \nem{ts\us operacion} & \nem{ts\us operacion} (cada partición) & Un archivo por partición en el directorio de datos \\
		\nem{bitacora} & \nem{ts\us bitacora} & \nem{ts\us bitacora} (cada partición) & Un archivo por partición en el directorio de datos \\
		\bottomrule
	\end{tabular}
\end{table}

\paragraph{Directorios.} Todos los tablespaces se ubican bajo \nem{C:\bs tablespaces}, con una carpeta por instancia y, dentro de ella, una subcarpeta por tablespace (tabla~\ref{tab:directorios-tablespaces} y figura~\ref{fig:tablespaces-archivos}).

\begin{table}[H]
	\centering
	\small
	\caption{Directorios de los tablespaces}
	\label{tab:directorios-tablespaces}
	\begin{tabular}{@{}L{4.6cm} L{11.2cm}@{}}
		\toprule
		\textbf{Instancia} & \textbf{Directorios} \\
		\midrule
		Primera de PostgreSQL & \nem{C:\bs tablespaces\bs pg1\bs catalogos}, \nem{\bs operacion}, \nem{\bs bitacora} \\
		Segunda de PostgreSQL & \nem{C:\bs tablespaces\bs pg2\bs catalogos}, \nem{\bs operacion}, \nem{\bs bitacora} \\
		Primera de MySQL & \nem{C:\bs tablespaces\bs mysql1\bs catalogos}, \nem{\bs operacion}, \nem{\bs bitacora} \\
		Segunda de MySQL & \nem{C:\bs tablespaces\bs mysql2\bs catalogos}, \nem{\bs operacion} \\
		\bottomrule
	\end{tabular}
\end{table}

Las secciones siguientes muestran los scripts de la primera instancia. En la segunda instancia son los mismos, cambiando \nem{pg1} por \nem{pg2} o \nem{mysql1} por \nem{mysql2}, y omitiendo \nem{pago} y \nem{bitacora}, que se ubican al particionarlas; en MySQL se omite también la creación de \nem{ts\us bitacora}.

\section{PostgreSQL}

\subsection{Directorios}

Cada tablespace de PostgreSQL es un directorio que debe existir y estar vacío al crearlo. El servicio del gestor se ejecuta con su propia cuenta de Windows, \nem{NT AUTHORITY\bs NetworkService} con el instalador por defecto, y no con la del usuario que administra el equipo, por lo que esa cuenta debe tener control total sobre los directorios. La cuenta del servicio se verifica en la columna \textit{Iniciar sesión como} de la consola de servicios (\nem{services.msc}).

Los directorios se crean desde una terminal ejecutada como administrador:

\begin{verbatim}
	mkdir C:\tablespaces\pg1\catalogos
	mkdir C:\tablespaces\pg1\operacion
	mkdir C:\tablespaces\pg1\bitacora
	icacls C:\tablespaces\pg1 /grant "NT AUTHORITY\NetworkService:(OI)(CI)F"
\end{verbatim}

Los indicadores \nem{(OI)(CI)} hacen que el permiso se herede a las subcarpetas y a los archivos que el gestor cree en ellas, y \nem{F} otorga control total.

\subsection{Creaci\'on de los tablespaces}

El superusuario crea los tablespaces y asigna como dueño al rol del DBA. Sin ese dueño, el DBA no podría trasladar tablas a ellos: el gestor rechaza la operación con el error \nem{permission denied for tablespace}. Las rutas se escriben con diagonales normales, que PostgreSQL acepta en Windows.

\begin{verbatim}
	-- Ejecutado por postgres
	CREATE TABLESPACE ts_catalogos OWNER rol_admin_bd
	  LOCATION 'C:/tablespaces/pg1/catalogos';
	CREATE TABLESPACE ts_operacion OWNER rol_admin_bd
	  LOCATION 'C:/tablespaces/pg1/operacion';
	CREATE TABLESPACE ts_bitacora OWNER rol_admin_bd
	  LOCATION 'C:/tablespaces/pg1/bitacora';
\end{verbatim}

Dentro de cada directorio, el gestor crea una subcarpeta propia de su versión (\nem{PG\us 17\us }\ldots) y, dentro de ella, una carpeta por base de datos, con un archivo por cada tabla e índice, nombrado con su identificador interno.

\subsection{Asignaci\'on de las tablas}

El DBA traslada cada tabla con \nem{ALTER TABLE \ldots{} SET TABLESPACE}. Esta sentencia traslada la tabla y su almacenamiento auxiliar (TOAST), donde PostgreSQL guarda los valores grandes, como los textos y JSON de la bitácora; pero \textbf{no traslada sus índices}, que deben moverse por separado con \nem{ALTER INDEX}. Los índices del esquema son los que crean las llaves primarias y alternas.

\begin{verbatim}
	-- Ejecutado por dlandero en la base de datos inversiones
	SET ROLE rol_admin_bd;

	-- Catálogos
	ALTER TABLE empresa  SET TABLESPACE ts_catalogos;
	ALTER TABLE sucursal SET TABLESPACE ts_catalogos;
	ALTER TABLE cliente  SET TABLESPACE ts_catalogos;
	ALTER TABLE aval     SET TABLESPACE ts_catalogos;
	ALTER INDEX pk_empresa     SET TABLESPACE ts_catalogos;
	ALTER INDEX pk_sucursal    SET TABLESPACE ts_catalogos;
	ALTER INDEX pk_cliente     SET TABLESPACE ts_catalogos;
	ALTER INDEX ak_cliente_rfc SET TABLESPACE ts_catalogos;
	ALTER INDEX pk_aval        SET TABLESPACE ts_catalogos;
	ALTER INDEX ak_aval_rfc    SET TABLESPACE ts_catalogos;

	-- Operación
	ALTER TABLE inversion   SET TABLESPACE ts_operacion;
	ALTER TABLE interes     SET TABLESPACE ts_operacion;
	ALTER TABLE pago        SET TABLESPACE ts_operacion;
	ALTER TABLE liquidacion SET TABLESPACE ts_operacion;
	ALTER INDEX pk_inversion        SET TABLESPACE ts_operacion;
	ALTER INDEX ak_inversion_pagare SET TABLESPACE ts_operacion;
	ALTER INDEX pk_interes          SET TABLESPACE ts_operacion;
	ALTER INDEX pk_pago             SET TABLESPACE ts_operacion;
	ALTER INDEX pk_liquidacion      SET TABLESPACE ts_operacion;

	-- Bitácora
	ALTER TABLE bitacora    SET TABLESPACE ts_bitacora;
	ALTER INDEX pk_bitacora SET TABLESPACE ts_bitacora;
\end{verbatim}

\section{MySQL}

\subsection{Directorios}

MySQL solo crea tablespaces fuera del directorio de datos en las rutas que se declaran en el parámetro \nem{innodb\us directories}; en cualquier otra, rechaza la operación con el error 3121 (\nem{The DATAFILE location must be in a known directory}). El parámetro acepta también las subcarpetas de la ruta declarada, por lo que basta con indicar la carpeta de la instancia. Se agrega en \nem{my.ini}, dentro de la sección \nem{[mysqld]}, y requiere reiniciar el servicio:

\begin{verbatim}
	innodb_directories="C:/tablespaces/mysql1"
\end{verbatim}

Las rutas de MySQL se escriben con diagonales normales, porque en sus cadenas de texto la diagonal invertida es un carácter de escape. Los directorios se crean y se les otorgan permisos igual que en PostgreSQL, para la cuenta con la que se ejecuta el servicio de MySQL, que también se verifica en \nem{services.msc}:

\begin{verbatim}
	mkdir C:\tablespaces\mysql1\catalogos
	mkdir C:\tablespaces\mysql1\operacion
	mkdir C:\tablespaces\mysql1\bitacora
	icacls C:\tablespaces\mysql1 /grant "NT AUTHORITY\NetworkService:(OI)(CI)F"
\end{verbatim}

\subsection{Creaci\'on de los tablespaces}

El superusuario crea los tablespaces, ya que la operación requiere el privilegio global \nem{CREATE TABLESPACE}. Cada tablespace es un archivo \nem{.ibd} dentro de su directorio.

\begin{verbatim}
	-- Ejecutado por root
	CREATE TABLESPACE ts_catalogos
	  ADD DATAFILE 'C:/tablespaces/mysql1/catalogos/ts_catalogos.ibd';
	CREATE TABLESPACE ts_operacion
	  ADD DATAFILE 'C:/tablespaces/mysql1/operacion/ts_operacion.ibd';
	CREATE TABLESPACE ts_bitacora
	  ADD DATAFILE 'C:/tablespaces/mysql1/bitacora/ts_bitacora.ibd';
\end{verbatim}

\subsection{Asignaci\'on de las tablas}

El DBA traslada cada tabla con \nem{ALTER TABLE \ldots{} TABLESPACE}. Los índices se trasladan junto con la tabla, y la operación no requiere privilegios sobre el tablespace, ya que en MySQL estos no tienen dueño.

\begin{verbatim}
	-- Ejecutado por dlandero en la base de datos inversiones
	-- Catálogos
	ALTER TABLE empresa  TABLESPACE ts_catalogos;
	ALTER TABLE sucursal TABLESPACE ts_catalogos;
	ALTER TABLE cliente  TABLESPACE ts_catalogos;
	ALTER TABLE aval     TABLESPACE ts_catalogos;

	-- Operación
	ALTER TABLE inversion   TABLESPACE ts_operacion;
	ALTER TABLE interes     TABLESPACE ts_operacion;
	ALTER TABLE pago        TABLESPACE ts_operacion;
	ALTER TABLE liquidacion TABLESPACE ts_operacion;

	-- Bitácora
	ALTER TABLE bitacora TABLESPACE ts_bitacora;
\end{verbatim}

La figura~\ref{fig:tablespaces-archivos} muestra el contenido de \nem{C:\bs tablespaces} con las tablas ya asignadas en las cuatro instancias, incluidas las particiones de la segunda instancia de PostgreSQL.

\captura{tablespaces-archivos.png}{Directorios de los tablespaces de todas las instancias}{tablespaces-archivos}

	\chapter{Multi-instancias}
\label{cap:multiinstancias}
\recursoscapitulo{10-multiinstancias}

Una instancia es un servidor de base de datos en ejecución: un proceso del gestor con su propio directorio de datos, su configuración y su puerto. Un mismo equipo puede ejecutar varias instancias del mismo gestor, siempre que cada una tenga su propio directorio de datos y escuche en un puerto distinto.

El sistema utiliza dos instancias de cada gestor. La primera contiene el esquema descrito en los capítulos anteriores. La segunda se construye como una copia exacta de la primera, con el mismo esquema y los mismos datos, y después se le aplican el particionamiento y la segmentación descritos en los capítulos siguientes. Así, ambas instancias parten del mismo estado, y las diferencias entre ellas se deben únicamente a esas dos técnicas.

\section{Arquitectura de las instancias}

Las dos instancias de cada gestor utilizan los mismos programas, instalados una sola vez; lo que las distingue es su directorio de datos, su configuración, su puerto y el servicio de Windows que las ejecuta (tabla~\ref{tab:instancias}). La primera instancia de cada gestor es la que creó el instalador, con las rutas descritas en la configuración del entorno. La segunda se crea manualmente, con sus archivos bajo \nem{C:\bs instancias}.

\begin{table}[H]
	\centering
	\small
	\caption{Instancias del sistema}
	\label{tab:instancias}
	\begin{tabular}{@{}L{2.7cm} L{3.1cm} L{3.1cm} L{3.1cm} L{3.1cm}@{}}
		\toprule
		& \textbf{PostgreSQL 1} & \textbf{PostgreSQL 2} & \textbf{MySQL 1} & \textbf{MySQL 2} \\
		\midrule
		Puerto & 5432 & 5434 & 3306 & 3308 \\
		Puerto del protocolo X & --- & --- & 33060 & 33070 \\
		Directorio de datos & El del instalador & \nem{C:\bs instancias\bs pg2\bs data} & El del instalador & \nem{C:\bs instancias\bs mysql2\bs data} \\
		Configuración & \nem{postgresql.conf} del directorio de datos & \nem{postgresql.conf} del directorio de datos & \nem{my.ini} del instalador & \nem{C:\bs instancias\bs mysql2\bs my.ini} \\
		Tablespaces & \nem{C:\bs tablespaces\bs pg1} & \nem{C:\bs tablespaces\bs pg2} & \nem{C:\bs tablespaces\bs mysql1} & \nem{C:\bs tablespaces\bs mysql2} \\
		Servicio & \nem{postgresql-\allowbreak x64-17} & \nem{postgresql-\allowbreak x64-17-2} & \nem{MySQL80} & \nem{MySQL80-2} \\
		Esquema & Base & Particionado y segmentado & Base & Particionado y segmentado \\
		\bottomrule
	\end{tabular}
\end{table}

\paragraph{Configuración idéntica.} Como se indicó en la configuración del entorno, la segunda instancia de cada gestor utiliza los mismos parámetros que la primera, excepto los que identifican a la instancia y no pueden repetirse en un mismo equipo. Para garantizarlo, su archivo de configuración no se escribe de nuevo: se copia el de la primera instancia y se modifican únicamente esos valores. Los archivos de registro de cada instancia quedan separados sin cambios adicionales, porque sus rutas son relativas al directorio de datos.

\paragraph{Conexión.} Los clientes eligen la instancia por su puerto: \nem{psql -p 5434} o \nem{mysql -P 3308}. Las reglas de acceso son las mismas en ambas instancias, porque \nem{pg\us hba.conf} se copia de la primera instancia y las cuentas de MySQL se crean con los mismos scripts. Para conectarse desde otro equipo, el firewall de Windows debe permitir también los puertos 5434 y 3308.

\paragraph{Responsabilidades.} De acuerdo con la sección~\ref{sec:admin-instancia}, el superusuario crea y configura las instancias, y el DBA construye el esquema y traslada los datos.

\section{Segunda instancia de PostgreSQL}

\paragraph{Directorio.} El superusuario del equipo crea el directorio de la instancia y otorga control total a la cuenta del servicio, igual que en los directorios de los tablespaces. Los comandos de este capítulo se ejecutan en una terminal de Windows (\nem{cmd}) abierta como administrador; los que no caben en una línea se dividen con el carácter \nem{\textasciicircum}, que indica que el comando continúa en la línea siguiente. El directorio de datos se crea dentro de él en el paso siguiente y hereda el permiso.

\begin{verbatim}
	mkdir C:\instancias\pg2
	icacls C:\instancias\pg2 /grant "NT AUTHORITY\NetworkService:(OI)(CI)F"
\end{verbatim}

\paragraph{Inicialización.} El programa \nem{initdb} crea el directorio de datos con su catálogo, las bases de datos iniciales y el superusuario \nem{postgres}, con la misma codificación y el mismo método de autenticación que la primera instancia. La opción \nem{-W} solicita la contraseña del superusuario.

\begin{verbatim}
	initdb -D C:\instancias\pg2\data -U postgres -E UTF8 -A scram-sha-256 -W
\end{verbatim}

\paragraph{Configuración.} El archivo \nem{postgresql.conf} que genera \nem{initdb} no coincide con el que deja el instalador: por ejemplo, \nem{listen\us addresses} se limita a \nem{localhost} y \nem{logging\us collector} está desactivado. Por ello, se reemplazan \nem{postgresql.conf} y \nem{pg\us hba.conf} por los de la primera instancia:

\begin{verbatim}
	copy /Y "C:\Program Files\PostgreSQL\17\data\postgresql.conf" ^
	  C:\instancias\pg2\data
	copy /Y "C:\Program Files\PostgreSQL\17\data\pg_hba.conf" ^
	  C:\instancias\pg2\data
\end{verbatim}

En la copia de \nem{postgresql.conf} se cambia únicamente el puerto:

\begin{verbatim}
	port = 5434
\end{verbatim}

\paragraph{Servicio.} El programa \nem{pg\us ctl} registra la instancia como servicio de Windows, con inicio automático y la misma cuenta que el servicio de la primera instancia:

\begin{verbatim}
	pg_ctl register -N postgresql-x64-17-2 -D C:\instancias\pg2\data ^
	  -U "NT AUTHORITY\NetworkService" -S auto
	net start postgresql-x64-17-2
\end{verbatim}

\paragraph{Verificación.} El superusuario se conecta al puerto de la segunda instancia y consulta sus parámetros (figura~\ref{fig:pg-instancia2}). El puerto y el directorio de datos son propios de la instancia; los demás parámetros coinciden con los de la primera.

\begin{verbatim}
	SELECT name, setting
	FROM pg_settings
	WHERE name IN ('port', 'data_directory', 'listen_addresses',
	               'logging_collector', 'log_statement', 'max_connections',
	               'shared_buffers', 'TimeZone')
	ORDER BY name;
\end{verbatim}

\captura{pg-instancia2.png}{Parámetros de la segunda instancia de PostgreSQL}{pg-instancia2}

\section{Segunda instancia de MySQL}

\paragraph{Directorio y configuración.} El superusuario del equipo crea el directorio de la instancia, le otorga permisos y copia en él el archivo \nem{my.ini} de la primera instancia:

\begin{verbatim}
	mkdir C:\instancias\mysql2
	icacls C:\instancias\mysql2 /grant "NT AUTHORITY\NetworkService:(OI)(CI)F"
	copy "C:\ProgramData\MySQL\MySQL Server 8.0\my.ini" C:\instancias\mysql2
\end{verbatim}

En la copia se modifican, dentro de la sección \nem{[mysqld]}, los valores de la tabla~\ref{tab:mysql2-parametros}. El directorio de datos no debe existir: lo crea la inicialización.

\begin{table}[H]
	\centering
	\small
	\caption{Parámetros propios de la segunda instancia de MySQL}
	\label{tab:mysql2-parametros}
	\begin{tabular}{@{}L{3.4cm} L{5.2cm} L{7.2cm}@{}}
		\toprule
		\textbf{Parámetro} & \textbf{Valor} & \textbf{Motivo} \\
		\midrule
		\nem{port} & \nem{3308} & Cada instancia escucha en su propio puerto. \\
		\nem{datadir} & \nem{"C:/instancias/\allowbreak mysql2/data"} & Dos instancias no pueden compartir el directorio de datos. \\
		\nem{innodb\us directories} & \nem{"C:/tablespaces/\allowbreak mysql2"} & Los tablespaces de la instancia (capítulo~\ref{cap:tablespaces}). \\
		\nem{mysqlx\us port} & \nem{33070} & El instalador activa el protocolo X en el puerto 33060; si ambas instancias lo usaran, la segunda no podría iniciar. \\
		\bottomrule
	\end{tabular}
\end{table}

\paragraph{Inicialización.} El programa \nem{mysqld}, con la opción \nem{-{}-initialize}, crea el directorio de datos y la cuenta \nem{'root'@'localhost'} con una contraseña temporal. Sin embargo, el \nem{my.ini} copiado ya activa \nem{skip\us name\us resolve}, de modo que esa cuenta no podría conectarse desde el cliente de Windows (sección~\ref{sec:conexion-mysql}). Por ello, la inicialización ejecuta un archivo con la opción \nem{-{}-init-file}, que asigna la contraseña definitiva de \nem{root} y crea sus cuentas de loopback antes del primer inicio:

\begin{verbatim}
	-- C:\instancias\mysql2\inicial.sql
	ALTER USER 'root'@'localhost' IDENTIFIED BY '********';
	CREATE USER 'root'@'127.0.0.1' IDENTIFIED BY '********';
	GRANT ALL PRIVILEGES ON *.* TO 'root'@'127.0.0.1' WITH GRANT OPTION;
	CREATE USER 'root'@'::1' IDENTIFIED BY '********';
	GRANT ALL PRIVILEGES ON *.* TO 'root'@'::1' WITH GRANT OPTION;
\end{verbatim}

\begin{verbatim}
	mysqld --defaults-file="C:\instancias\mysql2\my.ini" --initialize ^
	  --init-file="C:/instancias/mysql2/inicial.sql" --console
	del C:\instancias\mysql2\inicial.sql
\end{verbatim}

La opción \nem{-{}-defaults-file} debe ser la primera de la línea de comandos. El archivo se elimina después de la inicialización, porque contiene la contraseña del superusuario en texto plano.

\paragraph{Servicio.} El servicio se registra con el mismo programa, indicando su archivo de configuración. MySQL registra sus servicios con la cuenta del sistema local, por lo que se cambia a la misma cuenta que la primera instancia:

\begin{verbatim}
	mysqld --install MySQL80-2 --defaults-file="C:\instancias\mysql2\my.ini"
	sc config MySQL80-2 obj= "NT AUTHORITY\NetworkService" password= ""
	net start MySQL80-2
\end{verbatim}

\paragraph{Verificación.} El superusuario se conecta al puerto de la segunda instancia y consulta sus parámetros (figura~\ref{fig:mysql-instancia2}):

\begin{verbatim}
	SELECT @@port, @@datadir, @@innodb_directories, @@mysqlx_port,
	       @@max_connections, @@skip_name_resolve, @@general_log;
\end{verbatim}

\captura{mysql-instancia2.png}{Parámetros de la segunda instancia de MySQL}{mysql-instancia2}

\section{Carga del esquema y los datos}

La segunda instancia se construye con el mismo procedimiento que la primera, hasta tener el mismo esquema y los mismos datos. El esquema se crea con los scripts de los capítulos anteriores; los datos no se cargan de nuevo, sino que se copian de la primera instancia. La tabla~\ref{tab:secuencia-instancia2} muestra la secuencia completa.

\begin{table}[H]
	\centering
	\small
	\caption{Construcción de la segunda instancia}
	\label{tab:secuencia-instancia2}
	\begin{tabular}{@{}r L{3.6cm} L{11.2cm}@{}}
		\toprule
		& \textbf{Ejecuta} & \textbf{Paso} \\
		\midrule
		1 & Superusuario & Rol del DBA y, en PostgreSQL, la base de datos (secciones~\ref{sec:roles-pg} y~\ref{sec:roles-mysql}, e implementación física). \\
		2 & DBA & Base de datos en MySQL, roles de puesto y usuarios. En MySQL, el superusuario otorga además los privilegios sobre la instancia. \\
		3 & DBA & Esquema, tablas, procedimientos, triggers de validación y permisos (implementación física). \\
		4 & DBA & Copia de los datos de la primera instancia. \\
		5 & DBA & Triggers de bitácora (capítulo~\ref{cap:bitacora}). \\
		6 & Superusuario y DBA & Tablespaces, sin trasladar \nem{pago} ni \nem{bitacora} (capítulo~\ref{cap:tablespaces}). \\
		7 & DBA & Verificación de que ambas instancias contienen los mismos datos. \\
		\bottomrule
	\end{tabular}
\end{table}

A partir de ese punto, la segunda instancia recibe el particionamiento y la segmentación de los capítulos siguientes, que no se aplican a la primera.

\paragraph{Por qué se copian los datos.} Volver a ejecutar la carga de datos de prueba no reproduciría el estado de la primera instancia, que también contiene los registros creados durante las pruebas de alcance: el cliente 26, la inversión 41 con sus periodos de interés, el pago de la inversión 16 y los 17 registros de bitácora que dejaron esas operaciones. La copia se toma después de las pruebas, por lo que la segunda instancia recibe exactamente los mismos registros, con las mismas llaves. Esto es posible gracias a que las columnas de identidad y \nem{AUTO\us INCREMENT} admiten valores explícitos (tabla~\ref{tab:equivalencia-tipos}). Después de la copia, cada gestor continúa la numeración a partir del valor más alto copiado.

\paragraph{Por qué la bitácora se activa después.} Los triggers de bitácora se crean después de la copia, igual que en la primera instancia se crearon después de la carga inicial. Si existieran durante la copia, cada registro copiado generaría un registro de auditoría, como si un usuario lo hubiera insertado. Los triggers de validación, en cambio, sí existen durante la copia y vuelven a validar los datos sin afectarlos.

\paragraph{PostgreSQL.} El DBA genera un respaldo de solo datos de la primera instancia con \nem{pg\us dump} y lo restaura en la segunda con \nem{pg\us restore}. El respaldo incluye el valor de las secuencias de las columnas de identidad, y \nem{pg\us restore} carga las tablas en el orden que exigen sus llaves foráneas.

\begin{verbatim}
	pg_dump -h localhost -p 5432 -U dlandero -d inversiones -n inversiones ^
	  --data-only -F c -f inversiones_datos.dump
	pg_restore -h localhost -p 5434 -U dlandero -d inversiones ^
	  --data-only --exit-on-error inversiones_datos.dump
\end{verbatim}

\paragraph{MySQL.} El DBA genera el respaldo con \nem{mysqldump} y lo ejecuta en la segunda instancia con el cliente \nem{mysql}:

\begin{verbatim}
	mysqldump -h localhost -P 3306 -u dlandero -p --no-create-info --skip-triggers ^
	  --no-tablespaces --single-transaction inversiones > inversiones_datos.sql
	mysql -h localhost -P 3308 -u dlandero -p inversiones < inversiones_datos.sql
\end{verbatim}

\nem{-{}-no-create-info} y \nem{-{}-skip-triggers} limitan el respaldo a los datos, ya que las tablas y los triggers existen en la segunda instancia; \nem{-{}-no-tablespaces} omite la información de los tablespaces, que requeriría el privilegio \nem{PROCESS}; y \nem{-{}-single-transaction} toma una copia consistente sin bloquear las tablas. \nem{mysqldump} escribe las tablas en orden alfabético, por ejemplo \nem{interes} antes que \nem{inversion}, por lo que el respaldo desactiva la verificación de llaves foráneas durante la carga (\nem{FOREIGN\us KEY\us CHECKS = 0}) y la restablece al final.

\paragraph{Verificación.} La siguiente consulta, igual en ambos gestores, se ejecuta en las dos instancias; ambas deben devolver los mismos valores (figuras~\ref{fig:pg-instancias-datos} y~\ref{fig:mysql-instancias-datos}).

\begin{verbatim}
	SELECT (SELECT COUNT(*) FROM cliente) AS clientes,
	       (SELECT COUNT(*) FROM inversion) AS inversiones,
	       (SELECT COUNT(*) FROM interes) AS periodos,
	       (SELECT COUNT(*) FROM pago) AS pagos,
	       (SELECT COUNT(*) FROM liquidacion) AS liquidaciones,
	       (SELECT COUNT(*) FROM bitacora) AS bitacora,
	       (SELECT SUM(pago_interes) + SUM(pago_mora) FROM pago) AS total_pagado;
\end{verbatim}

Ambas instancias contienen 26 clientes, 41 inversiones, 651 periodos de interés, 479 pagos, 15 liquidaciones y 17 registros de bitácora, con el mismo total pagado.

\captura{pg-instancias-datos.png}{Datos de las dos instancias de PostgreSQL}{pg-instancias-datos}
\captura{mysql-instancias-datos.png}{Datos de las dos instancias de MySQL}{mysql-instancias-datos}

	\chapter{Particionamiento}
\label{cap:particionamiento}
\recursoscapitulo{11-particiones}

El particionamiento divide físicamente una tabla en partes llamadas particiones, cada una con un subconjunto de sus filas, según el valor de una columna llamada clave de partición. Para los usuarios, la tabla sigue siendo una sola: se consulta y se modifica con su nombre, y el gestor dirige cada fila a su partición. Cuando una consulta filtra por la clave de partición, el gestor lee únicamente las particiones que pueden contener las filas buscadas, lo que se conoce como poda de particiones.

El particionamiento se aplica en la segunda instancia de cada gestor, que hasta este punto es una copia exacta de la primera (capítulo~\ref{cap:multiinstancias}). Por ello, no se crean tablas nuevas, sino que se convierten las tablas existentes, con sus datos.

\section{Dise\~no de las particiones}

\paragraph{Tablas particionadas.} Se particionan las dos tablas que crecen de forma continua y que se consultan por periodos:

\begin{itemize}[leftmargin=*]
	\item \nem{pago}: recibe un registro por cada mes de cada inversión, por lo que es la tabla que más crece con la operación. Los estados de cuenta y las tablas de interés se consultan por periodos de fechas.
	\item \nem{bitacora}: recibe un registro por cada cambio en los datos y nunca elimina registros. El auditor la revisa por periodos.
\end{itemize}

\paragraph{Criterio.} Ambas tablas se particionan por rango de fechas, con una partición por año. La clave de partición es \nem{pago\us fecha} en \nem{pago} y \nem{bitacora\us fecha} en \nem{bitacora}. Ambos gestores exigen que la clave de partición forme parte de la llave primaria, ya que cada partición verifica la unicidad de sus propias filas; por ese motivo, las llaves primarias de ambas tablas incluyen su fecha desde el modelo conceptual.

\paragraph{Particiones.} Las particiones de \nem{pago} corresponden a los años de los datos de prueba, y las de \nem{bitacora}, a los años en que el sistema ha registrado o registrará cambios, ya que la bitácora inicia su operación en 2026. La primera partición de cada tabla recibe también las fechas anteriores a su año, y la última, todas las posteriores, de modo que ninguna fecha queda sin partición (tabla~\ref{tab:particiones}). Las particiones tienen el mismo nombre en ambos gestores.

\begin{table}[H]
	\centering
	\small
	\caption{Particiones de \nem{pago} y \nem{bitacora}}
	\label{tab:particiones}
	\begin{tabular}{@{}L{3.4cm} L{6cm} L{2.6cm} r@{}}
		\toprule
		\textbf{Partición} & \textbf{Fechas} & \textbf{Tablespace (PostgreSQL)} & \textbf{Registros} \\
		\midrule
		\nem{pago\us 2024} & Anteriores al 1 de enero de 2025 & \nem{ts\us operacion} & 74 \\
		\nem{pago\us 2025} & Del 1 de enero al 31 de diciembre de 2025 & \nem{ts\us operacion} & 237 \\
		\nem{pago\us 2026} & Del 1 de enero al 31 de diciembre de 2026 & \nem{ts\us operacion} & 168 \\
		\nem{pago\us posterior} & A partir del 1 de enero de 2027 & \nem{ts\us operacion} & 0 \\
		\midrule
		\nem{bitacora\us 2026} & Anteriores al 1 de enero de 2027 & \nem{ts\us bitacora} & 17 \\
		\nem{bitacora\us 2027} & Del 1 de enero al 31 de diciembre de 2027 & \nem{ts\us bitacora} & 0 \\
		\nem{bitacora\us posterior} & A partir del 1 de enero de 2028 & \nem{ts\us bitacora} & 0 \\
		\bottomrule
	\end{tabular}
\end{table}

Los 168 pagos de 2026 incluyen el que registró la tesorería durante las pruebas de alcance, y los 17 registros de la bitácora son los que dejaron esas pruebas en 2026. En MySQL, las particiones permanecen en el directorio de datos, ya que una tabla particionada no puede ubicarse en un tablespace general (capítulo~\ref{cap:tablespaces}).

\paragraph{Mantenimiento.} La partición posterior evita que una inserción falle por no tener partición, pero no debe acumular registros de varios años. Antes de que inicie un nuevo año, el DBA divide la partición posterior para crear la del año siguiente: en \nem{pago}, antes del 1 de enero de 2027; en \nem{bitacora}, antes del 1 de enero de 2028. Mientras la partición posterior está vacía, la división es inmediata. Del mismo modo, cuando la bitácora de un año ya no se necesita en línea, su partición puede desprenderse de la tabla y respaldarse por separado. Esta es la única forma de retirar registros de la bitácora, ya que su trigger de protección impide eliminarlos con \nem{DELETE}; al tratarse de un cambio de estructura, queda en el registro del servidor.

\paragraph{Conversión de las tablas.} Los gestores difieren en la forma de convertir una tabla existente:

\begin{itemize}[leftmargin=*]
	\item \textbf{PostgreSQL} no permite convertir una tabla en particionada. Se crea una tabla particionada nueva, se copian los datos y se elimina la original, dentro de una sola transacción. La nueva tabla conserva la llave foránea de \nem{pago}, porque PostgreSQL admite llaves foráneas en tablas particionadas, pero sus triggers y privilegios deben volver a crearse.
	\item \textbf{MySQL} convierte la tabla sin reemplazarla, con \nem{ALTER TABLE \ldots{} PARTITION BY}, y conserva sus triggers, privilegios y numeración. Sin embargo, InnoDB no admite llaves foráneas en tablas particionadas, por lo que la llave foránea de \nem{pago} se elimina y su regla se valida con triggers.
\end{itemize}

\section{Diagrama relacional con particiones}

La figura~\ref{fig:diagrama-particiones} muestra las tablas particionadas de la segunda instancia, sus particiones y el tablespace de cada una. Las demás tablas y relaciones no cambian respecto al diagrama relacional, con una excepción en MySQL: la relación entre \nem{interes} y \nem{pago} ya no se garantiza con una llave foránea, sino con triggers.

\captura{diagrama_particiones.jpg}{Diagrama relacional con particiones}{diagrama-particiones}

\section{PostgreSQL}

\paragraph{Conversión.} El DBA convierte ambas tablas en una sola transacción. Si cualquier paso falla, la transacción se revierte completa y las tablas originales quedan intactas; además, ningún usuario puede registrar cambios a mitad de la conversión, cuando las tablas originales ya cambiaron de nombre. El script sigue estos pasos:

\begin{enumerate}[leftmargin=*]
	\item Cambia el nombre de las tablas originales, de sus llaves primarias y de la secuencia de la bitácora, para liberar esos nombres. Las llaves primarias deben renombrarse porque crean un índice, y dos índices del mismo esquema no pueden tener el mismo nombre.
	\item Crea las tablas particionadas, con la misma definición que en la implementación física más la cláusula \nem{PARTITION BY RANGE}, y sus particiones en los tablespaces del capítulo~\ref{cap:tablespaces}. Cada partición recibe un índice para la llave primaria, que el gestor nombra automáticamente (por ejemplo, \nem{pago\us 2024\us pkey}) y ubica en el tablespace indicado por \nem{USING INDEX TABLESPACE}.
	\item Copia los datos, conservando sus llaves, y ajusta la secuencia de la bitácora al número más alto copiado.
	\item Elimina las tablas originales.
	\item Vuelve a crear los triggers de las nuevas tablas. La función de bitácora se reemplaza porque, en una tabla particionada, los triggers se ejecutan en cada partición, y \nem{TG\us TABLE\us NAME} devolvería, por ejemplo, \nem{pago\us 2026} en lugar de \nem{pago}. La función \nem{pg\us partition\us root} devuelve la tabla particionada a la que pertenece una partición, o un valor nulo si la tabla no está particionada; con \nem{COALESCE}, la función registra el nombre correcto en ambos casos.
	\item Otorga los privilegios sobre \nem{pago}. El auditor y el usuario de respaldo reciben los de las nuevas tablas por los privilegios por defecto de la implementación física, pero estos no abarcan secuencias, y \nem{pg\us dump} necesita leer la nueva secuencia de la bitácora.
\end{enumerate}

Al final, \nem{ANALYZE} actualiza las estadísticas de las nuevas tablas, que el planificador utiliza para elegir sus planes de ejecución.

\begin{verbatim}
	-- Ejecutado por dlandero en la segunda instancia, base de datos inversiones
	BEGIN;
	SET ROLE rol_admin_bd;
	
	-- 1. Las tablas originales se conservan con otro nombre
	ALTER TABLE pago RENAME TO pago_original;
	ALTER TABLE pago_original RENAME CONSTRAINT pk_pago TO pk_pago_original;
	ALTER TABLE bitacora RENAME TO bitacora_original;
	ALTER TABLE bitacora_original
	  RENAME CONSTRAINT pk_bitacora TO pk_bitacora_original;
	ALTER SEQUENCE bitacora_bitacora_num_seq
	  RENAME TO bitacora_original_bitacora_num_seq;
	
	-- 2. Tablas particionadas, con la misma definición
	CREATE TABLE pago (
	  inversion_num          INT           NOT NULL,
	  interes_fecha_inicial  DATE          NOT NULL,
	  pago_fecha             DATE          NOT NULL,
	  pago_interes           NUMERIC(14,2) NOT NULL,
	  pago_mora              NUMERIC(14,2) NOT NULL DEFAULT 0,
	  CONSTRAINT pk_pago
	    PRIMARY KEY (inversion_num, interes_fecha_inicial, pago_fecha)
	    USING INDEX TABLESPACE ts_operacion,
	  CONSTRAINT fk_pago_interes FOREIGN KEY (inversion_num, interes_fecha_inicial)
	    REFERENCES interes (inversion_num, interes_fecha_inicial)
	    ON DELETE RESTRICT ON UPDATE RESTRICT,
	  CONSTRAINT ck_pago_montos CHECK (pago_interes >= 0 AND pago_mora >= 0)
	) PARTITION BY RANGE (pago_fecha);
	
	CREATE TABLE pago_2024 PARTITION OF pago
	  FOR VALUES FROM (MINVALUE) TO ('2025-01-01') TABLESPACE ts_operacion;
	CREATE TABLE pago_2025 PARTITION OF pago
	  FOR VALUES FROM ('2025-01-01') TO ('2026-01-01') TABLESPACE ts_operacion;
	CREATE TABLE pago_2026 PARTITION OF pago
	  FOR VALUES FROM ('2026-01-01') TO ('2027-01-01') TABLESPACE ts_operacion;
	CREATE TABLE pago_posterior PARTITION OF pago
	  FOR VALUES FROM ('2027-01-01') TO (MAXVALUE) TABLESPACE ts_operacion;
	
	CREATE TABLE bitacora (
	  bitacora_num            BIGINT GENERATED BY DEFAULT AS IDENTITY,
	  bitacora_fecha          TIMESTAMP    NOT NULL DEFAULT LOCALTIMESTAMP,
	  bitacora_usuario        VARCHAR(100) NOT NULL,
	  bitacora_tabla          VARCHAR(63)  NOT NULL,
	  bitacora_operacion      VARCHAR(6)   NOT NULL,
	  bitacora_ip             VARCHAR(45)  NULL,
	  bitacora_sesion         BIGINT       NULL,
	  bitacora_query          TEXT         NULL,
	  bitacora_dato_anterior  JSONB        NULL,
	  bitacora_dato_nuevo     JSONB        NULL,
	  CONSTRAINT pk_bitacora PRIMARY KEY (bitacora_num, bitacora_fecha)
	    USING INDEX TABLESPACE ts_bitacora,
	  CONSTRAINT ck_bitacora_operacion
	    CHECK (bitacora_operacion IN ('INSERT', 'UPDATE', 'DELETE'))
	) PARTITION BY RANGE (bitacora_fecha);
	
	CREATE TABLE bitacora_2026 PARTITION OF bitacora
	  FOR VALUES FROM (MINVALUE) TO ('2027-01-01') TABLESPACE ts_bitacora;
	CREATE TABLE bitacora_2027 PARTITION OF bitacora
	  FOR VALUES FROM ('2027-01-01') TO ('2028-01-01') TABLESPACE ts_bitacora;
	CREATE TABLE bitacora_posterior PARTITION OF bitacora
	  FOR VALUES FROM ('2028-01-01') TO (MAXVALUE) TABLESPACE ts_bitacora;
	
	-- 3. Copia de los datos y de la numeración de la bitácora
	INSERT INTO pago SELECT * FROM pago_original;
	INSERT INTO bitacora SELECT * FROM bitacora_original;
	SELECT setval(pg_get_serial_sequence('bitacora', 'bitacora_num'),
	              (SELECT MAX(bitacora_num) FROM bitacora));
	
	-- 4. Eliminación de las tablas originales
	DROP TABLE pago_original;
	DROP TABLE bitacora_original;
	
	-- 5. Bitácora: nombre de la tabla raíz y triggers de las nuevas tablas
	CREATE OR REPLACE FUNCTION fn_bitacora()
	RETURNS trigger
	LANGUAGE plpgsql
	SECURITY DEFINER
	SET search_path = inversiones
	AS $$
	BEGIN
	  INSERT INTO bitacora (bitacora_fecha, bitacora_usuario, bitacora_tabla,
	                        bitacora_operacion, bitacora_ip, bitacora_sesion,
	                        bitacora_query, bitacora_dato_anterior,
	                        bitacora_dato_nuevo)
	  VALUES (LOCALTIMESTAMP, session_user,
	          COALESCE((SELECT relname FROM pg_class
	                    WHERE oid = pg_partition_root(TG_RELID)),
	                   TG_TABLE_NAME),
	          TG_OP,
	          host(inet_client_addr()), pg_backend_pid(), current_query(),
	          CASE WHEN TG_OP IN ('UPDATE', 'DELETE') THEN to_jsonb(OLD) END,
	          CASE WHEN TG_OP IN ('INSERT', 'UPDATE') THEN to_jsonb(NEW) END);
	  RETURN NULL;
	END;
	$$;
	
	CREATE TRIGGER trg_pago_bitacora
	  AFTER INSERT OR UPDATE OR DELETE ON pago
	  FOR EACH ROW EXECUTE FUNCTION fn_bitacora();
	
	CREATE TRIGGER trg_bitacora_proteger
	  BEFORE UPDATE OR DELETE ON bitacora
	  FOR EACH ROW EXECUTE FUNCTION fn_bitacora_proteger();
	
	CREATE TRIGGER trg_bitacora_truncate
	  BEFORE TRUNCATE ON bitacora
	  FOR EACH STATEMENT EXECUTE FUNCTION fn_bitacora_proteger();
	
	-- 6. Privilegios de las nuevas tablas. El auditor y el usuario de respaldo
	--    reciben los de las tablas por los privilegios por defecto, pero no
	--    el de la nueva secuencia de la bitácora, que necesita pg_dump.
	GRANT SELECT ON pago TO rol_gerente, rol_ejecutivo;
	GRANT SELECT, INSERT ON pago TO rol_tesoreria;
	GRANT SELECT ON SEQUENCE bitacora_bitacora_num_seq TO rol_respaldo;
	
	COMMIT;
	
	ANALYZE pago;
	ANALYZE bitacora;
\end{verbatim}

\paragraph{Verificación.} La primera consulta muestra las particiones de cada tabla, su rango y su tablespace; la segunda cuenta los registros de cada partición (figura~\ref{fig:pg-particiones}). Los registros suman 479 pagos y 17 registros de bitácora, los mismos que antes de la conversión.

\begin{verbatim}
	SELECT p.relname AS tabla, c.relname AS particion,
	       pg_get_expr(c.relpartbound, c.oid) AS rango,
	       t.spcname AS tablespace
	FROM pg_inherits i
	  JOIN pg_class p ON p.oid = i.inhparent
	  JOIN pg_class c ON c.oid = i.inhrelid
	  LEFT JOIN pg_tablespace t ON t.oid = c.reltablespace
	WHERE p.relkind = 'p'
	  AND p.relnamespace = 'inversiones'::regnamespace
	ORDER BY p.relname, c.relname;
	
	SELECT 'pago_2024' AS particion, COUNT(*) AS registros FROM pago_2024
	UNION ALL SELECT 'pago_2025', COUNT(*) FROM pago_2025
	UNION ALL SELECT 'pago_2026', COUNT(*) FROM pago_2026
	UNION ALL SELECT 'pago_posterior', COUNT(*) FROM pago_posterior
	UNION ALL SELECT 'bitacora_2026', COUNT(*) FROM bitacora_2026
	UNION ALL SELECT 'bitacora_2027', COUNT(*) FROM bitacora_2027
	UNION ALL SELECT 'bitacora_posterior', COUNT(*) FROM bitacora_posterior;
\end{verbatim}

\captura{pg-particiones.png}{Particiones y registros por partición en PostgreSQL}{pg-particiones}

La poda de particiones se comprueba con el plan de ejecución de una consulta que filtra los pagos de 2025 (figura~\ref{fig:pg-poda}). El plan recorre únicamente la partición \nem{pago\us 2025}.

\begin{verbatim}
	EXPLAIN
	SELECT COUNT(*), SUM(pago_interes)
	FROM pago
	WHERE pago_fecha BETWEEN DATE '2025-01-01' AND DATE '2025-12-31';
\end{verbatim}

\captura{pg-poda.png}{Poda de particiones en PostgreSQL}{pg-poda}

Por último, se comprueba que las reglas del esquema siguen vigentes (figura~\ref{fig:pg-integridad}). La tesorería intenta registrar el pago de un periodo de interés inexistente, que la llave foránea rechaza. Después, el DBA registra un pago válido dentro de una transacción que se revierte al final, y consulta el último registro de la bitácora, que debe indicar la tabla \nem{pago} y no el nombre de la partición. Ambos scripts son iguales en los dos gestores.

\begin{verbatim}
	-- Como xzaleta: pago de un periodo de interés inexistente
	INSERT INTO pago (inversion_num, interes_fecha_inicial, pago_fecha,
	                  pago_interes)
	VALUES (16, '2030-01-01', '2026-10-30', 100.00);
\end{verbatim}

\begin{verbatim}
	-- Como dlandero: pago válido dentro de una transacción que se revierte
	START TRANSACTION;
	INSERT INTO pago (inversion_num, interes_fecha_inicial, pago_fecha,
	                  pago_interes)
	VALUES (16, '2025-06-05', '2026-10-30', 4250.00);
	SELECT bitacora_num, bitacora_tabla, bitacora_operacion
	FROM bitacora
	ORDER BY bitacora_num DESC
	LIMIT 1;
	ROLLBACK;
\end{verbatim}

\captura{pg-integridad.png}{Integridad y bitácora con las tablas particionadas en PostgreSQL}{pg-integridad}

\section{MySQL}

\paragraph{Conversión.} El DBA convierte las tablas con \nem{ALTER TABLE}. A diferencia de PostgreSQL, las sentencias de definición de datos de MySQL no pueden revertirse dentro de una transacción, pero cada \nem{ALTER TABLE} es atómico: si falla, la tabla queda como estaba. El script sigue estos pasos:

\begin{enumerate}[leftmargin=*]
	\item Elimina la llave foránea \nem{fk\us pago\us interes}, ya que InnoDB no admite llaves foráneas en tablas particionadas.
	\item Particiona \nem{pago} y \nem{bitacora} por rango de su columna de fecha. La partición \nem{RANGE COLUMNS} compara directamente los valores de la columna, sin convertirlos en números. El gestor reconstruye cada tabla repartiendo sus filas entre las particiones; la reconstrucción no activa los triggers, por lo que no genera registros en la bitácora.
	\item Sustituye la llave foránea por dos procedimientos y cuatro triggers, con el mismo esquema que las validaciones de la implementación física. Los triggers de \nem{pago} rechazan un pago cuyo periodo de interés no existe, y los de \nem{interes} rechazan eliminar un periodo con pagos o cambiar su llave, que es el comportamiento de \nem{ON DELETE RESTRICT} y \nem{ON UPDATE RESTRICT}.
\end{enumerate}

A diferencia de una llave foránea, los triggers no bloquean el registro referenciado mientras se completa la operación, por lo que no protegen contra dos operaciones simultáneas que se contradigan. En este sistema el riesgo no existe en la práctica, porque ningún rol de operación puede eliminar ni modificar periodos de interés.

\begin{verbatim}
	-- Ejecutado por dlandero en la segunda instancia, base de datos inversiones
	
	-- 1. Las tablas particionadas no admiten llaves foráneas
	ALTER TABLE pago DROP FOREIGN KEY fk_pago_interes;
	
	-- 2. Particionamiento de las tablas existentes
	ALTER TABLE pago
	  PARTITION BY RANGE COLUMNS (pago_fecha) (
	    PARTITION pago_2024      VALUES LESS THAN ('2025-01-01'),
	    PARTITION pago_2025      VALUES LESS THAN ('2026-01-01'),
	    PARTITION pago_2026      VALUES LESS THAN ('2027-01-01'),
	    PARTITION pago_posterior VALUES LESS THAN (MAXVALUE)
	  );
	
	ALTER TABLE bitacora
	  PARTITION BY RANGE COLUMNS (bitacora_fecha) (
	    PARTITION bitacora_2026      VALUES LESS THAN ('2027-01-01'),
	    PARTITION bitacora_2027      VALUES LESS THAN ('2028-01-01'),
	    PARTITION bitacora_posterior VALUES LESS THAN (MAXVALUE)
	  );
	
	-- 3. La regla de la llave foránea, mediante triggers
	DELIMITER //
	
	CREATE DEFINER = CURRENT_USER PROCEDURE validar_pago_interes(
	  IN p_inversion_num INT, IN p_interes_fecha_inicial DATE)
	READS SQL DATA
	BEGIN
	  IF NOT EXISTS (SELECT 1 FROM interes
	                 WHERE inversion_num = p_inversion_num
	                   AND interes_fecha_inicial = p_interes_fecha_inicial) THEN
	    SIGNAL SQLSTATE '45000'
	      SET MESSAGE_TEXT = 'El periodo de interés del pago no existe';
	  END IF;
	END //
	
	CREATE DEFINER = CURRENT_USER PROCEDURE validar_interes_sin_pagos(
	  IN p_inversion_num INT, IN p_interes_fecha_inicial DATE)
	READS SQL DATA
	BEGIN
	  IF EXISTS (SELECT 1 FROM pago
	             WHERE inversion_num = p_inversion_num
	               AND interes_fecha_inicial = p_interes_fecha_inicial) THEN
	    SIGNAL SQLSTATE '45000'
	      SET MESSAGE_TEXT = 'El periodo de interés tiene pagos registrados';
	  END IF;
	END //
	
	CREATE DEFINER = CURRENT_USER TRIGGER trg_pago_interes_ins
	  BEFORE INSERT ON pago FOR EACH ROW
	BEGIN
	  CALL validar_pago_interes(NEW.inversion_num, NEW.interes_fecha_inicial);
	END //
	
	CREATE DEFINER = CURRENT_USER TRIGGER trg_pago_interes_upd
	  BEFORE UPDATE ON pago FOR EACH ROW
	BEGIN
	  CALL validar_pago_interes(NEW.inversion_num, NEW.interes_fecha_inicial);
	END //
	
	CREATE DEFINER = CURRENT_USER TRIGGER trg_interes_pago_del
	  BEFORE DELETE ON interes FOR EACH ROW
	BEGIN
	  CALL validar_interes_sin_pagos(OLD.inversion_num, OLD.interes_fecha_inicial);
	END //
	
	CREATE DEFINER = CURRENT_USER TRIGGER trg_interes_pago_upd
	  BEFORE UPDATE ON interes FOR EACH ROW
	BEGIN
	  IF NEW.inversion_num <> OLD.inversion_num
	     OR NEW.interes_fecha_inicial <> OLD.interes_fecha_inicial THEN
	    CALL validar_interes_sin_pagos(OLD.inversion_num,
	                                   OLD.interes_fecha_inicial);
	  END IF;
	END //
	
	DELIMITER ;
\end{verbatim}

\paragraph{Verificación.} La primera consulta muestra las particiones de cada tabla desde el diccionario de datos; la segunda cuenta los registros de cada partición, que se seleccionan con la cláusula \nem{PARTITION} (figura~\ref{fig:mysql-particiones}). Los resultados coinciden con los de PostgreSQL.

\begin{verbatim}
	SELECT TABLE_NAME AS tabla, PARTITION_NAME AS particion,
	       PARTITION_METHOD AS metodo, PARTITION_EXPRESSION AS columna,
	       PARTITION_DESCRIPTION AS menor_que
	FROM information_schema.PARTITIONS
	WHERE TABLE_SCHEMA = 'inversiones'
	  AND PARTITION_NAME IS NOT NULL
	ORDER BY TABLE_NAME, PARTITION_ORDINAL_POSITION;
	
	SELECT 'pago_2024' AS particion, COUNT(*) AS registros
	FROM pago PARTITION (pago_2024)
	UNION ALL SELECT 'pago_2025', COUNT(*) FROM pago PARTITION (pago_2025)
	UNION ALL SELECT 'pago_2026', COUNT(*) FROM pago PARTITION (pago_2026)
	UNION ALL SELECT 'pago_posterior', COUNT(*) FROM pago PARTITION (pago_posterior)
	UNION ALL SELECT 'bitacora_2026', COUNT(*) FROM bitacora PARTITION (bitacora_2026)
	UNION ALL SELECT 'bitacora_2027', COUNT(*) FROM bitacora PARTITION (bitacora_2027)
	UNION ALL SELECT 'bitacora_posterior', COUNT(*)
	          FROM bitacora PARTITION (bitacora_posterior);
\end{verbatim}

\captura{mysql-particiones.png}{Particiones y registros por partición en MySQL}{mysql-particiones}

La poda se comprueba con la misma consulta que en PostgreSQL. En MySQL, la columna \nem{partitions} del plan muestra únicamente \nem{pago\us 2025} (figura~\ref{fig:mysql-poda}).

\captura{mysql-poda.png}{Poda de particiones en MySQL}{mysql-poda}

Los scripts de integridad también son los mismos (figura~\ref{fig:mysql-integridad}). En MySQL, el pago del periodo inexistente lo rechaza el trigger \nem{trg\us pago\us interes\us ins}, con el mensaje \nem{El periodo de interés del pago no existe}.

\captura{mysql-integridad.png}{Integridad y bitácora con las tablas particionadas en MySQL}{mysql-integridad}

	\chapter{Segmentaci\'on}
\label{cap:segmentacion}

La segmentación divide una tabla en fragmentos lógicos mediante vistas. A diferencia del particionamiento, la tabla no se modifica físicamente: sus datos permanecen en un solo lugar y cada fragmento es una vista que selecciona una parte de sus filas o de sus columnas. Los usuarios consultan el fragmento que necesitan como si fuera una tabla, y el gestor obtiene sus datos de la tabla original en cada consulta.

Como se indicó en el capítulo~\ref{cap:multiinstancias}, la segmentación se implementa en la segunda instancia de cada gestor, junto con el particionamiento.

\section{Dise\~no de la segmentaci\'on}

\paragraph{Tipos de fragmentación.} El esquema aplica los dos tipos de fragmentación, cada uno sobre la tabla en la que tiene sentido (tabla~\ref{tab:fragmentacion}).

\begin{table}[H]
	\centering
	\small
	\caption{Tipos de fragmentación aplicados}
	\label{tab:fragmentacion}
	\begin{tabular}{@{}L{2.4cm} L{2.2cm} L{5.2cm} L{6cm}@{}}
		\toprule
		\textbf{Tipo} & \textbf{Tabla} & \textbf{Criterio} & \textbf{Cada fragmento contiene} \\
		\midrule
		Horizontal & \nem{inversion} & Estado de la inversión & Todas las columnas de las inversiones que se encuentran en un estado \\
		Vertical & \nem{bitacora} & Uso de las columnas & Todos los registros, con un subconjunto de las columnas y la llave primaria \\
		\bottomrule
	\end{tabular}
\end{table}

\paragraph{Reglas de corrección.} Una fragmentación es correcta si cumple tres reglas, que el diseño de los fragmentos de este esquema garantiza:

\begin{itemize}[leftmargin=*]
	\item \textbf{Completitud:} todo dato de la tabla aparece en algún fragmento. En la fragmentación horizontal, toda fila pertenece a un fragmento; en la vertical, toda columna.
	\item \textbf{Disyunción:} ningún dato aparece en más de un fragmento. En la fragmentación vertical, la excepción es la llave primaria, que se repite en todos los fragmentos para poder reunirlos.
	\item \textbf{Reconstrucción:} la tabla original puede obtenerse a partir de sus fragmentos: con la unión de los fragmentos horizontales (\nem{UNION ALL}) o con la reunión de los fragmentos verticales por su llave primaria (\nem{JOIN}).
\end{itemize}

\subsection{Fragmentación horizontal de \nem{inversion}}

La tabla \nem{inversion} se divide en tres fragmentos, uno por cada estado de una inversión definido en las reglas de negocio (tabla~\ref{tab:fragmentos-inversion}).

\begin{table}[H]
	\centering
	\small
	\caption{Fragmentos horizontales de \nem{inversion}}
	\label{tab:fragmentos-inversion}
	\begin{tabular}{@{}L{3.6cm} L{7.4cm} L{4.8cm}@{}}
		\toprule
		\textbf{Vista} & \textbf{Condición} & \textbf{Uso principal} \\
		\midrule
		\nem{inversion\us vigente} & Sin liquidación y con vencimiento igual o posterior a la fecha actual & Seguimiento del ejecutivo y del gerente \\
		\nem{inversion\us vencida} & Sin liquidación y con vencimiento anterior a la fecha actual & Cobro de adeudos de la tesorería \\
		\nem{inversion\us liquidada} & Con liquidación registrada & Revisión del auditor \\
		\bottomrule
	\end{tabular}
\end{table}

Las condiciones son excluyentes y cubren todos los casos: una inversión tiene o no tiene liquidación, y si no la tiene, su vencimiento es anterior a la fecha actual o no lo es. Por ello, la fragmentación es completa y disjunta por construcción.

\paragraph{Por qué vistas y no particiones.} El estado de una inversión no se almacena en ninguna columna: depende de la fecha actual y de la existencia de un registro en otra tabla, \nem{liquidacion}. Una inversión vigente pasa a estar vencida el día siguiente a su vencimiento sin que nadie modifique un solo registro. Por ello, el estado no puede servir como clave de partición: el gestor ubica cada fila en una partición según el valor almacenado en una columna de la propia tabla, y lo hace al insertarla o modificarla. Agregar una columna de estado obligaría a actualizarla diariamente con un proceso externo y repetiría información que ya se deduce de los datos. Una vista, en cambio, evalúa la condición en cada consulta, por lo que su contenido siempre corresponde al estado actual.

\paragraph{Contenido variable.} Como consecuencia, el contenido de las vistas cambia con el paso del tiempo. Con los datos de prueba, las 40 inversiones se reparten en 15 vigentes, 10 vencidas y 15 liquidadas. Mientras no se registren nuevas liquidaciones, estos conteos se mantienen hasta el 19 de enero de 2027, fecha de vencimiento de la inversión 30; al día siguiente, esa inversión pasa al fragmento de vencidas. Las reglas de corrección, en cambio, se cumplen en cualquier fecha.

\subsection{Fragmentación vertical de \nem{bitacora}}

Las columnas de \nem{bitacora} se usan de dos formas distintas. Las que identifican el evento (quién, cuándo, desde dónde, en qué tabla y con qué operación) se consultan constantemente para revisar la actividad del sistema. Las que contienen la sentencia y los valores anteriores y nuevos solo se consultan al revisar un cambio en particular, y son además las más pesadas, de tipo \nem{TEXT} y \nem{JSON}. La tabla se divide en dos fragmentos según ese uso (tabla~\ref{tab:fragmentos-bitacora}).

\begin{table}[H]
	\centering
	\small
	\caption{Fragmentos verticales de \nem{bitacora}}
	\label{tab:fragmentos-bitacora}
	\begin{tabular}{@{}L{3.4cm} L{8.2cm} L{4.2cm}@{}}
		\toprule
		\textbf{Vista} & \textbf{Columnas} & \textbf{Uso principal} \\
		\midrule
		\nem{bitacora\us resumen} & \nem{bitacora\us num}, \nem{bitacora\us fecha}, \nem{bitacora\us usuario}, \nem{bitacora\us tabla}, \nem{bitacora\us operacion}, \nem{bitacora\us ip}, \nem{bitacora\us sesion} & Revisión de la actividad: quién hizo qué y cuándo \\
		\nem{bitacora\us detalle} & \nem{bitacora\us num}, \nem{bitacora\us fecha}, \nem{bitacora\us query}, \nem{bitacora\us dato\us anterior}, \nem{bitacora\us dato\us nuevo} & Revisión del contenido de un cambio \\
		\bottomrule
	\end{tabular}
\end{table}

Ambos fragmentos conservan la llave primaria, (\nem{bitacora\us num}, \nem{bitacora\us fecha}), que permite reunirlos para reconstruir cada registro. Cada una de las demás columnas aparece en un solo fragmento, por lo que la fragmentación vertical también es completa y disjunta por construcción. El auditor localiza primero los eventos en \nem{bitacora\us resumen} y después consulta el detalle de los que le interesan.

\paragraph{Por qué vistas y no tablas separadas.} Dividir la bitácora en dos tablas físicas obligaría a cada trigger a escribir en ambas, y en MySQL no podría haber una llave foránea entre ellas, porque las tablas particionadas no admiten llaves foráneas. Con vistas, los triggers siguen escribiendo un solo registro en una sola tabla, y la fragmentación no modifica la forma de auditar descrita en el capítulo~\ref{cap:bitacora}. Las vistas tampoco dependen de que la tabla esté particionada: el gestor resuelve las particiones al consultarlas.

\subsection{Uso y privilegios de las vistas}

Las vistas son de solo lectura para los usuarios del sistema: las altas y los cambios se siguen haciendo sobre las tablas, con los privilegios del capítulo~\ref{cap:implementacion}. Cada vista se ejecuta con los privilegios de su dueño, el DBA, por lo que otorgar el privilegio \nem{SELECT} sobre una vista da acceso a sus datos aunque el usuario no tenga acceso a la tabla. Por ello, los privilegios sobre las vistas siguen la misma matriz que los de las tablas:

\begin{itemize}[leftmargin=*]
	\item Los fragmentos de \nem{inversion} se otorgan al gerente, al ejecutivo y a la tesorería, que ya pueden consultar la tabla completa.
	\item Los fragmentos de \nem{bitacora} no se otorgan a ningún rol de operación, ya que solo el auditor puede consultar la bitácora.
	\item El auditor y el usuario de respaldo los reciben sin otorgamientos adicionales: en PostgreSQL, por los privilegios por defecto definidos en el capítulo~\ref{cap:implementacion}; en MySQL, porque sus privilegios abarcan toda la base de datos (\nem{inversiones.*}).
\end{itemize}

Las vistas se definen con \nem{SELECT i.*} en los fragmentos horizontales, ya que estos conservan todas las columnas. Ambos gestores expanden el asterisco al crear la vista, por lo que, si en el futuro se agrega una columna a \nem{inversion}, las vistas deben volver a crearse para incluirla.

\section{Diagrama relacional con segmentaci\'on}

Las figuras~\ref{fig:segmentacion-horizontal} y~\ref{fig:segmentacion-vertical} muestran las vistas que se derivan de cada tabla. Las tablas se dibujan con borde continuo y encabezado oscuro, y las vistas con borde punteado y encabezado claro. Las flechas continuas indican de qué tabla obtiene sus datos cada vista; las punteadas, la tabla que se consulta para evaluar su condición. Las demás tablas y relaciones del esquema no cambian respecto al diagrama relacional (figura~\ref{fig:diagrama-er}).

\captura{segmentacion_horizontal.jpg}{Fragmentación horizontal de \nem{inversion}}{segmentacion-horizontal}
Las flechas continuas indican el origen de los datos; las punteadas, la tabla que se consulta para evaluar la condición.

\captura[0.85\textwidth]{segmentacion_vertical.jpg}{Fragmentación vertical de \nem{bitacora}}{segmentacion-vertical}
Ambas vistas conservan la llave primaria, que permite reconstruir cada registro.

\section{PostgreSQL}

El DBA crea las vistas en la segunda instancia, con el rol \nem{rol\us admin\us bd} activo para que el rol sea su dueño, igual que el de las tablas. En la misma secuencia otorga los privilegios sobre los fragmentos de \nem{inversion}.

\begin{verbatim}
	-- Ejecutado por dlandero en la segunda instancia, base de datos inversiones
	SET ROLE rol_admin_bd;

	-- Fragmentación horizontal de inversion, por estado
	CREATE VIEW inversion_vigente AS
	SELECT i.*
	FROM inversion i
	WHERE i.inversion_fecha_vencimiento >= CURRENT_DATE
	  AND NOT EXISTS (SELECT 1 FROM liquidacion l
	                  WHERE l.inversion_num = i.inversion_num);

	CREATE VIEW inversion_vencida AS
	SELECT i.*
	FROM inversion i
	WHERE i.inversion_fecha_vencimiento < CURRENT_DATE
	  AND NOT EXISTS (SELECT 1 FROM liquidacion l
	                  WHERE l.inversion_num = i.inversion_num);

	CREATE VIEW inversion_liquidada AS
	SELECT i.*
	FROM inversion i
	WHERE EXISTS (SELECT 1 FROM liquidacion l
	              WHERE l.inversion_num = i.inversion_num);

	-- Fragmentación vertical de bitacora
	CREATE VIEW bitacora_resumen AS
	SELECT bitacora_num, bitacora_fecha, bitacora_usuario, bitacora_tabla,
	       bitacora_operacion, bitacora_ip, bitacora_sesion
	FROM bitacora;

	CREATE VIEW bitacora_detalle AS
	SELECT bitacora_num, bitacora_fecha, bitacora_query,
	       bitacora_dato_anterior, bitacora_dato_nuevo
	FROM bitacora;

	-- Privilegios
	GRANT SELECT ON inversion_vigente, inversion_vencida, inversion_liquidada
	  TO rol_gerente, rol_ejecutivo, rol_tesoreria;
\end{verbatim}

\section{MySQL}

Las vistas tienen la misma definición que en PostgreSQL. Se crean con \nem{DEFINER = CURRENT\us USER} y \nem{SQL SECURITY DEFINER}, de modo que se ejecutan con los privilegios de la cuenta del DBA, igual que los procedimientos y triggers del esquema. Como MySQL solo permite otorgar privilegios sobre un objeto por sentencia, cada fragmento de \nem{inversion} requiere su propio \nem{GRANT}.

\begin{verbatim}
	-- Ejecutado por dlandero en la segunda instancia, base de datos inversiones

	-- Fragmentación horizontal de inversion, por estado
	CREATE DEFINER = CURRENT_USER SQL SECURITY DEFINER VIEW inversion_vigente AS
	SELECT i.*
	FROM inversion i
	WHERE i.inversion_fecha_vencimiento >= CURRENT_DATE
	  AND NOT EXISTS (SELECT 1 FROM liquidacion l
	                  WHERE l.inversion_num = i.inversion_num);

	CREATE DEFINER = CURRENT_USER SQL SECURITY DEFINER VIEW inversion_vencida AS
	SELECT i.*
	FROM inversion i
	WHERE i.inversion_fecha_vencimiento < CURRENT_DATE
	  AND NOT EXISTS (SELECT 1 FROM liquidacion l
	                  WHERE l.inversion_num = i.inversion_num);

	CREATE DEFINER = CURRENT_USER SQL SECURITY DEFINER VIEW inversion_liquidada AS
	SELECT i.*
	FROM inversion i
	WHERE EXISTS (SELECT 1 FROM liquidacion l
	              WHERE l.inversion_num = i.inversion_num);

	-- Fragmentación vertical de bitacora
	CREATE DEFINER = CURRENT_USER SQL SECURITY DEFINER VIEW bitacora_resumen AS
	SELECT bitacora_num, bitacora_fecha, bitacora_usuario, bitacora_tabla,
	       bitacora_operacion, bitacora_ip, bitacora_sesion
	FROM bitacora;

	CREATE DEFINER = CURRENT_USER SQL SECURITY DEFINER VIEW bitacora_detalle AS
	SELECT bitacora_num, bitacora_fecha, bitacora_query,
	       bitacora_dato_anterior, bitacora_dato_nuevo
	FROM bitacora;

	-- Privilegios
	GRANT SELECT ON inversiones.inversion_vigente
	  TO rol_gerente, rol_ejecutivo, rol_tesoreria;
	GRANT SELECT ON inversiones.inversion_vencida
	  TO rol_gerente, rol_ejecutivo, rol_tesoreria;
	GRANT SELECT ON inversiones.inversion_liquidada
	  TO rol_gerente, rol_ejecutivo, rol_tesoreria;
\end{verbatim}


	\chapter{\'Indices}
\section{An\'alisis de columnas candidatas}
\section{PostgreSQL}
\section{MySQL}
\end{document}
